\documentclass[11pt, reqno]{amsart}

% Abridged version of main.tex. The preamble and the macros are those of the long version,
% copied into summary/preamble_summary.tex and summary/macros_summary.tex.
% preamble_summary.tex: packages, theorem environments and general macros of the abridged version.
% Copied from the preamble of the long version (main.tex); the draft-note machinery is omitted.
% GENUS workshop style: page geometry, fonts, hyperref, fancyhdr and the claim-status tags (genus-collab/meta)
\usepackage{genus}
% text width of the original note (170 mm), so that its displayed formulas still fit the line
\geometry{margin=2cm}
\usepackage[utf8]{inputenc}
\usepackage[T1]{fontenc}

%-------------------------------------------------------------------------------
% Packages (modelled on the preamble of [BS] and of the companion paper [LP2]; only what is needed)
\usepackage[dvipsnames]{xcolor}
\usepackage{amssymb}
\usepackage[mathscr]{eucal}
\usepackage{amsthm}
\usepackage{amsmath}
\usepackage{mathtools}
\definecolor{highlight}{HTML}{3465a4}
% hyperref is loaded by genus.sty, with the GENUS link colours
\hypersetup{hypertexnames=false, hyperfootnotes=false, linktocpage=true, bookmarksdepth=2}
\usepackage{cleveref}
\usepackage{url}
\usepackage{array}

%-------------------------------------------------------------------------------
% Typefaces. [BS] uses libertine with newtxmath; libertine needs fontaxes, which is not assumed
% here, so mathpazo (Palatino) is the substitute, as in the companion paper.
% (mathpazo and courier dropped: genus.sty sets Latin Modern)
\usepackage{microtype}
\linespread{1.06}
\setlength{\parskip}{4pt}
\setlength{\parindent}{0pt}
\setcounter{secnumdepth}{3}
% enumitem is not assumed: first-level enumerate is labelled (i), (ii), ... ; for named parts use
% description or \item[(F1)] with an explicit label.
\renewcommand{\labelenumi}{(\roman{enumi})}
\renewcommand{\theenumi}{(\roman{enumi})}
\numberwithin{equation}{section}

%-------------------------------------------------------------------------------
% Theorem environments (names as in [BS])
\theoremstyle{plain}
\newtheorem{thm}{Theorem}[section]
\newtheorem{lem}[thm]{Lemma}
\newtheorem{prop}[thm]{Proposition}
\newtheorem{cor}[thm]{Corollary}
\theoremstyle{definition}
\newtheorem{defn}[thm]{Definition}
\newtheorem{rmk}[thm]{Remark}
\newtheorem{notation}[thm]{Notation}
% genus.sty defines an example environment on its own counter; redefine it on the counter of thm
\let\example\undefined\let\endexample\undefined\expandafter\let\csname c@example\endcsname\undefined
\newtheorem{example}[thm]{Example}
\newtheorem{conv}[thm]{Convention}

\crefname{equation}{Eq.}{Eqs.}
\crefname{thm}{Theorem}{Theorems}
\crefname{lem}{Lemma}{Lemmas}
\crefname{prop}{Proposition}{Propositions}
\crefname{cor}{Corollary}{Corollaries}
\crefname{defn}{Definition}{Definitions}
\crefname{rmk}{Remark}{Remarks}
\crefname{notation}{Notation}{Notations}
\crefname{example}{Example}{Examples}
\crefname{conv}{Convention}{Conventions}
\crefname{section}{Section}{Sections}
\crefname{subsection}{Section}{Sections}
\crefname{appendix}{Appendix}{Appendices}

%-------------------------------------------------------------------------------
% Letters
\newcommand{\bbA}{\mathbb{A}}
\newcommand{\bbC}{\mathbb{C}}
\newcommand{\bbD}{\mathbb{D}}
\newcommand{\bbE}{\mathbb{E}}
\newcommand{\bbF}{\mathbb{F}}
\newcommand{\bbK}{\mathbb{K}}
\newcommand{\bbL}{\mathbb{L}}
\newcommand{\bbP}{\mathbb{P}}
\newcommand{\bbQ}{\mathbb{Q}}
\newcommand{\bbR}{\mathbb{R}}
\newcommand{\bbT}{\mathbb{T}}
\newcommand{\bbZ}{\mathbb{Z}}
\newcommand{\cA}{\mathcal{A}}
\newcommand{\cC}{\mathcal{C}}
\newcommand{\cE}{\mathcal{E}}
\newcommand{\cF}{\mathcal{F}}
\newcommand{\cG}{\mathcal{G}}
\newcommand{\cL}{\mathcal{L}}
\newcommand{\cN}{\mathcal{N}}
\newcommand{\cO}{\mathcal{O}}
\newcommand{\cP}{\mathcal{P}}
\newcommand{\cQ}{\mathcal{Q}}
\newcommand{\cW}{\mathcal{W}}
\newcommand{\cZ}{\mathcal{Z}}
\newcommand{\mgp}[1]{\mathsf{#1}}

%-------------------------------------------------------------------------------
% General
\newcommand{\dd}{\mathrm{d}}
\newcommand{\ri}{\mathrm{i}}
\newcommand{\re}{\mathrm{e}}
\newcommand{\de}{\partial}
% \llbracket: from newtxmath if available, otherwise a substitute (stmaryrd is not assumed)
\providecommand{\llbracket}{[\![}
\providecommand{\rrbracket}{]\!]}
\newcommand{\fps}[1]{\llbracket #1 \rrbracket}
\newcommand{\lser}[1]{(\!(#1)\!)}
\newcommand{\Frac}{\mathrm{Frac}}
\DeclareMathOperator{\Tr}{Tr}
\DeclareMathOperator{\End}{End}
\DeclareMathOperator{\Aut}{Aut}
\DeclareMathOperator{\Spec}{Spec}
\DeclareMathOperator*{\Res}{Res}
\DeclareMathOperator{\Disc}{Disc}
\DeclareMathOperator{\Li}{Li}
\DeclareMathOperator{\val}{val}
\DeclareMathOperator{\ev}{ev}
\newcommand{\tensor}{\otimes}
\newcommand{\coef}[1]{[#1]}                    % coefficient extraction, e.g. \coef{t^k}

%-------------------------------------------------------------------------------
% Geometry, tori, weights (weights as in [BS]: \alpha_i, \epsilon_i)
\newcommand{\Tmax}{\mgp{T}}                    % the five-torus
\newcommand{\TCY}{\mgp{T}_{\mathrm{CY}}}        % the Calabi--Yau four-torus E = \epsilon_1 + \epsilon_2
\newcommand{\TF}{\mgp{T}_{\mathrm{F}}}          % fibrewise torus \alpha_0 = 0, \alpha_i = \epsilon_i
\newcommand{\Tthree}{\mgp{T}_{3}}               % \alpha_0 = 0 inside \TCY
\newcommand{\wA}{\mathsf{A}}                    % \alpha_0 + \alpha_2 = E - \alpha_1
\newcommand{\wB}{\mathsf{B}}                    % \alpha_0 + \alpha_1 = E - \alpha_2
\newcommand{\uq}{\mathsf{u}}                    % (E - 2\alpha_1)(E - 2\alpha_2)
\newcommand{\Mbar}{\overline{M}}
\newcommand{\vir}{\mathrm{virt}}
\newcommand{\GW}{\mathsf{GW}}
\newcommand{\Hdg}{\bbE}                        % Hodge bundle
\newcommand{\Lamv}{\Lambda^{\!\vee}}           % \Lamv(u) = \sum_i (-1)^i \lambda_i u^{g-i}
\newcommand{\Hg}{\mathsf{Hg}}                  % \Hg_{\bw} = \prod_a \Lamv(w_a)/w_a^g (normalised Hodge class)
\newcommand{\Nb}{\cN_{\mathrm{cf}}}             % \cO(-1)_{\alpha_1} + \cO(-1)_{\alpha_2} + \cO_{-\epsilon_1} + \cO_{-\epsilon_2}
\newcommand{\bw}{\mathbf{w}}                   % (E,-\epsilon_1,-\epsilon_2)
\newcommand{\nw}{\mathsf{nw}}                  % (\alpha_1,\wA,-\epsilon_1,-\epsilon_2)
\newcommand{\wts}{\mathrm{wt}}                 % \wts_\pm = the five tangent weights at p_\pm
\newcommand{\eul}{\mathrm{e}}                  % \eul(\wts_\pm): product of the tangent weights
\newcommand{\sF}{\mathscr{F}}                  % \sF_g(\cQ) = \sum_{d\ge 0} \GW_{g,d}\cQ^d
\newcommand{\sG}{\mathscr{G}}                  % \sG_g = \sF_g - \sF_g(0)

%-------------------------------------------------------------------------------
% References to the long version of this paper: \Lref{Theorem~4.1} prints [L, Theorem 4.1].
\newcommand{\Lref}[1]{[L,~#1]}


% macros_summary.tex: notation macros of the abridged version, copied from the long version
% (main.tex and the local macros of its section files).
%-------------------------------------------------------------------------------
% Fields and rings
\newcommand{\Kal}{\bbK_{\alpha}}               % Q(\alpha_0,\alpha_1,\alpha_2,\epsilon_1,\epsilon_2)
\newcommand{\KCY}{\bbK_{\mathrm{CY}}}          % Q(\alpha_0,\alpha_1,\alpha_2,\epsilon_1), \epsilon_2 = E - \epsilon_1
\newcommand{\Kred}{\bbK}                       % Q(\alpha_1,\alpha_2,\epsilon_1) (a_0-adic arguments)
\newcommand{\Kthree}{\bbK_{3}}                 % Q(\alpha_1,\alpha_2,\epsilon_1) as the function field of \Tthree
\newcommand{\Cz}{\bbC\lser{z}}                 % formal Laurent series in z
\newcommand{\Rz}{\mathsf{R}_{z}}               % \Cz((\tau))
\newcommand{\Rzo}{\mathsf{R}_{z}^{\circ}}      % \bbC[[z]]((\tau))
\newcommand{\Oan}{\mathscr{O}}                 % convergent power series in t (or \tau)
\newcommand{\Man}{\mathscr{M}}                 % meromorphic germs
\newcommand{\QA}{\mathsf{QA}}                  % quadratic algebra K(\cQ)[U]/(\Pq)
\newcommand{\RRing}{\mathbf{R}}                % \Kal[[t]] (Hodge side)
\newcommand{\Ccov}{\mathcal{C}}                % double cover \mathsf{y}^2 = \Disc(\cQ)
\newcommand{\ydc}{\mathsf{y}}                  % \ydc^2 = \Disc(\cQ)

%-------------------------------------------------------------------------------
% Base and genus-zero data
\newcommand{\thQ}{\theta_{\cQ}}                % \cQ d/d\cQ = d/dt
\newcommand{\thU}{\theta_{U}}                  % U d/dU
\newcommand{\Pq}{\mathsf{P}}                   % \Pq(U) = \wB\cQ U^2 - (\alpha_0 + \cQ(\alpha_1-\alpha_2))U - \alpha_2
\newcommand{\Ns}{\mathsf{N}}                   % \Ns(s) = \wB\tau s^2 - (E+\alpha_0\tau)s + \wA
\newcommand{\sast}{s^{\ast}}                   % singular root of \Ns
\newcommand{\stil}{\tilde{s}}                  % regular root of \Ns
\newcommand{\Lv}{\mathsf{L}}                   % eigenvalue function \Lv(U); \Lv_i = \Lv(U_i)
\newcommand{\qprod}{\circ}                     % quantum product
\newcommand{\eid}{\mathfrak{e}}                % idempotents
\newcommand{\fid}{\mathfrak{f}}                % normalised idempotents \Delta_i^{1/2}\eid_i
\newcommand{\Lsum}{\mathsf{S}}                 % \Lv_1 + \Lv_2 = E/\tau - \alpha_0
\newcommand{\Lprod}{\mathsf{\Pi}}              % \Lv_1\Lv_2 = \alpha_1\alpha_2/\tau
\newcommand{\lrem}{\ell}                       % \lrem(t) = \Lv_1 - E/t
\newcommand{\muq}{\mu}                         % E/\tau - \wA
\newcommand{\omq}{\omega}                      % (E+z)/\tau - (\wA+z)
\newcommand{\Aop}{\bbA}                        % matrix of the quantum differential equation on covectors
\newcommand{\chiv}{\boldsymbol{\chi}}          % \chiv_i = (1,\Lv_i)
\newcommand{\pim}{\varpi}                      % \pim(a) = \langle a,\chiv_2\rangle
\newcommand{\Thf}{\Theta}                      % \Delta_2/(\Lv_2-\Lv_1)
\newcommand{\Pg}{\mathsf{P}_{g}}               % [t^{2g}] \epsilon_1\epsilon_2 t^2/(4\sinh(\epsilon_1 t/2)\sinh(\epsilon_2 t/2))
\newcommand{\Ph}[1]{\mathsf{P}_{#1}}           % \Ph{h} for other genera
\newcommand{\gapg}{\mathrm{gap}_g}
\newcommand{\Ssh}{\mathsf{S}_{\circ}}          % \Ssh(y) = y/(2\sinh(y/2))
\newcommand{\Aeul}{\cA}                        % \Aeul_n(\cQ): \Li_{-n} = \Aeul_n/(1-\cQ)^{n+1}

%-------------------------------------------------------------------------------
% Euler integral, Laplace calculus, Gauss
\newcommand{\dphi}{\breve{\phi}}               % \phi_\tau(\sast) - \phi_0(\wA/E)
\newcommand{\alg}{\mathfrak{a}}                % the unit \alg(\tau) with \alg(0)=1
\newcommand{\Acal}{\mathscr{A}}                % ring of formal sums \sum K_{pq} v^p z^q
\newcommand{\gauss}[2]{\langle\!\langle #1\rangle\!\rangle_{#2}}   % Gaussian functional
\newcommand{\Lnorm}{\mathscr{N}}               % \Lnorm[f,a;u_0]: Gaussian-normalised formal Laplace expansion
\newcommand{\Lhat}{\widehat{\mathscr{N}}}      % normalised to leading coefficient 1
\newcommand{\Lfull}{\mathscr{N}^{\mathrm{full}}}
\newcommand{\Rr}{\mathsf{Rr}}                  % remainder of the phase, \Rr_f(v)
\newcommand{\Nlap}{\mathcal{N}}                % \Nlap(z,\tau) = \Lhat[\phi_\tau,(1-\tau s)^{-1};\sast]
\newcommand{\wick}{\mathsf{wick}}              % \wick_m: Wick polynomials
\newcommand{\Abeta}{\mathsf{A}_{\beta}}        % \Lnorm[\phi_0,1;\wA/E]
\newcommand{\ga}{\mathsf{a}}                   % 1 + \alpha_2/z
\newcommand{\gb}{\mathsf{b}}                   % 1 + \wA/z
\newcommand{\gc}{\mathsf{c}}                   % 2 + E/z
\newcommand{\Gauss}{\mathsf{Gauss}}            % Gauss's operator
\newcommand{\Fhyp}{\mathsf{F}}                 % {}_2F_1(\ga,\gb;\gc;\tau)
\newcommand{\Ephi}{\mathsf{X}}                 % \exp(-\dphi/z); coefficients \Ephi_n
\newcommand{\Hcoef}{\mathsf{H}}                % \Hcoef_n = [\tau^n](\Fhyp\Ephi)
\newcommand{\Kcoef}{\mathsf{K}}                % \Kcoef_n = [\tau^n](\Fhyp\Ephi\alg)
\newcommand{\Dz}{\mathsf{D}_{z}}               % z d/dt
\newcommand{\Dd}{\mathsf{D}}                   % abstract derivation symbol in the operator identity
\newcommand{\Dop}{\mathsf{D}}                  % \Dop_i = d/dt - \Delta_i'/(2\Delta_i)
\newcommand{\Eik}{\mathsf{Eik}}
\newcommand{\Ptr}{\mathsf{P}_{\mathrm{tr}}}
\newcommand{\Rtr}{\mathsf{R}_{\mathrm{tr}}}
\newcommand{\Ttr}{\mathsf{T}_{\mathrm{tr}}}
\newcommand{\Sph}{\mathsf{S}_{\phi}}           % eikonal phase
\newcommand{\Phia}{\Phi}                       % transport amplitude
\newcommand{\Ddisc}{\mathbb{D}}                % discs \Ddisc, \Ddisc', \Ddisc''
\newcommand{\Pol}{\mathrm{Pol}}                % \Pol_n(p)
\newcommand{\wrk}{\mathsf{w}}                  % \wrk_k = \tau^k \cW_k(U_1)  (\wr is the wreath-product symbol)
\newcommand{\zq}{\zeta}                        % \tau e^{\dphi/E}
\newcommand{\amp}{\mathsf{m}}                  % amplitude of a Laplace expansion (s03, s04)
\newcommand{\Dam}{\mathsf{D}_{\mathrm{am}}}     % \Dam q = d_tau q + q (d_tau f)/z (s03, s04)
\newcommand{\sQ}{\mathsf{s}}                    % 1/\cQ, the local coordinate at \cQ = \infty (s02, s07, appendix; not the Euler variable)

%-------------------------------------------------------------------------------
% Rows, factors, graph sums (one frame: frame B; point-frame objects carry a superscript pt)
\newcommand{\rpt}{\mathrm{r}^{\mathrm{pt}}}    % \rpt_{0,i}: point-frame identity row
\newcommand{\rZ}{\mathrm{r}^{Z}}               % \rZ_{0,i} = h_{\nw}\rpt_{0,i}: the row of R_Z
\newcommand{\rr}{\mathrm{r}}                   % \rr_{a,i}: frame-B rows
\newcommand{\vv}{\mathbf{r}}                   % \vv_i = (\rr_{0,i},\rr_{1,i})
\newcommand{\Rm}{\mathsf{R}}                   % R-matrix at generic weights
\newcommand{\RZ}{\mathsf{R}_{Z}}               % R-matrix with rows \rZ
\newcommand{\RB}{\mathsf{R}'}                  % frame-B R-matrix \RZ h_{\bw}^{-1}
\newcommand{\hw}{h_{\bw}}                      % h_{(E,-\epsilon_1,-\epsilon_2)}
\newcommand{\hnw}{h_{\nw}}                     % h_{(\alpha_1,\wA,-\epsilon_1,-\epsilon_2)}
\newcommand{\hreg}{h_{\mathrm{reg}}}           % h_{(\alpha_1,\alpha_2,\wA,\wB,-E,-E)}
\newcommand{\bk}{\mathsf{b}}                   % \bk_k = B_{2k}/(2k(2k-1))
\newcommand{\TFT}{\boldsymbol{\omega}}         % semisimple TFT
\newcommand{\Rinv}{\mathbf{B}}                 % \Rm^{-1}(\psi)
\newcommand{\Mdef}{\mathbf{M}}                 % defect matrix of the flatness proof
\newcommand{\Dil}{\mathsf{Dil}}                % \Dil_i(d) frame-B dilaton factors
\newcommand{\Edge}{\cE}                        % \Edge_{ij}(m,n) frame-B edge factors
\newcommand{\Ver}{\mathrm{Ver}}
\newcommand{\Edg}{\mathrm{Edg}}
\newcommand{\col}{\mathfrak{c}}                % colouring \col : \Ver(\Gamma) \to \{1,2\}
\newcommand{\GCont}{\mathrm{Cont}}             % \GCont(\Gamma,\col)
\newcommand{\Gsum}{\mathscr{G}}                % frame-B Hodge-vertex graph sum \Gsum_g
\newcommand{\Pij}{\cP}                         % \Pij_{ij}(z,w) = -\langle \vv_i(-z),\vv_j(-w)\rangle
\newcommand{\Kone}{\mathsf{K}_{1}}             % symplectic constant at idempotent 1
\newcommand{\Ktwo}{\mathsf{K}_{2}}
\newcommand{\brow}{\boldsymbol{\sigma}}        % \rho\vv_1(-\psi)
\newcommand{\Yv}{\mathbf{Y}}                   % t\brow
\newcommand{\Gk}{\mathsf{G}}                   % \Gk(z;t) = \rho\rr_{0,1}; \Gk_k = [t^k]\Gk
\newcommand{\gl}{\mathsf{g}}                   % \gl(\psi;t) = \Gk(-\psi;t)
\newcommand{\Dg}{\mathrm{D}\mathsf{g}}         % d/dt (t\gl)
\newcommand{\Rk}{\mathsf{V}}                   % \Rk_k = [t^k](z\Gk)
\newcommand{\Ebar}{\overline{\cE}}             % t^2\rho^2\Edge_{11}
\newcommand{\Nbar}{\overline{\mathsf{N}}}      % numerator of \Ebar
\newcommand{\Psh}{\mathfrak{p}}                % \Psh_k = \prod_{i\le k}((k+1)\alpha_2 - iE)((k+1)\wA - iE)
\newcommand{\Cj}{\mathsf{C}}                   % \Cj_j = \Psh_{j-1}/(j!(j-1)!E^{2j-4})
\newcommand{\Aleaf}{A}                         % A_j(t) leaf coefficient series
\newcommand{\Ahat}{\widehat{A}}                % t A_j

%-------------------------------------------------------------------------------
% Leaves and Hodge side
\newcommand{\lf}{\boldsymbol{\lambda}}         % \lf_j(\psi) = 1/(E - j\psi)
\newcommand{\Lsp}{\mathscr{L}}                 % \Lsp_m = span{\lf_1,...,\lf_m}
\newcommand{\qo}{\mathsf{q}_{\circ}}           % E - E^2\lf_1
\newcommand{\qt}{\mathsf{q}_{t}}               % dilaton series -\psi\rr_{0,1}(-\psi)
\newcommand{\Tdil}{T_{\mathrm{dil}}}           % \qo + \psi = -\psi^2/(E-\psi)
\newcommand{\Fpot}{\mathsf{F}^{\mathrm{pot}}}  % \Fpot_h
\newcommand{\dbase}{\mathfrak{d}}              % \dbase^s
\newcommand{\Dom}{\mathfrak{D}}                % neighbourhood of \qo
\newcommand{\Xiz}{\Xi}                         % \Xiz_h(\msM), scale E
\newcommand{\Xit}{\Xi^{\tau}}                  % scale-1 integral of [LP2]
\newcommand{\ZZ}{\mathrm{Z}}                   % \ZZ[h,n] = \Xiz_h(1^n)
\newcommand{\Mirr}{\mathsf{M}_{\mathrm{irr}}}  % \Mirr(h)
\newcommand{\dg}{\mathrm{dg}}                  % \dg(\msM) = \sum (j-1)
\newcommand{\msM}{\mathsf{M}}                  % multiset of integers >= 2
\newcommand{\msU}{\mathsf{U}}                  % multiset of positive integers
\newcommand{\tauw}{\tau_{\bw}}                 % -\epsilon_1/E, the scale-1 parameter of [LP2]
\newcommand{\st}{\sigma_{\!\bw}}               % -\epsilon_1\epsilon_2/E^2

%-------------------------------------------------------------------------------
% Three-torus and flop
\newcommand{\xc}{\mathsf{x}}                   % \alpha_1/\alpha_2 on \Tthree
\newcommand{\cc}{\mathsf{c}}                   % 1 + \xc = E/\alpha_2
\newcommand{\kap}{\varkappa}                   % \epsilon_1\epsilon_2\alpha_2^4/\alpha_1
\newcommand{\Zfl}{Z'}                          % the flop
\newcommand{\Xfl}{X'}


%-------------------------------------------------------------------------------
% Local macros of the section files of the long version
\providecommand{\pc}{\mathsf{p}}
\providecommand{\kE}{k_E}
\providecommand{\kph}{\kappa_{\mathrm{ph}}}
\providecommand{\Lt}{\Lnorm_{\tau}}
\providecommand{\Ngauss}{\Nlap_{\mathrm{G}}}
\providecommand{\Spht}{\mathsf{S}_{\phi,\tau}}
\providecommand{\Sphtt}{\mathsf{S}_{\phi,\tau\tau}}
\providecommand{\Abetam}[1]{\mathsf{A}_{\beta,#1}}
\providecommand{\Bfn}{\mathrm{B}}
\providecommand{\leaVn}{\mathbf{n}}
\providecommand{\leaPr}{\mathrm{Pr}}
\providecommand{\leaPhiF}{\boldsymbol{\Phi}_{\Fhyp}}
\providecommand{\leaLt}{\Lnorm_{2}}
\providecommand{\leaDt}{\mathsf{D}_{\stil}}
\providecommand{\leaDtau}{\mathsf{D}_{\stil,\tau}}
\providecommand{\leaNtw}{\mathsf{N}_{12}}
\providecommand{\leaDisc}{\Ddisc_{1}}
\providecommand{\TFsup}{\mathrm{F}}
\providecommand{\ideKt}{\widetilde{\bbK}_{\mathrm{CY}}}
\providecommand{\ideFz}{\bbF_{\circ}}
\providecommand{\ideone}{\mathbf{1}}
\providecommand{\ideV}{\mathcal{V}}
\providecommand{\idePF}{\mathsf{PF}}
\providecommand{\idekm}{\mathsf{k}_{\mathrm{m}}}
\providecommand{\ideCh}{\mathcal{X}}
\providecommand{\ideCurve}{\mathfrak{C}}
\providecommand{\idesec}{\varsigma}
\providecommand{\ideLam}{\boldsymbol{\Lambda}}
\providecommand{\idecan}{\boldsymbol{u}}
\providecommand{\ideS}{\mathscr{S}}
\providecommand{\idelg}[1]{\mathrm{a}^{#1}}
\providecommand{\ideP}{\mathfrak{P}}
\providecommand{\ideN}{\mathfrak{N}}
\providecommand{\ideG}{\mathfrak{G}}
\providecommand{\ideW}{\mathsf{W}_{\Psi}}
\providecommand{\ideRf}{\mathsf{R}^{\mathrm{f}}}
\providecommand{\ideY}{\mathsf{Y}}
\providecommand{\ideZ}{\mathsf{Z}}
\providecommand{\idesp}{\mathfrak{s}}
\providecommand{\ideQU}{\mathrm{V}}
\providecommand{\ideT}{\mathsf{T}_{\!\Rm}}
\providecommand{\regbra}[2]{\langle #1\mid #2\rangle}
\providecommand{\regD}{\dg}


%-------------------------------------------------------------------------------
% Title

\title[Refined GW theory of the resolved conifold on the CY torus (abridged)]{Refined Gromov--Witten theory of the resolved conifold on the Calabi--Yau torus: conifold regularity and rigidity}
\author{The GENUS collaboration}
\genusproject{solutions/resolved-conifold-cy-torus (abridged)}
\thanks{This is the abridged version of the long document of the same title (output/main.pdf of this GENUS project), referred to as [L] throughout. It states the results and sketches the arguments; complete proofs are in the long version.}

\begin{document}

\begin{abstract}
We study the equivariant Gromov--Witten theory of the resolved conifold times $\bbC^2$ on its Calabi--Yau four-torus, written as a Givental--Teleman graph sum with the normalised Hodge class of weights $(E,-\epsilon_1,-\epsilon_2)$ at every vertex. At the conifold point $\cQ=1$ one of the two idempotents is singular. We state, and sketch the proofs of, four results: every factor of the graph sum at the singular idempotent has a leaf structure, that is, after one rescaling the coefficient of $t^k$, $\cQ=\re^t$, is a combination of the functions $1/(E-j\psi)$ with $j\leq k+1$ (Theorem~A, at every non-degenerate weight point of the five-torus); the graph sum equals the Gromov--Witten potential on the Calabi--Yau torus (Theorem~B); every coloured stable graph with an edge is regular at the conifold point, the whole polar part being carried by the single-vertex graph with the value $\beta_{2g-2}(\epsilon_1,\epsilon_2)\,t^{2-2g}$ (Theorem~C); and for $g\geq2$ the genus-$g$ potential equals $\Pg\Li_{3-2g}(\cQ)/(\epsilon_1\epsilon_2)$ up to its constant term, that is, $\epsilon_1\epsilon_2\GW_{g,d}=\Pg d^{2g-3}$ for $d\geq1$ (Theorem~D). With genus $0$ and $1$ this proves the rigidity conjecture of \cite{BS} for the resolved conifold, for every non-zero class and every group acting through the Calabi--Yau torus. The inputs are Teleman's classification, the $R$-matrix action of Pandharipande--Pixton--Zvonkine, virtual localisation, Mumford's formula, Stirling's series, and two theorems on Hodge integrals from the companion paper \cite{LP2}; in genus one and degree $d\geq2$ the formula of \cite{BS} on the subtori $\alpha_i=\epsilon_j$ is used as an input. The constant-map term is not determined. This is the abridged version of the paper: it gives the statements and the key points of the arguments, and refers to the long version for the complete proofs.
\end{abstract}

\maketitle
\genusdisclaimer

\setcounter{tocdepth}{1}
\tableofcontents

% s1_intro (abridged version). Label prefix: int. The main theorems are stated here exactly as in the long version.
\section{Introduction}
\label{sec:int}

This paper is the abridged version of a longer paper with the same title, which we call \emph{the long version} and cite as [L]: a reference such as \Lref{Theorem~4.1} is to the statement with that number in the long version. The statements below are those of the long version; the proofs are replaced by sketches which give the mechanism of each step, and the proofs are in [L]. A few steps of [L] are marked there as written out in more detail than in its sources and still to be checked; they are listed in \cref{sec:evi}.

%-------------------------------------------------------------------------------
\subsection{The setting}
\label{sec:int-setting}

Let $X$ be the resolved conifold, the total space of $\cO(-1)\oplus\cO(-1)\to\bbP^1$, and let $Z=X\times\bbC^2$. The five-torus $\Tmax=(\bbC^*)^5$ acts on $Z$ with two fixed points $p_\pm$ on the zero section, with tangent weights $(\alpha_0,\alpha_1,\alpha_2,-\epsilon_1,-\epsilon_2)$ at $p_+$ and $(-\alpha_0,\alpha_0+\alpha_1,\alpha_0+\alpha_2,-\epsilon_1,-\epsilon_2)$ at $p_-$: $\alpha_0$ is the weight of the tangent line of $\bbP^1$, $\alpha_1,\alpha_2$ are the weights of the two $\cO(-1)$ summands, and $-\epsilon_1,-\epsilon_2$ those of the two trivial directions. The \emph{Calabi--Yau torus} $\TCY\subset\Tmax$ is the four-torus on which the five tangent weights sum to zero,
\[
E\coloneqq\alpha_0+\alpha_1+\alpha_2=\epsilon_1+\epsilon_2 ,
\]
and we write $\wA\coloneqq\alpha_0+\alpha_2$, $\wB\coloneqq\alpha_0+\alpha_1$. The equivariant Gromov--Witten invariants $\GW_{g,d}$ of $Z$ in genus $g$ and degree $d[\bbP^1]$ are defined, as in \cite{BS}, by virtual localisation on the proper fixed loci of the moduli spaces of stable maps to $\bbP^1$ (\cref{def:set-inv}); on $\TCY$ they are elements of the function field $\KCY=\bbQ(\alpha_0,\alpha_1,\alpha_2,\epsilon_1)$, $\epsilon_2=E-\epsilon_1$. The genus-$g$ series is
\[
\sF_g(\cQ)=\sum_{d\geq0}\GW_{g,d}\,\cQ^d ,\qquad \cQ=\re^{t},
\]
with the constant-map term $d=0$ included; the series of \cite{BS} runs over $d\geq1$, and the difference is the constant $\sF_g(0)$. The \emph{conifold point} is $\cQ=1$, that is $t=0$, where the two critical points of the superpotential governing the genus-zero theory collide with a logarithmic singularity.

The tool is the Givental--Teleman expression of the higher-genus theory as a sum over stable graphs. On $\TCY$ the point $\cQ=0$ is semisimple (because $\alpha_0\neq0$), the equivariant classes of $Z$ form a semisimple cohomological field theory of rank two, and Teleman's classification writes it as the action of an $R$-matrix on the topological field theory with vertex weights $\Delta_i^{g_v-1}$, $i\in\{1,2\}$. The two idempotents behave very differently at the conifold point: one, the \emph{singular} idempotent (colour $1$), has $\Delta_1\sim-\epsilon_1\epsilon_2E^3/t^2$, the other, the \emph{regular} idempotent (colour $2$), has $\Delta_2$ analytic and non-zero at $t=0$. We write the graph sum in a particular frame, \emph{frame~B}: the $R$-matrix is multiplied on the right by the scalar series $\hw^{-1}$, $\hw=h_{(E,-\epsilon_1,-\epsilon_2)}$, which by Mumford's formula transfers the normalised Hodge class
\[
\Hg_{\bw}=\prod_{a=1}^{3}\frac{\Lamv(\bw_a)}{\bw_a^{g_v}},\qquad \bw=(E,-\epsilon_1,-\epsilon_2),\qquad \Lamv(u)=\sum_i(-1)^i\lambda_iu^{g_v-i},
\]
onto every vertex. The resulting \emph{frame-B Hodge-vertex graph sum} $\Gsum_g$ (\cref{def:set-graphsum}) is a sum over connected stable graphs of genus $g$ without legs, over colourings of their vertices by $\{1,2\}$, and over decorations, of products of vertex weights, dilaton factors and edge factors, all built from two rows $\rr_{0,i}(z;t)$, $\rr_{1,i}(z;t)$ per idempotent. The class $\bw$ is forced: it is the unique class, up to pairs $(v,-v)$, for which the singular row takes the value $E/(E+z)$ at the conifold point (\cref{lem:set-hmult}, \cref{cor:row-F2}). The rows themselves are normalised stationary-phase expansions of the Euler-type integrand $\re^{\cW/z}\dd U/U$ at the two critical points of the superpotential $\cW$, multiplied by one constant $\hnw=h_{(\alpha_1,\wA,-\epsilon_1,-\epsilon_2)}$ fixed at $\cQ=0$.

%-------------------------------------------------------------------------------
\subsection{Statements}
\label{sec:int-statements}

The notation is that of \cref{sec:set}: the non-degeneracy condition (ND) is $\alpha_1\alpha_2\wA\wB\epsilon_1\epsilon_2E\neq0$ (\cref{def:set-ND}), a \emph{weight point with (ND)} is a point of $\mathrm{Lie}(\TCY)$ at which (ND) holds, $\alpha_0=0$ being allowed, and ``at the generic point of $\TCY$'' means ``in $\KCY$''. The frame-B rows and factors are those of \cref{def:set-rows,def:set-factors}; $\rho(t)=(\Delta_1t^2/(-\epsilon_1\epsilon_2E^3))^{-1/2}$ with $\rho(0)=1$; the \emph{leaves} are $\lf_j(\psi)=1/(E-j\psi)$, $j\geq1$, the \emph{leaf spaces} are $\Lsp_m=\mathrm{span}(\lf_1,\dots,\lf_m)$, $\qo=E-E^2\lf_1$, $\qt(\psi)=\sum_{d\geq1}\Dil_1(d)\psi^d$ is the dilaton series of the singular idempotent, and $f_{j,b}(\psi)=\coef{w^b}\Edge_{1j}(\psi,w)$ are the legs (\cref{def:set-rho}); $\tau=1-1/\cQ$, and $\delta$, $\sast$, $\Lv_2$, $\alg$ are the algebraic functions of $\tau$ of \cref{def:set-alg}, \cref{lem:set-eig} and \cref{lem:row-hess}. The first theorem is the structural statement from which everything else follows.

\begin{thm}[{Theorem~A: leaf structure at generic weights; \Lref{Theorems~1.1 and~4.1}}]\genusstatus{AI-proposed, unchecked}
\label{thm:int-leaf}
Let $(\alpha_0,\alpha_1,\alpha_2,\allowbreak\epsilon_1,\epsilon_2)$ be a weight point of $\TCY$ with (ND), or indeterminates with $E=\epsilon_1+\epsilon_2$ imposed. All factors are the frame-B factors at their true idempotents. Then:
\begin{description}
\item[(F1)] $\Delta_1=-\epsilon_1\epsilon_2E^3/(t^2\rho^2)$, where $\rho=\tau/(t\alg)$ is analytic at $t=0$ with $\rho(0)=1$, so that $\Delta_1=-(\epsilon_1\epsilon_2E^3/t^2)(1-\uq t/E^2+O(t^2))$, $\uq=(E-2\alpha_1)(E-2\alpha_2)$; $\Delta_2=\epsilon_1\epsilon_2\wB\delta\sast\Lv_2$ is analytic at $t=0$ with $\Delta_2(0)=\epsilon_1\epsilon_2\alpha_1\alpha_2\wA\wB/E\neq0$ and $\Delta_2'(0)=0$.
\item[(F2)] $\rr_{0,1}(z;0)=E/(E+z)$, and $\rho\,\rr_{1,1}(z;t)=E/t+\alpha_1\alpha_2/(E+z)+O(t)$.
\item[(F3)] Every frame-B dilaton factor $\Dil_i(d)$ ($i=1,2$, $d\geq1$) and every edge factor $\Edge_{ij}(m,n)$ ($i,j\in\{1,2\}$, $m,n\geq0$) is analytic at $t=0$.
\item[(L1)] $\coef{t^0}(\rho\qt)=\qo=E-E^2\lf_1$, and $\coef{t^k}(\rho\qt)\in\Lsp_{k+1}$ for every $k\geq1$.
\item[(L2)] For $j\in\{1,2\}$, every $b\geq0$ and every $k\geq0$, $\coef{t^k}(\rho f_{j,b})\in\Lsp_{k+1}$. Moreover $\Edge_{21}(b,a)=\Edge_{12}(a,b)$, and $\coef{t^k}f_{j,b}\in\Lsp_{k+1}$ as well.
\item[(L3)] For every $k\geq0$, $\coef{t^k}(\rho^2\Edge_{11}(\psi,\psi'))\in\mathrm{span}\{\lf_i(\psi)\lf_l(\psi'):i,l\geq1,\ (i-1)+(l-1)\leq k\}$, and $\Edge_{11}(a,b)=\Edge_{11}(b,a)$. The same support bound holds for $\coef{t^k}(\rho^e\Edge_{11})$ for every integer $e$.
\end{description}
Every membership is an identity of formal power series in $\psi$ (and $\psi'$), valid at every $\psi$-order and every $t$-order; the coefficients are rational functions of the weights, regular on (ND). The bounds in (L1) and (L3) are sharp at generic weights.
\end{thm}

In words: at the conifold point the vertex weight of the singular idempotent has a double pole with an exactly known unit, every factor is regular, and the $t$-expansion of every rescaled factor at the singular idempotent is \emph{leaf-structured}: the coefficient of $t^k$ is a combination of the simple-pole functions $1/(E-j\psi)$ with $j\leq k+1$, and of their products for loops, with no constant term. The proof is sketched in \cref{sec:row,sec:lea}, where a closed form of every leaf coefficient is also given (\cref{thm:lea-closed}). The proof does not use the Calabi--Yau condition: the frame-B rows depend on $(\alpha_0,\alpha_1,\alpha_2)$ only, and $\epsilon_1\epsilon_2$ enters the factors as an overall constant, so the theorem holds at every weight point of the five-torus with (ND), with $\bw=(E,-\epsilon_1,-\epsilon_2)$ and $E=\alpha_0+\alpha_1+\alpha_2$ (\cref{rmk:lea-cy}). The Calabi--Yau condition enters the main chain once, in Theorem~C below, where it makes $\bw$ a sum-zero triple whose first weight is the leaf scale $E$.

\begin{thm}[{Theorem~B: the graph sum is the potential; \Lref{Theorem~1.2 and Proposition~5.1}}]\genusstatus{AI-proposed, unchecked}
\label{thm:int-ident}
Let $g\geq2$. At the generic point of $\TCY$, $\sF_g$ lies in $\KCY(\cQ)$, is regular at $\cQ=0$ with expansion $\sum_{d\geq0}\GW_{g,d}\cQ^d$ there, and
\[
\sF_g=\Gsum_g ,
\]
the frame-B Hodge-vertex graph sum \eqref{eq:set-Gsum} with vertex class $\Hg_{\bw}$, $\bw=(E,-\epsilon_1,-\epsilon_2)$, and the frame-B factors of \cref{def:set-rows,def:set-factors}, in which $\RB=\RZ\hw^{-1}$ and $\RZ$ has the rows $\hnw\rpt_{0,i}$, $\nw=(\alpha_1,\wA,-\epsilon_1,-\epsilon_2)$. Each coloured term $\GCont(\Gamma,\col)$ is an element of the function field $\KCY(\Ccov)$ of the double cover $\ydc^2=\Disc(\cQ)$ of $\bbP^1_\cQ$, regular over every $\cQ\notin\{0,1,\infty\}$ with $\Disc(\cQ)\neq0$.
\end{thm}

The proof is sketched in \cref{sec:ide}. It is an identity in $\KCY(\cQ)$, proved at the generic point of $\TCY$ without any specialisation: the field $\KCY$ embeds in $\bbC$, Teleman's classification applies there, and the $R$-matrix descends to $\KCY\lser{\cQ}$.

\begin{thm}[{Theorem~C: regularity and the polar part; \Lref{Theorem~1.3 and Corollary~6.13}}]\genusstatus{AI-proposed, unchecked}
\label{thm:int-reg}
Fix $g\geq2$. In frame B with $\bw=(E,-\epsilon_1,-\epsilon_2)$, at the generic point of $\TCY$ or at any weight point with (ND):
\begin{enumerate}
\item[(a)] every coloured stable graph $(\Gamma,\col)$ of genus $g$ with at least one edge is regular at $t=0$ after summing over decorations;
\item[(b)] the single-vertex graph is regular at colour $2$, and at colour $1$ its polar part is exactly
\[
\gapg\,t^{2-2g},\qquad \gapg\coloneqq(-\epsilon_1\epsilon_2E^3)^{g-1}\ZZ[g,0]=\beta_{2g-2}(\epsilon_1,\epsilon_2)=\frac{(2g-3)!\,\Pg}{\epsilon_1\epsilon_2},
\]
with no polar term of order $t^{3-2g},\dots,t^{-1}$; here $\ZZ[g,0]$ is the Hodge integral of \cref{def:set-fpot}, and
\[
\Pg=\coef{t^{2g}}\frac{\epsilon_1\epsilon_2\,t^2}{4\sinh(\epsilon_1t/2)\sinh(\epsilon_2t/2)},\qquad
\beta_n(\epsilon_1,\epsilon_2)=(n-1)!\,\coef{t^n}\frac{1}{4\sinh(\epsilon_1t/2)\sinh(\epsilon_2t/2)} ;
\]
\item[(c)] hence the expansion of $\Gsum_g$ in $\KCY\lser t$ has polar part exactly $\gapg t^{2-2g}$, and, by \cref{thm:int-ident}, so has the image of $\sF_g$ under $\cQ=\re^t$.
\end{enumerate}
\end{thm}

The proof is sketched in \cref{sec:reg}. Parts (a), (b) and the first half of (c) are statements about the decoration sums as defined in \cref{def:set-graphsum}; they rest on \cref{thm:int-leaf}, on the two theorems of the companion paper \cite{LP2} quoted as \cref{thm:reg-L,thm:reg-VI}, and on the string and dilaton equations. The whole polar part of the graph sum sits in the single-vertex graph at the singular colour.

\begin{thm}[{Theorem~D: rigidity; \Lref{Theorems~1.4 and~7.1}}]\genusstatus{AI-proposed, unchecked}
\label{thm:int-rig}
For $g\geq2$, in $\KCY(\cQ)$,
\[
\sF_g(\cQ)-\sF_g(0)=\frac{\Pg\,\Li_{3-2g}(\cQ)}{\epsilon_1\epsilon_2},\qquad\text{that is}\qquad \epsilon_1\epsilon_2\,\GW_{g,d}=\Pg\,d^{2g-3}\quad\text{for every }d\geq1,
\]
with $\Li_{-n}(\cQ)=\sum_{d\geq1}d^n\cQ^d$. Both sides of the second identity are polynomials in the weights on $\mathrm{Lie}(\TCY)$, so it holds at every point of $\mathrm{Lie}(\TCY)$, including $\alpha_1=\alpha_2$, $E=0$, $\epsilon_1=0$ or $\epsilon_2=0$, and $\alpha_0=0$; in particular $\epsilon_1\epsilon_2\GW_{g,d}$ depends on $(\epsilon_1,\epsilon_2)$ only, symmetrically and through $\epsilon_1^2,\epsilon_2^2$.
\end{thm}

The proof is sketched in \cref{sec:rig}. In genus $0$ one has $\epsilon_1\epsilon_2\GW_{0,d}=d^{-3}$ at every weight, by the Aspinwall--Morrison formula (\cref{prop:rig-g0}), and in genus $1$ one has $\epsilon_1\epsilon_2\GW_{1,d}=\Ph{1}/d=-(\epsilon_1^2+\epsilon_2^2)/(24d)$ on $\TCY$ (\cref{lem:rig-g1}); the genus-one statement is proved by direct localisation for $d=1$, and for $d\geq2$ it uses, as an input, the formula of \cite{BS} for the resolved conifold on the subtori on which $\alpha_i=\epsilon_j$. Together:

\begin{cor}[{the rigidity conjecture of \cite{BS} for the resolved conifold; \Lref{Corollaries~1.5 and~7.13}}]\genusstatus{AI-proposed, unchecked}
\label{cor:int-conj}
For every $d\geq1$ and every $g\geq0$, $\epsilon_1\epsilon_2\,\GW_g(Z,\TCY;d[\bbP^1])=\Pg\,d^{2g-3}$, so that
\[
\sum_{g\geq0}\GW_g(Z,\TCY;d[\bbP^1])=\frac1d\Big(2\sinh\frac{d\epsilon_1}2\Big)^{-1}\Big(2\sinh\frac{d\epsilon_2}2\Big)^{-1}.
\]
Hence the rigidity conjecture of \cite{BS} holds for $X$ the resolved conifold, every non-zero class $d[\bbP^1]$, and the group $\TCY$; by contravariance it holds for every subtorus of $\TCY$ on which the invariants are defined by localisation (the normal bundle $\Nb$ moving on the fixed loci) and for every connected group whose maximal torus acts on $Z$ through such a subtorus, in particular for every Calabi--Yau subtorus of the five-torus with this property.
\end{cor}

The three-torus $\Tthree=\{\alpha_0=0\}\cap\TCY$, which contains the fibrewise torus $\TF$ ($\alpha_i=\epsilon_i$) on which the zero section is fixed pointwise, is not covered by the proof of \cref{thm:int-ident} directly, because $\cQ=0$ is not semisimple there; Theorems~A and~C hold on $\Tthree$ by a rescaling of the rows to the fibrewise ones and a ring map applied to the identity of \cref{thm:int-ident} (\cref{cor:rig-3T}), and rigidity on $\Tthree$ is the restriction of \cref{thm:int-rig} by polynomiality.

%-------------------------------------------------------------------------------
\subsection{Inputs}
\label{sec:int-inputs}

The results are conditional only on the following cited theorems, each named at the place where it enters.
\begin{enumerate}
\item Stirling's series for $\log\Gamma$ with remainder \cite[5.11.1 and Section~5.11(ii)]{DLMF}, together with the Beta integral \cite[5.12.1]{DLMF} and Laplace's method, which is proved in full in \Lref{Lemma~3.11}. This is the only analytic input of Theorem~A; it enters through the normalised expansion of the Beta phase (\cref{lem:row-betaphase}, \cref{prop:row-beta}).
\item Teleman's classification of semisimple cohomological field theories \cite[Theorem~(2), (1.15), Prop.~8.7(ii)]{Tel12}, with the stability range of \cite[2.17--2.19, (2.15)]{Tel12} as reported there from \cite{Loo96,MW07,Har85,Iva90}; the $R$-matrix action and its group property \cite[Defs.~2.2, 2.13, Prop.~2.14]{PPZ15}; the virtual localisation formula \cite{GP99}; the twisted-theory axioms \cite[App.~1]{CG07}; Mumford's relation and formula \cite{Mum83}, \cite[Section~1.1]{FP00}; the non-vanishing of $\lambda_g\lambda_{g-1}\lambda_{g-2}$ \cite[Theorem~4]{FP00}; and the Aspinwall--Morrison formula \cite[p.~3]{FP00}. These enter Theorem~B; the localisation formula \cite{GP99} enters also Theorem~D (\cref{lem:rig-poly}) and genus $1$ (\cref{lem:rig-g1}), Mumford's relation enters genus $1$, and the Aspinwall--Morrison formula is the genus-zero statement (\cref{prop:rig-g0}). The locators in \cite{Tel12} are those of the arXiv version; they have not been compared with the journal numbering.
\item The two theorems of the companion paper \cite{LP2}: the vanishing theorem for the Hodge integrals $\Xiz_h(\msM)$ with the class $\Hg_{(1,\tau,-1-\tau)}$ at leaf degree $\dg(\msM)\leq2h-3$ (\cite[Theorem~8.1, Corollaries~8.2--8.4]{LP2}, quoted as \cref{thm:reg-L}), and the value of the single vertex, $(-\epsilon_1\epsilon_2E^3)^{g-1}\ZZ[g,0]=\beta_{2g-2}(\epsilon_1,\epsilon_2)$ (\cite[Theorem~9.1]{LP2}, quoted as \cref{thm:reg-VI}). Both rest on the Mari\~no--Vafa formula \cite[(1)--(3)]{LLZ03}. They enter Theorem~C only. The numbering of the statements of \cite{LP2} quoted here is provisional.
\item For genus $1$ and $d\geq2$ only, the formula of \cite{BS} for the resolved conifold on the subtori on which $\alpha_i=\epsilon_j$ (\cref{lem:rig-g1}).
\end{enumerate}
Which result uses which: Theorem~A uses (i); Theorem~C uses Theorem~A and (iii); Theorem~B uses (i), through the normalisation constant at $\cQ=0$, and (ii); Theorem~D uses Theorems~B and~C and, in the polynomiality lemma, \cite{GP99}; genus $0$ uses the Aspinwall--Morrison formula; genus $1$ uses \cite{GP99}, Mumford's relation and, for $d\geq2$, (iv). None of the arguments uses topological recursion, a mirror theorem, a degeneration argument or a finite-order computation; the one input proved by degeneration is (iv), used in genus $1$ for $d\geq2$ only; the finite-order checks that accompanied the work are summarised in \cref{sec:evi}, and no proof step refers to them.

%-------------------------------------------------------------------------------
\subsection{What is not claimed}
\label{sec:int-notclaimed}

\begin{enumerate}
\item Nothing is claimed off $\TCY$. Theorem~A holds at every weight point of the five-torus with (ND), but Theorem~C does not: off $\TCY$ the class $\bw$ does not sum to zero, the vanishing theorem does not apply, and the coloured graphs with an edge acquire poles at $t=0$ (\cref{sec:evi}, as evidence only).
\item The constant-map term $\sF_g(0)=\GW_{g,0}$ is not determined. It is a rational function of all five weights \Lref{Remark~7.5}, it depends on $\alpha$, as its computed values in genus $2,3,4$ show \Lref{Remark~7.16}, and it is not part of the conjecture, which concerns non-zero classes. The only statement about it proved in [L] is the flop identity $\sF_g(\infty)=\sF_g(0)$ (\cref{prop:rig-flop}).
\item The conjecture for group actions on $Z$ that are not conjugate into the five-torus is not addressed.
\item The proofs of the cited theorems of \cref{sec:int-inputs} are not reproduced, neither here nor in [L]. In particular the two theorems of \cite{LP2} are quoted, with the scale conversion written out (\cref{lem:reg-scale}).
\end{enumerate}

%-------------------------------------------------------------------------------
\subsection{Plan and dependencies}
\label{sec:int-strategy}

\Cref{sec:set} fixes the genus-zero data, the rows and factors of the frame-B graph sum, the leaf spaces and the Hodge-side integrals. \Cref{sec:row} computes the singular row in closed form: in the Euler coordinate $U=1-\tau s$ the two critical points stay separated at $t=0$, the point-frame row is a normalised formal Laplace expansion which satisfies Gauss's hypergeometric equation in $\tau$, and its constant is a Beta integral evaluated by Stirling's series. \Cref{sec:lea} sketches the proof of Theorem~A: the closed form gives (L1) by partial fractions, the quantum differential equation and the self-adjointness of the quantum product give a first-order flow for the loop $\Edge_{11}$ with a rank-one right side of leaf series, hence (L3), and the same calculus at the second critical point gives (L2). \Cref{sec:ide} sketches the proof of Theorem~B: Teleman's $R$-matrix over $\KCY$ is flat, the rows form a second flat symplectic solution, and the two agree once the constants are fixed at $\cQ=0$. \Cref{sec:reg} quotes the two theorems of \cite{LP2} and sketches the proof of Theorem~C: dilaton homogeneity rescales each singular vertex, the leaf structure turns its expansion into Hodge integrals with leaves, and the vanishing theorem kills every polar term except that of the single vertex. \Cref{sec:rig} sketches the proof of Theorem~D, by locating the possible poles of the rational function $\sF_g$ and comparing with the polylogarithm at $\cQ=1$, together with genus $0$ and $1$, \cref{cor:int-conj}, the flop identity and the three-torus. \Cref{sec:evi} summarises the finite-order checks and what remains conditional.

\emph{Dependencies.} \Cref{sec:row} uses \cref{sec:set}; \cref{sec:lea} uses \cref{sec:set,sec:row}; \cref{sec:ide} uses \cref{sec:set} and the Laplace calculus and Beta-phase lemma of \cref{sec:row}, not Theorem~A; \cref{sec:reg} uses \cref{sec:set}, Theorem~A and the cited theorems of \cite{LP2}, and uses Theorem~B only for the second half of \cref{thm:int-reg}(c); \cref{sec:rig} uses Theorems~B and~C, the genus-zero data of \cref{sec:set}, the localisation formula \cite{GP99} and, for the flop identity and the three-torus only, \cref{lem:row-betaphase} and \cref{cor:lea-a0}. Teleman's theorem and the other inputs of item (ii) of \cref{sec:int-inputs} enter only through \cref{sec:ide}.

%-------------------------------------------------------------------------------
\subsection{Relation to the results of \texorpdfstring{\cite{BS}}{[BS]}}
\label{sec:int-BS}

The paper \cite{BS} defines the refined invariants by localisation, formulates the rigidity conjecture for a rigid class on $X\times\bbC^2$ under a Calabi--Yau group action, proves the case in which the action on $X$ is itself Calabi--Yau, and proves for the resolved conifold the closed formula of \cref{cor:int-conj} on the subtori on which $\alpha_i=\epsilon_j$, by degeneration of the base curve and the relative invariants of local curves; it also verifies the formula on $\TCY$ by computer in low degree and genus. What is new here is the statement on all of $\TCY$, for every $d\geq1$ and every genus $g\geq2$, genus by genus, and the method: for $g\geq2$ no degeneration is used, and the result comes from the structure of the Hodge-vertex graph sum at the conifold point, that is, from Theorems~A, B and~C. The genus-one input of \cite{BS} at $d\geq2$ is the one place where a result of \cite{BS} is used.

\subsection*{Acknowledgements}

This work was supported by the UKRI grant UKRI2394 and by the Leverhulme Trust grant LIP-2023-007.

% s2_setup_a (abridged version). Label prefix: set. Geometry, invariants, genus-zero data.
\section{Set-up: the geometry, the torus and the genus-zero data}
\label{sec:set}

Everything in this section is a definition or an elementary identity; the verifications are direct computations, given in \Lref{Section~2}. Throughout, idempotent (colour) $1$ is the singular one and idempotent (colour) $2$ the regular one, in the sense of \cref{def:set-alg}.

%-------------------------------------------------------------------------------
\subsection{The target, the tori and the invariants}
\label{sec:set-target}

Let $X$ be the resolved conifold and $Z\coloneqq X\times\bbC^2$, the total space of $\cO(-1)^{\oplus2}\oplus\cO^{\oplus2}\to\bbP^1$. The five-torus $\Tmax=(\bbC^*)^5$ acts on $Z$ with the two fixed points $p_+,p_-$ of the zero section, with tangent weights
\begin{equation}
\label{eq:set-wts}
\wts_+=(\alpha_0,\alpha_1,\alpha_2,-\epsilon_1,-\epsilon_2),\qquad
\wts_-=(-\alpha_0,\alpha_0+\alpha_1,\alpha_0+\alpha_2,-\epsilon_1,-\epsilon_2) ,
\end{equation}
and the $\Tmax$-linearised normal bundle of the zero section is $\Nb=\cO(-1)_{\alpha_1}\oplus\cO(-1)_{\alpha_2}\oplus\cO_{-\epsilon_1}\oplus\cO_{-\epsilon_2}$. We write
\[
E\coloneqq\alpha_0+\alpha_1+\alpha_2,\qquad \wA\coloneqq\alpha_0+\alpha_2=E-\alpha_1,\qquad \wB\coloneqq\alpha_0+\alpha_1=E-\alpha_2,\qquad \uq\coloneqq(E-2\alpha_1)(E-2\alpha_2),
\]
so that $\wA\wB=\alpha_1\alpha_2+\alpha_0E$ and $2\alpha_0E+4\alpha_1\alpha_2=E^2+\uq$. The products of the tangent weights are $\eul(\wts_+)=\alpha_0\alpha_1\alpha_2\epsilon_1\epsilon_2$ and $\eul(\wts_-)=-\alpha_0\wA\wB\epsilon_1\epsilon_2$.

\begin{defn}[{\Lref{Definition~2.1}: the Calabi--Yau torus; non-degeneracy}]
\label{def:set-ND}
The \emph{Calabi--Yau four-torus} $\TCY\subset\Tmax$ is the subtorus on which $E=\alpha_0+\alpha_1+\alpha_2=\epsilon_1+\epsilon_2$. Inside $\TCY$ we use the three-torus $\Tthree:\ \alpha_0=0$ and the \emph{fibrewise torus} $\TF:\ \alpha_0=0,\ \alpha_i=\epsilon_i$, which fixes the zero section pointwise. The \emph{non-degeneracy condition} is
\[
\text{(ND)}\qquad \alpha_1\alpha_2\wA\wB\epsilon_1\epsilon_2E\neq0 .
\]
A \emph{weight point with (ND)} is a point of $\mathrm{Lie}(\TCY)$, with complex (or rational) coordinates, at which (ND) holds; $\alpha_0=0$ is allowed. The function fields in use are $\Kal\coloneqq\bbQ(\alpha_0,\alpha_1,\alpha_2,\epsilon_1,\epsilon_2)$, $\KCY\coloneqq\bbQ(\alpha_0,\alpha_1,\alpha_2,\epsilon_1)$ with $\epsilon_2\coloneqq E-\epsilon_1$, and $\Kred\coloneqq\bbQ(\alpha_1,\alpha_2,\epsilon_1)$, so that $\KCY=\Kred(\alpha_0)$; $\Kred$ regarded as the function field of $\Tthree$ is written $\Kthree$. ``At the generic point of $\TCY$'' means ``in $\KCY$'', where (ND) holds and $\alpha_0$, $\alpha_1-\alpha_2$ are non-zero.
\end{defn}

\begin{defn}[{\Lref{Definition~2.3}: invariants by localisation; the series}]
\label{def:set-inv}
For $g\geq2$ and $d\geq0$ let $\GW_{g,d}\coloneqq\GW^{Z,\TCY}_{g,d[\bbP^1]}\in\KCY$ be the genus-$g$, degree $d[\bbP^1]$ $\TCY$-equivariant Gromov--Witten invariant of $Z$ in the definition by localisation: the $\TCY$-fixed loci of $\Mbar_{g,0}(Z,d)$ are the $\alpha_0$-fixed loci of the proper moduli space $\Mbar_{g,0}(\bbP^1,d)$ (as $\alpha_0\neq0$ in $\KCY$), and the invariant is the sum over them of the residues of the virtual localisation formula \cite{GP99}, with $\Nb$ moving. For the five-torus the same definition gives $\GW^{\Tmax}_{g,d}\in\Kal$. The genus-$g$ series is
\[
\sF_g(\cQ)\coloneqq\sum_{d\geq0}\GW_{g,d}\,\cQ^d,\qquad \sG_g\coloneqq\sF_g-\sF_g(0)=\sum_{d\geq1}\GW_{g,d}\,\cQ^d .
\]
We put $t\coloneqq\log\cQ$ and $\tau\coloneqq1-1/\cQ=1-\re^{-t}$, so that $\cQ\tau=\cQ-1$ and $\dd\tau/\dd t=1-\tau$. The \emph{conifold point} is $\cQ=1$, that is, $t=\tau=0$. Finally, for $g\geq0$ and $n\geq1$,
$\Li_{-n}(\cQ)\coloneqq\sum_{d\geq1}d^n\cQ^d$ and
\[
\Pg\coloneqq\coef{t^{2g}}\,\frac{\epsilon_1\epsilon_2\,t^2}{4\sinh(\epsilon_1t/2)\sinh(\epsilon_2t/2)},\qquad
\beta_n(\epsilon_1,\epsilon_2)\coloneqq(n-1)!\,\coef{t^n}\,\frac{1}{4\sinh(\epsilon_1t/2)\sinh(\epsilon_2t/2)} .
\]
\end{defn}

On functions of $t$ a prime is $\dd/\dd t=\thQ=\cQ\,\dd/\dd\cQ$ and $f_\tau\coloneqq\dd f/\dd\tau$, so $f'=(1-\tau)f_\tau$; on functions of the Euler variable $s$ below a prime is $\dd/\dd s$. Further $\Dz\coloneqq z\,\dd/\dd t$, and $\coef{t^k}$, $\coef{z^m}$, $\coef{\psi^a}$ denote coefficients, taken coefficientwise in the other variables.

%-------------------------------------------------------------------------------
\subsection{The superpotential and the Frobenius data}
\label{sec:set-crit}

\begin{lem}[{\Lref{Lemma~2.5}: superpotential, critical points; parts omitted}]\genusstatus{AI-proposed, unchecked}
\label{lem:set-crit}
Let $\thU\coloneqq U\,\dd/\dd U$, $\cW(U)\coloneqq\alpha_1\log(1-\cQ U)-\alpha_2\log U+\wA\log(1-U)$ and $\cW_k\coloneqq\thU^k\cW$ for $k\geq1$; the $\cW_k$ are rational in $U$ and $\cQ$, and for $k\geq2$ they are polynomials without constant term in $(1-U)^{-1}$ and $(1-\cQ U)^{-1}$ with coefficients in $\bbZ[\alpha_1,\alpha_2,\wA]$.
\begin{enumerate}
\item The critical equation $U\cW_U=0$ is equivalent to
\begin{equation}
\label{eq:set-P}
\Pq(U)\coloneqq\wB\cQ U^2-\big(\alpha_0+\cQ(\alpha_1-\alpha_2)\big)U-\alpha_2=0 ,
\end{equation}
and $(1-\cQ U)(1-U)\,U\cW_U=\Pq(U)$ identically. The discriminant of $\Pq$ in $U$ is $\Disc(\cQ)$ of \cref{lem:set-disc}.
\item $\Pq(1)=\wA(\cQ-1)$ and $\Pq(1/\cQ)=-\alpha_1\tau$. So, for $\cQ$ finite and non-zero, a root $U_i$ of $\Pq$ satisfies $U_i=1$ or $\cQ U_i=1$ only over $\cQ=1$.
\item Put $\Lv(U)\coloneqq-\alpha_1\cQ U/(1-\cQ U)=\de_t\cW$. At a critical point $\Lv(U)=\alpha_2+\wA U/(1-U)$, and in $\KCY(\cQ)[U]/(\Pq)$ one has $\Lv(U)=(\wB U+\alpha_2)/\tau$.
\end{enumerate}
\end{lem}

\begin{prop}[{\Lref{Proposition~2.6}: Frobenius data in the flat basis}]\genusstatus{AI-proposed, unchecked}
\label{prop:set-frob}
Let $\{1,H\}$ be the flat basis, $H$ the equivariant hyperplane class of $\bbP^1$ with $H|_{p_+}=0$, $H|_{p_-}=-\alpha_0$. The genus-zero potential of the $\TCY$-equivariant theory of $Z$ is $\mathsf{Prep}^{\mathrm{cl}}+\Li_3(\re^{t_1})/(\epsilon_1\epsilon_2)$ with $\mathsf{Prep}^{\mathrm{cl}}$ the classical cubic term (\cref{lem:ide-g0}). The inverse of the metric is
\[
\eta^{-1}=\epsilon_1\epsilon_2\begin{pmatrix}\alpha_0\alpha_1\alpha_2&\alpha_1\alpha_2\\\alpha_1\alpha_2&-E\end{pmatrix},
\qquad\text{independent of }t;
\]
$1$ is the unit of the quantum product, and $H\qprod H=-(\alpha_1\alpha_2/\tau)\,1+(E/\tau-\alpha_0)H$. The eigenvalues $\Lv_1,\Lv_2$ of $H\qprod$ are the roots of $\Lv^2-(E/\tau-\alpha_0)\Lv+\alpha_1\alpha_2/\tau$, and the product is commutative, associative and Frobenius. Covectors $a=(a_0,a_1)$ (components paired with $1$ and $H$) are paired by $\eta^{-1}$,
\begin{equation}
\label{eq:set-pairing}
\langle a,b\rangle\coloneqq\epsilon_1\epsilon_2\big[\alpha_0\alpha_1\alpha_2\,a_0b_0+\alpha_1\alpha_2(a_0b_1+a_1b_0)-E\,a_1b_1\big],
\end{equation}
symmetric with constant coefficients.
\end{prop}

\begin{defn}[{\Lref{Definition~2.7}: the quadratic algebra and its involution}]
\label{def:set-QA}
Let $\QA$ be the ring $\KCY(\cQ)[U]/(\Pq(U))$, with its two characters $\ev_i:U\mapsto U_i$, which are jointly injective. The \emph{involution} of $\QA$ is the ring automorphism $U\mapsto U_1+U_2-U$, which exchanges the two characters. By \cref{lem:set-disc}, $\QA$ is the function field $\KCY(\Ccov)$ of the double cover $\Ccov\to\bbP^1_\cQ$, $\ydc^2=\Disc(\cQ)$, and the involution is its deck involution; $\Lv(U)\in\QA$ and $\ev_i\Lv(U)=\Lv_i$.
\end{defn}

\begin{prop}[{\Lref{Proposition~2.8}: idempotents and vertex weights}]\genusstatus{AI-proposed, unchecked}
\label{prop:set-idem}
The idempotents of the quantum product are $\eid_1=(H-\Lv_2)/(\Lv_1-\Lv_2)$ and $\eid_2=(H-\Lv_1)/(\Lv_2-\Lv_1)$, with $H\qprod\eid_i=\Lv_i\eid_i$ and $\eta(\eid_i,\eid_j)=0$ for $i\neq j$. Put $\Delta_i\coloneqq1/\eta(\eid_i,\eid_i)$; the \emph{vertex weight} of a vertex of genus $g_v$ and colour $i$ is $\Delta_i^{g_v-1}$. For $\{i,j\}=\{1,2\}$,
\begin{equation}
\label{eq:set-Delta}
\begin{split}
\Delta_i&=\frac{\alpha_1\alpha_2\wA\wB\epsilon_1\epsilon_2(\Lv_i-\Lv_j)}{\alpha_1\alpha_2-E\Lv_j}
=\epsilon_1\epsilon_2(\Lv_i-\Lv_j)(\alpha_1\alpha_2-E\Lv_i)\\
&=\epsilon_1\epsilon_2(\alpha_0\alpha_1\alpha_2+2\alpha_1\alpha_2\Lv_i-E\Lv_i^2)
=\alpha_1\wA\epsilon_1\epsilon_2\,\cW_2(U_i),
\end{split}
\end{equation}
these four expressions being equal as elements of $\QA$; moreover $(\alpha_1\alpha_2-E\Lv_1)(\alpha_1\alpha_2-E\Lv_2)=\alpha_1\alpha_2\wA\wB$, $\Delta_1\Delta_2=-\epsilon_1^2\epsilon_2^2\,\alpha_1\alpha_2\wA\wB(\Lv_1-\Lv_2)^2$ and $1/\Delta_1+1/\Delta_2=E/(\alpha_1\alpha_2\wA\wB\epsilon_1\epsilon_2)$.
\end{prop}

Every factor of idempotent $2$ defined below is obtained from the corresponding factor of idempotent $1$ by applying the involution, that is, by evaluating the same expression in $U$ at the other root \Lref{Lemma~4.23}.

%-------------------------------------------------------------------------------
\subsection{Algebraic functions of \texorpdfstring{$\tau$}{tau}; the discriminant}
\label{sec:set-alg}

\begin{defn}[{\Lref{Definition~2.9}: algebraic functions of $\tau$}]
\label{def:set-alg}
Put $\delta^2\coloneqq(E+\alpha_0\tau)^2-4\wA\wB\tau=E^2-(E^2+\uq)\tau+\alpha_0^2\tau^2$, and let $\delta$ be the root analytic near $\tau=0$ with $\delta(0)=E$; it is analytic on the disc $\Ddisc\coloneqq\{|\tau|<r_\Ddisc\}$, $r_\Ddisc$ the smallest modulus of a root of $\delta^2$ (finite under (ND)). Further $\Ddisc'\coloneqq\Ddisc\cap\{|\tau|<1\}$ and $\Ddisc''\coloneqq\Ddisc'\setminus\{0\}$. The \emph{critical polynomial} is $\Ns(s)\coloneqq\wB\tau s^2-(E+\alpha_0\tau)s+\wA$, with roots
\[
\sast\coloneqq\frac{E+\alpha_0\tau-\delta}{2\wB\tau}=\frac{2\wA}{E+\alpha_0\tau+\delta},\qquad
\stil\coloneqq\frac{E+\alpha_0\tau+\delta}{2\wB\tau}=\frac{\wA}{\wB\tau\sast};
\]
$\sast$ is analytic on $\Ddisc$ with $\sast(0)=\wA/E$, and $\tau\stil$ is analytic with value $E/\wB$ at $0$. The \emph{critical points} are $U_1\coloneqq1-\tau\sast$ and $U_2\coloneqq1-\tau\stil$, the two roots of $\Pq$, with $U_1\to1$ and $U_2\to-\alpha_2/\wB$ as $t\to0$; idempotent $1$ is the \emph{singular} idempotent and idempotent $2$ the \emph{regular} one. The \emph{Euler phase} is $\phi_\tau(s)\coloneqq\wA\log s+\alpha_1\log(1-s)-\alpha_2\log(1-\tau s)$, with principal branches near $(\tau,s)=(0,\wA/E)$, and $\dphi(\tau)\coloneqq\phi_\tau(\sast(\tau))-\phi_0(\wA/E)$, an analytic germ at $\tau=0$ with $\dphi(0)=0$.
\end{defn}

\begin{lem}[{\Lref{Lemma~2.10}: the critical polynomial}]\genusstatus{AI-proposed, unchecked}
\label{lem:set-N}
Let $\tau\neq0$. Then $\phi_\tau'(s)=\Ns(s)/\big(s(1-s)(1-\tau s)\big)$, so the critical points of $\phi_\tau$ in $\bbC\setminus\{0,1,1/\tau\}$ are $\sast$ and $\stil$; on $\Ddisc'$, $\sast\notin\{0,1,1/\tau\}$, and on $\Ddisc''$, $\stil\notin\{0,1,1/\tau\}$ and $U_2\notin\{0,1,1/\cQ\}$. At $s=\sast$, with $\Lv_2\coloneqq\wB\sast-\alpha_0$:
\begin{enumerate}
\item $\Ns'(\sast)=2\wB\tau\sast-(E+\alpha_0\tau)=-\delta$ (and $\Ns'(\stil)=+\delta$);
\item $(1-\tau)\alpha_2s=(\wA-\wB s)(1-\tau s)$, equivalently $(E-\wB\tau s)(1-s)=\alpha_1(1-\tau s)$ (also at $s=\stil$); moreover $E-\wB\tau\sast=\delta+\tau\Lv_2$;
\item $\dd\sast/\dd\tau=\sast\Lv_2/\delta$;
\item $\wA-Es=-\tau s(\wB s-\alpha_0)$ (also at $s=\stil$).
\end{enumerate}
Consequently (envelope formula)
\begin{equation}
\label{eq:set-envelope}
\dphi_\tau=\frac{\alpha_2\sast}{1-\tau\sast}=\frac{\wA-\wB\sast}{1-\tau},\qquad \dphi'=\wA-\wB\sast ,
\end{equation}
and $\dphi=(\alpha_2\wA/E)\tau+O(\tau^2)$.
\end{lem}

\begin{lem}[{\Lref{Lemmas~2.11 and~2.12}: eigenvalues and $\Delta_2$ in terms of $\sast$}]\genusstatus{AI-proposed, unchecked}
\label{lem:set-eig}
With $U=1-\tau s$ ($\tau\neq0$), $\Pq(1-\tau s)=\cQ\tau\Ns(s)$ as polynomials in $s$. The eigenvalues are
\[
\Lv_1=\frac{\wA}{\tau\sast}-\alpha_0,\qquad \Lv_2=\wB\sast-\alpha_0,\qquad
\Lv_1+\Lv_2=\frac E\tau-\alpha_0,\qquad \Lv_1\Lv_2=\frac{\alpha_1\alpha_2}{\tau},\qquad \Lv_1-\Lv_2=\frac\delta\tau .
\]
Moreover $\Lv_1=E/t+\lrem(t)$ with $\lrem$ analytic at $t=0$, $\lrem(0)=E/2-\alpha_0-\alpha_1\alpha_2/E$, and $\Lv_2$ is analytic with $\Lv_2(0)=\alpha_1\alpha_2/E$. Finally $\alpha_1\alpha_2-E\Lv_2=-\wB\tau\sast\Lv_2$,
\[
\Delta_2=\epsilon_1\epsilon_2\wB\,\delta\,\sast\Lv_2,\qquad
\Delta_1\Delta_2=-\epsilon_1^2\epsilon_2^2\,\alpha_1\alpha_2\wA\wB\,\frac{\delta^2}{\tau^2},\qquad
\frac{\Delta_2}{\Lv_2-\Lv_1}=\epsilon_1\epsilon_2(\alpha_1\alpha_2-E\Lv_2),
\]
and $\Delta_2$ is analytic at $t=0$ with $\Delta_2(0)=\epsilon_1\epsilon_2\alpha_1\alpha_2\wA\wB/E\neq0$ and $\Delta_2'(0)=0$.
\end{lem}

\begin{lem}[{\Lref{Lemma~2.13}: the discriminant}]\genusstatus{AI-proposed, unchecked}
\label{lem:set-disc}
Let $\Disc(\cQ)\coloneqq\alpha_0^2+\big(2\alpha_0(\alpha_1+\alpha_2)+4\alpha_1\alpha_2\big)\cQ+(\alpha_1-\alpha_2)^2\cQ^2$. Then $\Disc=\cQ^2\tau^2(\Lv_1-\Lv_2)^2=\cQ^2\delta^2$ is the discriminant of $\Pq(U)$ in $U$; its own discriminant in $\cQ$ is $16\alpha_1\alpha_2\wA\wB$; $\Disc(1)=E^2$, $\Disc(0)=\alpha_0^2$, so at the generic point of $\TCY$, $\Disc$ has two distinct finite roots, none over $\cQ\in\{0,1\}$; $\Disc$ is irreducible over $\KCY=\Kred(\alpha_0)$. The double cover $\Ccov\to\bbP^1_\cQ$, $\ydc^2=\Disc(\cQ)$, has function field $\KCY(\Ccov)=\QA$, by $2\wB\cQ U-(\alpha_0+\cQ(\alpha_1-\alpha_2))=\pm\ydc$; it is ramified exactly over the two roots of $\Disc$ and unramified over $\cQ=0,1,\infty$.
\end{lem}

\begin{lem}[{\Lref{Lemma~2.14}: behaviour at $\cQ=1$, $\cQ=0$, $\cQ=\infty$}]\genusstatus{AI-proposed, unchecked}
\label{lem:set-limits}
The following hold.
\begin{enumerate}
\item At $\cQ=1$: $\Lv_1=E/\tau+O(1)$, $\Lv_2\to\alpha_1\alpha_2/E$, $t^2\Delta_1$ is analytic with $t^2\Delta_1|_{t=0}=-\epsilon_1\epsilon_2E^3$, and $\Delta_2(0)\neq0$. Idempotent $1$ is therefore the singular one.
\item At $\cQ=0$, label the roots of $\Pq$ by $U_+\to-\alpha_2/\alpha_0$ and $\cQ U_-\to\alpha_0/\wB$. Then $\Lv_+\to0$, $\Lv_-\to-\alpha_0$, and $\Delta_\pm(0)=\eul(\wts_\pm)$, both non-zero in $\KCY$ because $\alpha_0\neq0$: the point $\cQ=0$ is semisimple at $\alpha_0\neq0$ (on $\Tthree$ and on $\TF$ it is not).
\item At $\cQ=\infty$ ($\sQ\coloneqq1/\cQ\to0$): $U_1(\sQ)$ and $V_2(\sQ)\coloneqq\cQ U_2$ are power series in $\sQ$ over $\KCY$ with $U_1\to(\alpha_1-\alpha_2)/\wB$, $V_2\to-\alpha_2/(\alpha_1-\alpha_2)$; $\Lv_1\to\alpha_1$, $\Lv_2\to\alpha_2$, and $\Delta_1(\infty)=-\alpha_1\wB(\alpha_1-\alpha_2)\epsilon_1\epsilon_2$, $\Delta_2(\infty)=\alpha_2\wA(\alpha_1-\alpha_2)\epsilon_1\epsilon_2$, both non-zero in $\KCY$. Here the labels $1,2$ are those of the two points of $\Ccov$ over $\cQ=\infty$ at which $U\to(\alpha_1-\alpha_2)/\wB$, respectively $U\to0$.
\end{enumerate}
\end{lem}

% s2_setup_b (abridged version). Label prefix: set. Rows, factors, graph sum, rings, leaves, Hodge side.
%-------------------------------------------------------------------------------
\subsection{The rows and the factors}
\label{sec:set-rows}

For a finite multiset $\mathsf{v}$ of non-zero weights put
\begin{equation}
\label{eq:set-h}
h_{\mathsf{v}}(z)\coloneqq\exp\Big[-\sum_{k\geq1}\bk_k\,z^{2k-1}\sum_{a}\mathsf{v}_a^{1-2k}\Big],\qquad \bk_k\coloneqq\frac{B_{2k}}{2k(2k-1)},
\end{equation}
with $B_{2k}$ the Bernoulli numbers. Then $h_{\mathsf{u}+\mathsf{v}}=h_{\mathsf{u}}h_{\mathsf{v}}$, $h_{-\mathsf{v}}=h_{\mathsf{v}}^{-1}$, $h_{(v,-v)}=1$ and $h_{\mathsf{v}}(z)h_{\mathsf{v}}(-z)=1$; $h_{\mathsf{v}}$ is independent of $t$. The series in use are $\hw=h_{(E,-\epsilon_1,-\epsilon_2)}$, $\hnw=h_{(\alpha_1,\wA,-\epsilon_1,-\epsilon_2)}$, $\hreg=h_{(\alpha_1,\alpha_2,\wA,\wB,-E,-E)}$ and $h_{\wts_\pm}$, with $\bw\coloneqq(E,-\epsilon_1,-\epsilon_2)$ and $\nw\coloneqq(\alpha_1,\wA,-\epsilon_1,-\epsilon_2)$.

\begin{defn}[{\Lref{Definition~2.15}: rows}]
\label{def:set-rows}
The \emph{point-frame identity row} at idempotent $i=1,2$ is
\begin{equation}
\label{eq:set-rpt}
\rpt_{0,i}(z;t)\coloneqq\sum_{m\geq0}z^m\,\wick_m\big(\cW_2(U_i),\cW_3(U_i),\ldots\big),
\end{equation}
where the universal polynomial $\wick_m$ is the coefficient of $z^m$ in
\[
\sum_{n\geq0}\frac{(2n)!}{n!\,2^n}\Big(\frac{-z}{\cW_2}\Big)^{n}\,\coef{v^{2n}}\exp\Big[\sum_{k\geq3}\frac{\cW_kv^k}{k!\,z}\Big];
\]
thus $\wick_0=1$, and $\wick_m$ is a polynomial in $\cW_2^{-1},\cW_3,\ldots,\cW_{2m+2}$ with rational coefficients. The row \eqref{eq:set-rpt} is the normalised formal Laplace expansion of $\re^{\cW/z}\dd U/U$ at the critical point $U_i$ (\cref{prop:row-wickrow} for $i=1$, \cref{thm:lea-row2reg} for $i=2$). The \emph{row with the $\cQ=0$ constant} is $\rZ_{0,i}\coloneqq\hnw(z)\,\rpt_{0,i}(z;t)$; the matrix with these rows is $\RZ$. The \emph{frame-B rows} are
\begin{equation}
\label{eq:set-rowB}
\rr_{0,i}(z;t)\coloneqq h_{(\alpha_1,\wA,-E)}(z)\,\rpt_{0,i}(z;t)=\hw(z)^{-1}\hnw(z)\,\rpt_{0,i}(z;t),
\end{equation}
the second equality by \cref{lem:set-hmult}; the matrix with these rows is $\RB=\RZ\hw^{-1}$, and \emph{frame B} means this frame. The frame-B rows depend on $(\alpha_0,\alpha_1,\alpha_2)$ only. The \emph{$H$-row} is
\begin{equation}
\label{eq:set-Hrow}
\rr_{1,i}(z)\coloneqq(\Lv_i+z\Dop_i)\,\rr_{0,i}(z),\qquad \Dop_i\coloneqq\frac{\dd}{\dd t}-\frac{\Delta_i'}{2\Delta_i},
\end{equation}
a rule which holds in every frame, the scalar series $h$ being independent of $t$. The \emph{covector} of idempotent $i$ is $\vv_i(z)\coloneqq(\rr_{0,i}(z),\rr_{1,i}(z))$, and $\rr[m]\coloneqq\coef{z^m}\rr$.
\end{defn}

\begin{lem}[{\Lref{Lemma~2.16}: the multiset identity}]\genusstatus{AI-proposed, unchecked}
\label{lem:set-hmult}
One has
\begin{equation}
\label{eq:set-hmult}
\hnw(z)\,h_{(E,-\alpha_1,-\wA)}(z)=\hw(z),\qquad\text{equivalently}\qquad \hw^{-1}\hnw=h_{(\alpha_1,\wA,-E)} .
\end{equation}
Moreover, for two multisets $\bw',\bw$ of non-zero weights, $h_{\bw'}=h_{\bw}$ if and only if $\bw'$ and $\bw$ agree up to pairs $(v,-v)$. In particular $\bw=(E,-\epsilon_1,-\epsilon_2)$ is, up to such pairs, the unique class with $h_{\bw}=\hnw h_{(E,-\alpha_1,-\wA)}$.
\end{lem}

The class $\bw$ makes the frame-B constant at $t=0$ the geometric series $\rr_{0,1}(z;0)=E/(E+z)$ (\cref{cor:row-F2}), because the point-frame row at $t=0$ is $h_{(E,-\alpha_1,-\wA)}E/(E+z)$ (\cref{prop:row-beta}); the same scalar $\hnw$ multiplies both idempotents (\cref{lem:ide-norm}).

\begin{defn}[{\Lref{Definition~2.17}: dilaton and edge factors, legs}]
\label{def:set-factors}
All rows are the frame-B rows. The \emph{dilaton factors} are $\Dil_i(d)\coloneqq(-1)^d\,\rr_{0,i}[d-1]$ for $d\geq1$, and the \emph{dilaton series} at idempotent $i$ is
\[
\qt^{(i)}(\psi)\coloneqq\sum_{d\geq1}\Dil_i(d;t)\,\psi^d=-\psi\,\rr_{0,i}(-\psi;t),\qquad \qt\coloneqq\qt^{(1)} ;
\]
since $\rr_{0,i}[0]=1$, $\Dil_i(1)=-1$ at every $t$: the term $d=1$ is the dilaton shift. The \emph{pairing of the rows} and the \emph{edge factors} are
\begin{equation}
\label{eq:set-Pij}
\Pij_{ij}(z,w)\coloneqq-\langle\vv_i(-z),\vv_j(-w)\rangle ,
\end{equation}
\begin{equation}
\label{eq:set-edge}
\Edge_{ij}(z,w)\coloneqq\frac{\Pij_{ij}(z,w)-\Pij_{ij}(-w,w)}{z+w}=\sum_{m,n\geq0}\Edge_{ij}(m,n)\,z^mw^n ,
\end{equation}
with the pairing \eqref{eq:set-pairing}; the numerator vanishes on $z=-w$, so the quotient is a power series in $z,w$ coefficientwise in $t$. In leaf language $\Edge_{11}(\psi,\psi')$ means $\Edge_{11}(z=\psi,w=\psi')$. The \emph{legs} are $f_{j,b}(\psi)\coloneqq\coef{w^b}\Edge_{1j}(\psi,w)$, for $j\in\{1,2\}$ and $b\geq0$.
\end{defn}

Once the symplectic relations $\langle\vv_i(z),\vv_j(-z)\rangle=\delta_{ij}\Delta_i$ are known (\cref{thm:lea-qde}), \eqref{eq:set-edge} reads $\Edge_{ij}(z,w)=[\delta_{ij}\Delta_i-\langle\vv_i(-z),\vv_j(-w)\rangle]/(z+w)$, and $\Edge_{ij}(m,n)=\Edge_{ji}(n,m)$.

\begin{defn}[{\Lref{Definition~2.20}: $\rho$, the leaves, the leaf spaces}]
\label{def:set-rho}
By \cref{lem:set-limits}(i), $t^2\Delta_1/(-\epsilon_1\epsilon_2E^3)$ is analytic at $t=0$ with value $1$; let $\rho(t)\coloneqq\big(\Delta_1t^2/(-\epsilon_1\epsilon_2E^3)\big)^{-1/2}$, the branch with $\rho(0)=1$. The \emph{leaves} are $\lf_j(\psi)\coloneqq1/(E-j\psi)$ for $j\geq1$; in the row variable $z=-\psi$, $\lf_j=1/(E+jz)$. The \emph{leaf spaces} are $\Lsp_m\coloneqq\mathrm{span}(\lf_1,\ldots,\lf_m)$ and $\Lsp_m^+\coloneqq(\text{constants})+\Lsp_m$. Further
\[
\qo\coloneqq E-E^2\lf_1=-\frac{E\psi}{E-\psi},\qquad \Tdil(\psi)\coloneqq\qo+\psi=-\frac{\psi^2}{E-\psi}.
\]
\end{defn}

%-------------------------------------------------------------------------------
\subsection{Coloured stable graphs and the graph sum}
\label{sec:set-graphsum}

For a stable curve of genus $h$, with $\lambda_i=c_i(\Hdg)$ the Chern classes of the Hodge bundle, let $\Lamv(u)\coloneqq\sum_i(-1)^i\lambda_iu^{h-i}$; for a multiset $\mathsf{v}$ of non-zero weights the \emph{normalised Hodge class} is $\Hg_{\mathsf{v}}\coloneqq\prod_{a\in\mathsf{v}}\Lamv(\mathsf{v}_a)/\mathsf{v}_a^{h}$. The vertex class is
\begin{equation}
\label{eq:set-Hg}
\Hg_{\bw}=\prod_{a=1}^{3}\frac{\Lamv(\bw_a)}{\bw_a^{h}},\qquad \bw=(E,-\epsilon_1,-\epsilon_2),
\end{equation}
the same at both colours and independent of $t$. It is a pull-back under the forgetful maps, so the string and dilaton equations hold for it when $h\geq2$.

\begin{defn}[{\Lref{Definition~2.22}: coloured stable graphs; the decoration sum; the graph sum}]
\label{def:set-graphsum}
Fix $g\geq2$. Let $\Gamma$ run over the connected stable graphs of genus $g$ without legs, with vertex set $\Ver(\Gamma)$, edge set $\Edg(\Gamma)$, vertex genera $g_v$ and $n_v$ half-edges at $v$. A \emph{colouring} is a map $\col:\Ver(\Gamma)\to\{1,2\}$; colour $1$ is singular, colour $2$ regular. A \emph{decoration} assigns an exponent $a_h\geq0$ to each half-edge $h$ ($a_e,b_e$ at the two half-edges of an edge $e$, in a chosen order of its ends $v(e),v'(e)$) and to each vertex $v$ an integer $m_v\geq0$ with an ordered tuple $(d_{v,1},\ldots,d_{v,m_v})$ of \emph{dilaton degrees} $d_{v,k}\geq2$, with weight $1/m_v!$. The \emph{decoration sum} of $(\Gamma,\col)$ is
\begin{multline}
\label{eq:set-graphsum}
\GCont(\Gamma,\col)\coloneqq\sum_{\mathrm{dec}}\ \prod_{v\in\Ver(\Gamma)}\Big[\Delta_{\col_v}^{\,g_v-1}\,\frac{1}{m_v!}\int_{\Mbar_{g_v,n_v+m_v}}\Hg_{\bw}\prod_{h\text{ at }v}\psi_h^{a_h}\prod_{k=1}^{m_v}\psi_{n_v+k}^{d_{v,k}}\ \prod_{k=1}^{m_v}\Dil_{\col_v}(d_{v,k})\Big]\\
\times\prod_{e\in\Edg(\Gamma)}\Edge_{\col_{v(e)}\col_{v'(e)}}(a_e,b_e),
\end{multline}
and the \emph{frame-B Hodge-vertex graph sum} is
\begin{equation}
\label{eq:set-Gsum}
\Gsum_g\coloneqq\sum_{\Gamma}\frac{1}{|\Aut\Gamma|}\sum_{\col:\Ver(\Gamma)\to\{1,2\}}\GCont(\Gamma,\col).
\end{equation}
\end{defn}

The decoration sum is finite, since the integral vanishes unless $\sum_ha_h+\sum_kd_{v,k}\leq3g_v-3+n_v+m_v$ at every vertex, which forces $m_v\leq3g_v-3+n_v$; ``regular at the conifold point after summing over decorations'' means that the Laurent expansion of $\GCont(\Gamma,\col)$ in $t$ has no negative powers. Equation \eqref{eq:set-graphsum} is a definition; the regularity statements of \cref{sec:reg} are statements about it, and its identification with $\sF_g$ is Theorem~B, used only in \cref{sec:rig}.

%-------------------------------------------------------------------------------
\subsection{Rings, principles and leaf identities}
\label{sec:set-rings}

\begin{lem}[{\Lref{Lemma~2.23}: rings and principles; condensed}]\genusstatus{AI-proposed, unchecked}
\label{lem:set-rings}
Let $\Oan\coloneqq\bbC\{t\}=\bbC\{\tau\}$, $\Man$ the field of meromorphic germs, $\Oan(\Ddisc)$ the functions analytic on $\Ddisc$, $\Cz=\bbC\lser z$, $\Rz\coloneqq\Cz\lser\tau$ with the derivation $\dd/\dd t=(1-\tau)\dd/\dd\tau$, and $\Rzo\coloneqq\bbC\fps z\lser\tau\subset\Rz$; $\iota:\Oan\fps z\to\bbC\fps{z,\tau}$ Taylor-expands every $z$-coefficient, and $\Man\lser z$ denotes Laurent series in $z$ with coefficients in $\Man$.
\begin{description}
\item[Transfer.] An identity in $\Rz$ between elements of $\Rzo$ is an identity of formal double series: reading off $z^m$ gives identities of Laurent series in $\tau$ and, when these are expansions of meromorphic germs (as for all rows), identities of meromorphic functions of $t$ near $0$.
\item[(P)] Every object built from the genus-zero data (the rows, the dilaton and edge factors, $\rho$, $\Lv_i$, $\Delta_i$) is obtained from the weights and $\tau$ by field operations, $\dd/\dd t$, and the expansions of the square roots $\delta$, $\alg$ and of logarithms of units; at $\tau=0$ the only denominators are $\alpha_1,\alpha_2,\wA,\wB,E,\epsilon_1,\epsilon_2$ and the Gauss denominators $E+jz$. Hence the $t$-, $z$- and $\psi$-coefficients are rational functions of the weights, regular on (ND); an identity between such coefficients proved at every complex weight point of a non-empty open subset of the weight space holds in $\Kal$ (or $\KCY$), and an identity or a membership proved over $\Kal$ specialises to every weight point with (ND).
\item[Leaf spaces.] A rational function $f$ of $\psi$ lies in $\Lsp_m^+$ if and only if its poles are simple, lie in $\{E/j:1\leq j\leq m\}$, and $f$ is bounded at $\infty$; it lies in $\Lsp_m$ if and only if moreover $f(\psi)\to0$ as $\psi\to\infty$. A rational function regular at $0$ is identified with its Taylor series at $0$.
\item[(T)] For an increasing chain $V_0\subset V_1\subset\cdots$ of subspaces of a vector space $V$ and $f\in V\fps\tau=V\fps t$: ($\coef{\tau^n}f\in V_n$ for $n\leq k$) if and only if ($\coef{t^n}f\in V_n$ for $n\leq k$); multiplication by a series in $t$ not involving the auxiliary variables preserves such a filtration.
\end{description}
\end{lem}

\begin{lem}[{\Lref{Lemma~2.24}: leaf identities, partial fractions}]\genusstatus{AI-proposed, unchecked}
\label{lem:set-leafid}
For $j,m\geq1$:
\textup{(LAM1)} $\coef{\psi^n}\lf_j=j^n/E^{n+1}$;
\textup{(LAM2)} $\psi\lf_j=(E\lf_j-1)/j$;
\textup{(LAM3)} $(E-m\psi)\lf_j=m/j+E(1-m/j)\lf_j$;
\textup{(LAM4)} $f=\sum_{j}f_j\lf_j$ is $O(\psi^{-2})$ at $\psi=\infty$ if and only if $\sum_jf_j/j=0$;
\textup{(LAM5)} $\lf_i\lf_j=(i\lf_i-j\lf_j)/(E(i-j))$ for $i\neq j$;
\textup{(LAM6)} $\lf_j(\psi)-\lf_j(-w)=j(\psi+w)\lf_j(\psi)\lf_j(-w)$;
\textup{(PF)} a rational function $f$ of $z$ whose poles are of order $\leq1$, lie in $\{-E/j:1\leq j\leq M\}$, and which is bounded at $\infty$, equals $f(\infty)+\sum_{j=1}^{M}\beta_j/(E+jz)$ with $\beta_j=\lim_{z\to-E/j}(E+jz)f(z)$.
\end{lem}

%-------------------------------------------------------------------------------
\subsection{The Hodge side}
\label{sec:set-hodge}

Fix $h\geq2$, and let $\RRing\coloneqq\RRing_0\fps t$ with $\RRing_0$ a $\bbQ$-algebra in which $\epsilon_1,\epsilon_2,E$ are invertible (for instance $\bbC$ or $\KCY$).

\begin{defn}[{\Lref{Definition~2.25}: the Hodge potential and the numbers $\Xiz_h$}]
\label{def:set-fpot}
For $\mathsf{b}_1,\ldots,\mathsf{b}_s\in\RRing\fps\psi$ put
\[
\dbase^s(\mathsf{b}_1,\ldots,\mathsf{b}_s)\coloneqq\sum_{m\geq0}\frac{1}{m!}\int_{\Mbar_{h,s+m}}\Hg_{\bw}\prod_{i=1}^{s}\mathsf{b}_i(\psi_i)\prod_{k=1}^{m}\Tdil(\psi_{s+k}),
\]
a finite sum. Let $\Dom\coloneqq\qo+t\RRing\fps\psi$. For $\mathbf{q}=\qo+\upsilon\in\Dom$ and $f_1,\ldots,f_n\in\RRing\fps\psi$ define
\[
\dd^n\Fpot_h(\mathbf{q})[f_1,\ldots,f_n]\coloneqq\sum_{r\geq0}\frac{1}{r!}\,\dbase^{n+r}(f_1,\ldots,f_n,\upsilon,\ldots,\upsilon)\quad(r\text{ entries }\upsilon),\qquad \Fpot_h(\mathbf{q})\coloneqq\dd^0\Fpot_h(\mathbf{q}),
\]
which converges $t$-adically. If $\mathbf{q}=-\psi+O(\psi^2)$, the \emph{dilaton-shifted form} is
\[
\Fpot_h(\mathbf{q})\coloneqq\sum_{m\geq0}\frac{1}{m!}\int_{\Mbar_{h,m}}\Hg_{\bw}\prod_{i=1}^{m}(\mathbf{q}+\psi)(\psi_i),
\]
and similarly for $\dd^n\Fpot_h(\mathbf{q})[f_1,\ldots,f_n]$; on the overlap with $\Dom$ the two definitions agree. In this form the decorations of \eqref{eq:set-graphsum} appear: since $\Dil_i(1)\psi=-\psi$, the vertex factor at a vertex of genus $h\geq2$ and colour $i$ with half-edge series $f_1,\ldots,f_n$ is $\Delta_i^{h-1}\dd^n\Fpot_h(\qt^{(i)})[f_1,\ldots,f_n]$. At idempotent $1$, $\qt\in\Dom$ and $\rho\qt\in\Dom$ (by (F1)--(F3) of \cref{thm:int-leaf}); at idempotent $2$, $\qt^{(2)}\notin\Dom$, only the dilaton-shifted form is used, and nothing is rescaled. For a multiset $\msU=\{j_1,\ldots,j_r\}$ of positive integers,
\[
\Xiz_h(\msU)\coloneqq\dd^r\Fpot_h(\qo)[\lf_{j_1},\ldots,\lf_{j_r}]=\dbase^r(\lf_{j_1},\ldots,\lf_{j_r}),
\]
the leaf-scale-$E$ integrals with the class $\Hg_{\bw}$; $J_{\msU}\coloneqq\sum_ij_i$, $\dg(\msU)\coloneqq\sum_i(j_i-1)$, and $\ZZ[h,n]\coloneqq\Xiz_h(1^n)$, so that $\ZZ[h,0]=\Fpot_h(\qo)$.
\end{defn}

\begin{defn}[{\Lref{Definition~2.26}: the statement $\Mirr(h)$}]
\label{def:set-mirr}
$\Mirr(h)$ is the statement: $\Xiz_h(\msM)=0$ for every \emph{non-empty} multiset $\msM$ of integers $\geq2$ with $\dg(\msM)\leq2h-3$.
\end{defn}

The empty multiset is excluded because $\Xiz_h(\emptyset)=\ZZ[h,0]$ is the value of the pole (\cref{thm:reg-VI}); parts equal to $1$ are removed by the string identity (\cref{lem:reg-string}). The statement $\Mirr(h)$ for the class $\Hg_{\bw}$, at the generic point of $\TCY$ and at every weight point with $\epsilon_1\epsilon_2E\neq0$, is the cited \cref{thm:reg-L}.

% s3_rows. Abridged Section 3 (long version: sections/s03_rows.tex); label prefix: row.
\section{The singular row: Euler integral, Gauss equation and the Beta constant}
\label{sec:row}

This section computes the frame-B identity row $\rr_{0,1}(z;t)$ at the singular idempotent in closed form, at weights with (ND), complex or indeterminate (principle (P) of \cref{lem:set-rings}). All statements are formal, except \cref{lem:row-betaphase,prop:row-beta}, whose analytic inputs are Laplace's method, proved in full in \Lref{Lemma~3.11}, the Beta integral \cite[5.12.1]{DLMF}, and Stirling's series with remainder \cite[5.11.1 and Section~5.11(ii)]{DLMF}.

%-------------------------------------------------------------------------------
\subsection{The Euler coordinate and the formal Laplace calculus}
\label{sec:row-euler}

For $\tau\neq0$ put $U=1-\tau s$. Then $1-U=\tau s$, $1-\cQ U=-\cQ\tau(1-s)$, $\dd U/U=-\tau\,\dd s/(1-\tau s)$, and $\cW(1-\tau s)$ is the Euler phase $\phi_\tau(s)$ of \cref{def:set-alg} plus a term independent of $s$ \Lref{Lemma~3.1}; so $\re^{\cW/z}\dd U/U$ is $\re^{\phi_\tau(s)/z}(1-\tau s)^{-1}\dd s$ times a factor independent of $s$. The logarithmic points $U=1$ and $U=1/\cQ$ of $\cW$, which collide at $t=0$, are $s=0$ and $s=1$, and the critical point $\sast$ of $\phi_\tau$ stays at $\sast(0)=\wA/E$.

\begin{lem}[{\Lref{Lemma~3.2}: Hessians, the unit $\alg$, the tensor rule; formal}]\genusstatus{AI-proposed, unchecked}
\label{lem:row-hess}
Let $\tau\neq0$. Then
\begin{equation}
\label{eq:row-hess}
\phi_\tau''(\sast)=-\frac{\delta}{\sast(1-\sast)U_1},\qquad
\phi_\tau''(\stil)=\frac{\delta}{\stil(1-\stil)U_2},\qquad
\phi_0''(\wA/E)=-\frac{E^3}{\alpha_1\wA},
\end{equation}
and both critical points are non-degenerate on $\Ddisc''$. Put $\alg(\tau)\coloneqq(1-\tau\sast)\big(\phi_\tau''(\sast)/\phi_0''(\wA/E)\big)^{1/2}$, the square root with $\alg(0)=1$. It is well defined:
\begin{equation}
\label{eq:row-alg}
\alg^2=\frac{U_1^2\,\phi_\tau''(\sast)}{\phi_0''(\wA/E)}=\frac{\alpha_1\wA\,\delta\,U_1}{E^3\,\sast(1-\sast)}=\frac{\wA\,\delta\,(E-\wB\tau\sast)}{E^3\,\sast}
\end{equation}
is a unit of $\Oan(\Ddisc)$ with value $1$ at $\tau=0$, so $\alg$ and $1/\alg$ are analytic on $\Ddisc$. Tensor rule: $\cW_2(U_i)=U_i^2\phi_\tau''(s_i)/\tau^2$ with $(s_1,s_2)\coloneqq(\sast,\stil)$.
\end{lem}

So $\cW_2(U_1)$ has a double pole at $\tau=0$. For $k\geq2$ the rescaled derivatives $\wrk_k(\tau)\coloneqq\tau^k\cW_k(U_1)$ extend analytically to $\Ddisc'$, $\tau=0$ included, with $\wrk_k(0)=(-1)^k\phi_0^{(k)}(\wA/E)$, and $\wrk_2$ has no zero on $\Ddisc'$ \Lref{Lemma~3.4}.

\begin{defn}[{\Lref{Definitions~3.5 and~3.7}: normalised formal Laplace expansions}]
\label{def:row-Lnorm}
Let $\Acal$ be the ring of formal sums $K=\sum_{p\geq0,\,q\in\bbZ}K_{pq}v^pz^q$ on whose support the weight $p/2+q$ is bounded below and takes each value $\leq N$ only finitely often. For $a\neq0$ the \emph{Gaussian functional} is
\[
\gauss{K}{a}\coloneqq\sum_{p\ \mathrm{even},\,q}K_{pq}\,(p-1)!!\,\Big(-\frac za\Big)^{p/2}z^q\ \in\Cz\qquad((-1)!!\coloneqq1).
\]
Data: a point $u_0$, a \emph{phase} $f$ analytic near $u_0$ with $f'(u_0)=0$ and $f_2\coloneqq f''(u_0)\neq0$, and an \emph{amplitude} $\amp\in\Oan_{u_0}[z,1/z]$, where $\Oan_{u_0}$ denotes the germs of analytic functions at $u_0$. With $\Rr_f(v)\coloneqq\sum_{k\geq3}f^{(k)}(u_0)v^k/k!$ and $\mathrm{Exp}[\Rr/z]\coloneqq\sum_j\Rr^j/(j!z^j)\in\Acal$, put $\Lnorm[f,\amp;u_0]\coloneqq\gauss{\amp(u_0+v)\,\mathrm{Exp}[\Rr_f(v)/z]}{f_2}\in\Cz$, and $\Lhat[f,\amp;u_0]\coloneqq\Lnorm[f,\amp;u_0]/\amp(u_0)$ for $\amp\in\Oan_{u_0}$ with $\amp(u_0)\neq0$. For $\amp=1$ the coefficient of $z^m$ in $\Lnorm[f,1;u_0]$ is the Wick polynomial $\wick_m$ of \cref{def:set-rows} with $\cW_k$ replaced by $f^{(k)}(u_0)$. For $\amp\in\Oan_{u_0}$, $\Lnorm[f,\amp;u_0]$ lies in $\bbC\fps z$, with constant term $\amp(u_0)$ and coefficients polynomial in $1/f_2$ and the Taylor coefficients of $f$ and $\amp$ at $u_0$. $\Lnorm$ is unchanged when a constant is added to $f$ and is linear in $\amp$. Rescaling of the phase: for a constant $\lambda\neq0$ and an amplitude $\amp$ independent of $z$, $\Lnorm[\lambda f,\amp;u_0](z)=\Lnorm[f,\amp;u_0](z/\lambda)$.
\end{defn}

\begin{prop}[{\Lref{Propositions~3.8, 3.9 and~3.10}: the three rules (P1), (P2), (P3); formal}]\genusstatus{AI-proposed, unchecked}
\label{prop:row-calc}
Let $f$, $u_0$ and $\amp$ be as in \cref{def:row-Lnorm}.
\begin{description}
\item[(P3) Total derivatives.] $\Lnorm[f,\mathsf{k}'+f'\mathsf{k}/z;u_0]=0$ for every $\mathsf{k}\in\Oan_{u_0}[z,1/z]$.
\item[(P2) Derivative in a parameter with a moving expansion point.] Here $\tau$ is any parameter and $s$ any coordinate. Let $f(\tau,s)$ and $\amp(\tau,s)$ be analytic near $(\tau_0,s_0)$ ($\amp$ may be a Laurent polynomial in $z$ with such coefficients), let $s_c(\tau)$ be analytic with $s_c(\tau_0)=s_0$, $\de_sf(\tau,s_c(\tau))=0$ and $f_2(\tau)\coloneqq\de_s^2f(\tau,s_c(\tau))\neq0$. Put $f_0(\tau)\coloneqq f(\tau,s_c(\tau))$ and $\mathsf{n}(\tau)\coloneqq\Lnorm[f(\tau,\cdot),\amp(\tau,\cdot);s_c(\tau)]$. Then every $z$-coefficient of $\mathsf{n}$ is analytic in $\tau$, and
\begin{equation}
\label{eq:row-P2}
\de_\tau\mathsf{n}+\Big(\frac{f_{0,\tau}}{z}-\frac{f_{2,\tau}}{2f_2}\Big)\mathsf{n}=\Lnorm[f(\tau,\cdot),\Dam\amp;s_c(\tau)],\qquad \Dam\amp\coloneqq\de_\tau\amp+\amp\,\frac{\de_\tau f}{z}.
\end{equation}
\item[(P1) Coordinate invariance.] Let $h$ be analytic near $v_0$ with $h(v_0)=u_0$, $h_1\coloneqq h'(v_0)\neq0$, and let $\amp\in\Oan_{u_0}$. Then $f\circ h$ has a non-degenerate critical point at $v_0$, and $\Lnorm[f\circ h,(\amp\circ h)\,h';v_0]=h_1\,\Lnorm[f,\amp;u_0]$; so $\Lhat[f\circ h,(\amp\circ h)h';v_0]=\Lhat[f,\amp;u_0]$ if $\amp(u_0)\neq0$.
\end{description}
\end{prop}

\begin{proof}[Sketch of proof]
(P3): $\re^{f/z}(\mathsf{k}'+f'\mathsf{k}/z)=\dd/\dd u\,(\re^{f/z}\mathsf{k})$, and the statement is Gaussian integration by parts, $\gauss{\dd K/\dd v}{a}=-(a/z)\gauss{vK}{a}$. (P2): $f_{0,\tau}/z$ comes from the moving critical value, $-f_{2,\tau}/(2f_2)$ from the Gaussian normalisation, and the motion of the expansion point contributes a total-derivative amplitude, killed by (P3). (P1): interpolate linearly between the affine map $u_0+h_1(v-v_0)$ and $h$; by (P2) and (P3) the expansion does not depend on the interpolation parameter.
\end{proof}

%-------------------------------------------------------------------------------
\subsection{The point-frame row and Gauss's equation}
\label{sec:row-wickrow}

Put
\[
\Nlap(z,\tau)\coloneqq\Lhat[\phi_\tau,(1-\tau s)^{-1};\sast(\tau)]=U_1\,\Lnorm[\phi_\tau,(1-\tau s)^{-1};\sast(\tau)],\qquad \Abeta(z)\coloneqq\Lnorm[\phi_0,1;\wA/E].
\]

\begin{prop}[{\Lref{Proposition~3.13}; formal}]\genusstatus{AI-proposed, unchecked}
\label{prop:row-wickrow}
\begin{enumerate}
\item For every $\tau\in\Ddisc''$ (every small $t\neq0$), coefficientwise in $z$, $\rpt_{0,1}(z;t)=\Nlap(z,\tau)$.
\item For every $m\geq0$ and $\tau\in\Ddisc''$, $\coef{z^m}\rpt_{0,1}(z;t)=\wick_m\big(\wrk_2(\tau),\wrk_3(\tau),\dots,\wrk_{2m+2}(\tau)\big)$. Hence the $z^m$ coefficient of $\rpt_{0,1}$ extends analytically to all of $\Ddisc'$, $\tau=0$ included, with value $\coef{z^m}\Abeta$ at $\tau=0$.
\item $\Nlap$ lies in $\Oan(\Ddisc')\fps z$, $\Nlap(z,0)=\Abeta(z)$, and the identity of (i) holds on all of $\Ddisc'$ and in $\bbC\fps z\fps t$.
\end{enumerate}
\end{prop}

\begin{proof}[Sketch of proof]
In the coordinate $y=\log U$ the integrand $\re^{\cW/z}\dd U/U$ has amplitude $1$ and $\dd/\dd y=\thU$, so by \cref{def:set-rows} $\rpt_{0,1}=\Lnorm[\cW(\re^y),1;\log U_1]$. The change of coordinate $y=\log(1-\tau s)$ turns the phase into $\phi_\tau+\mathrm{const}$ and the amplitude into $-\tau(1-\tau s)^{-1}$; (P1) in its normalised form gives (i). For (ii): every monomial of $\wick_m$ is $\cW_2^{-n}\prod_i\cW_{k_i}$ with $\sum_ik_i=2n$, hence invariant under $\cW_k\mapsto\lambda^k\cW_k$ \Lref{Lemma~3.12}; with $\lambda=\tau$ it is the same monomial in the $\wrk_k$, so it is analytic at $\tau=0$, with no cancellation between monomials.
\end{proof}

Gauss's parameters, series and operator (acting on $\Cz\fps\tau$) are
\begin{equation}
\label{eq:row-gaussdef}
\begin{gathered}
\ga\coloneqq1+\frac{\alpha_2}{z},\quad \gb\coloneqq1+\frac{\wA}{z},\quad \gc\coloneqq2+\frac Ez,\qquad
\Fhyp\coloneqq{}_2F_1(\ga,\gb;\gc;\tau)=\sum_{n\geq0}\frac{\prod_{j=1}^{n}(\alpha_2+jz)(\wA+jz)}{n!\,z^n\prod_{j=2}^{n+1}(E+jz)}\,\tau^n,\\
\Gauss[Y]\coloneqq\tau(1-\tau)Y_{\tau\tau}+\big(\gc-(\ga+\gb+1)\tau\big)Y_\tau-\ga\gb\,Y ,
\end{gathered}
\end{equation}
and $\Ephi\coloneqq\mathrm{Exp}[-\dphi/z]\in\bbC[1/z]\fps\tau$.

\begin{prop}[{\Lref{Proposition~3.16}: Gauss's equation for the formal expansion; formal}]\genusstatus{AI-proposed, unchecked}
\label{prop:row-gaussode}
Let $\Lt\coloneqq\Lnorm[\phi_\tau,(1-\tau s)^{-1};\sast(\tau)]$, an element of $\Oan(\Ddisc')\fps z$ by \cref{prop:row-wickrow}, and
\[
\Ngauss(z,\tau)\coloneqq\mathrm{Exp}[\dphi(\tau)/z]\,\Big(\frac{\phi_\tau''(\sast)}{\phi_0''(\wA/E)}\Big)^{-1/2}\iota(\Lt)\ \in\Cz\fps\tau
\]
(root with value $1$ at $\tau=0$; $\iota$ as in \cref{lem:set-rings}). Then $\Gauss[\Ngauss]=0$ in $\Cz\fps\tau$, $\Ngauss(z,0)=\Abeta(z)$, which lies in $1+z\Kal\fps z$, and $\Ngauss=\mathrm{Exp}[\dphi/z]\,\Nlap/\alg$.
\end{prop}

\begin{proof}[Sketch of proof]
The decisive identity is pointwise in $(s,\tau)$ \Lref{Lemma~3.15}: with $f=\phi_\tau$, $\amp_0=(1-\tau s)^{-1}$, $\Dam$ as in \eqref{eq:row-P2} and $\kE\coloneqq s(1-s)(1-\tau s)^{-2}$,
\[
\tau(1-\tau)\Dam^2\amp_0+\big(\gc-(\ga+\gb+1)\tau\big)\Dam\amp_0-\ga\gb\,\amp_0=-\ga\,\big(\kE'+f'\kE/z\big),
\]
which is Euler's integration by parts, written for amplitudes. By (P2), $\nabla\coloneqq\dd/\dd\tau+\dphi_\tau/z-f_{2,\tau}/(2f_2)$ satisfies $\nabla\Lnorm[q]=\Lnorm[\Dam q]$ at the moving point $\sast(\tau)$; so Gauss's operator with $\nabla$ in place of $\dd/\dd\tau$, applied to $\Lt$, is the expansion of a total-derivative amplitude, which vanishes by (P3). Conjugation by $\mathrm{Exp}[\dphi/z](f_2(\tau)/f_2(0))^{-1/2}$ turns $\nabla$ into $\dd/\dd\tau$. No boundary term occurs.
\end{proof}

\begin{thm}[{\Lref{Theorem~3.18}: the closed form; formal}]\genusstatus{AI-proposed, unchecked}
\label{thm:row-closed}
In $\Cz\fps\tau$, $\Nlap=\Abeta\,\Fhyp\,\Ephi\,\alg$.
\end{thm}

\begin{proof}[Sketch of proof]
If $Y=\sum_kY_k\tau^k\in\Cz\fps\tau$ satisfies $\Gauss[Y]=0$, then $(k+1)(k+\gc)Y_{k+1}=(k+\ga)(k+\gb)Y_k$, and $k+\gc=(E+(k+2)z)/z$ is a unit of $\Cz$ because $E\neq0$; so $Y=Y_0\Fhyp$ \Lref{Lemma~3.17}. Apply this to $\Ngauss$, with $\Ngauss(z,0)=\Abeta$, and multiply by $\alg\,\Ephi$.
\end{proof}

%-------------------------------------------------------------------------------
\subsection{The Beta constant, the frame-B row, and (F2)}
\label{sec:row-beta}

The evaluation of $\Abeta$ is the one analytic input of this section and of \cref{sec:lea} (input (i) of \cref{sec:int-inputs}); $h_{\mathsf{v}}$ is the series \eqref{eq:set-h}. \Cref{lem:row-betaphase} is used again in \cref{prop:lea-regrows}, \cref{lem:ide-norm} and \cref{prop:rig-flop}.

\begin{lem}[{\Lref{Lemma~3.19}: the Beta phase; analytic input}]\genusstatus{AI-proposed, unchecked}
\label{lem:row-betaphase}
Let $a$, $b$ be indeterminates, or complex numbers with $ab(a+b)\neq0$, and let $f(y)\coloneqq ay+b\log(1-\re^y)$, the phase $a\log U+b\log(1-U)$ in the coordinate $y=\log U$. Then $f$ has a critical point at $y_0\coloneqq\log(a/(a+b))$, with $f''(y_0)=-a(a+b)/b\neq0$, and $\Lnorm[f,1;y_0]=h_{(-a,-b,a+b)}(z)$ in $\bbQ(a,b)\fps z$.
\end{lem}

\begin{prop}[{\Lref{Proposition~3.20}: the Beta constant; analytic input}]\genusstatus{AI-proposed, unchecked}
\label{prop:row-beta}
\begin{equation}
\label{eq:row-beta}
\Abeta(z)=\Lnorm[\phi_0,1;\wA/E]=h_{(E,-\alpha_1,-\wA)}(z)\,\frac{E}{E+z}\qquad\text{in }\bbQ(\alpha_1,\wA)\fps z .
\end{equation}
\end{prop}

\begin{proof}[Sketch of proof of \cref{lem:row-betaphase,prop:row-beta}]
Both sides have $z$-coefficients rational in the parameters, so it suffices to take them real and positive. Then $\Bfn(a/z,b/z+1)=\int_{-\infty}^0\re^{f(y)/z}\dd y$ and $\Bfn(\wA/z+1,\alpha_1/z+1)=\int_0^1\re^{\phi_0(s)/z}\dd s$ \cite[5.12.1]{DLMF}, and by Laplace's method \Lref{Lemma~3.11} such an integral is asymptotic, as $z\to0^+$, to $\re^{f(u_0)/z}\big(2\pi z/(-f''(u_0))\big)^{1/2}\Lnorm[f,1;u_0]$ at the maximum $u_0$ of the phase. On the other hand Stirling's series with remainder, applied to $\log\Gamma(a/z)+\log\Gamma(b/z)-\log\Gamma((a+b)/z)$, gives the prefactor $\re^{f(u_0)/z}\big(2\pi z(a+b)/(ab)\big)^{1/2}$ times the exponential of $\sum_k\bk_kz^{2k-1}\big(a^{1-2k}+b^{1-2k}-(a+b)^{1-2k}\big)$, which is $\log h_{(-a,-b,a+b)}$ because $1-2k$ is odd. The shifts of the arguments by $1$ contribute, through $\Gamma(u+1)=u\Gamma(u)$, the factor $b/(a+b)$, independent of $z$, for the lemma, and $\wA\alpha_1/(E(E+z))$ for the proposition, where $(a,b)=(\wA,\alpha_1)$; this factor turns the prefactor into that of Laplace's method, times $E/(E+z)$ in the case of the proposition, the origin of the factor $E/(E+z)$. Uniqueness of asymptotic power series concludes.
\end{proof}

\begin{cor}[{\Lref{Corollary~3.23}: regularity of the $\tau$-coefficients; formal given \eqref{eq:row-beta}}]\genusstatus{AI-proposed, unchecked}
\label{cor:row-reg}
For every $k\geq0$, $\Kcoef_k\coloneqq\coef{\tau^k}(\Fhyp\Ephi\alg)$ and $\Hcoef_k\coloneqq\coef{\tau^k}(\Fhyp\Ephi)$ are rational functions of $z$ with coefficients in $\Kal$, without pole at $z=0$.
\end{cor}

Indeed $\Kcoef_k=\coef{\tau^k}(\Nlap/\Abeta)$ by \cref{thm:row-closed}, with $\Abeta\in1+z\bbC\fps z$ and $\Nlap\in\bbC\fps z\fps\tau$ (\cref{prop:row-wickrow}); the value of $\Abeta$ is not used here.

\begin{thm}[{\Lref{Theorem~3.24}: the frame-B singular row; formal given \eqref{eq:row-beta}}]\genusstatus{AI-proposed, unchecked}
\label{thm:row-main}
In $\bbC\fps z\fps t$, and coefficientwise in $z$ as analytic functions of $t$ near $0$,
\begin{equation}
\label{eq:row-main}
\rr_{0,1}(z;t)=\frac{E}{E+z}\,\Fhyp(z,\tau)\,\Ephi(z,\tau)\,\alg(\tau),\qquad \tau=1-\re^{-t}.
\end{equation}
\end{thm}

\begin{proof}[Sketch of proof]
By \cref{def:set-rows}, \cref{prop:row-wickrow} and \cref{thm:row-closed}, $\rr_{0,1}=\hw^{-1}\hnw\rpt_{0,1}=\hw^{-1}\hnw\Abeta\,\Fhyp\Ephi\alg$, and by \cref{prop:row-beta} and the multiset identity of \cref{lem:set-hmult}, $\hw^{-1}\hnw\Abeta=\hw^{-1}\hnw h_{(E,-\alpha_1,-\wA)}\,E/(E+z)=E/(E+z)$; this is where the class $\bw$ is forced. The passage from $\tau$ to $t$ is the principle (T) of \cref{lem:set-rings}.
\end{proof}

\begin{cor}[{\Lref{Corollary~3.25}: (F2) at $t=0$}]\genusstatus{AI-proposed, unchecked}
\label{cor:row-F2}
$\rr_{0,1}(z;0)=E/(E+z)$. Also $\rpt_{0,1}(z;0)=\Abeta(z)=h_{(E,-\alpha_1,-\wA)}(z)\,E/(E+z)$, and for the row $\rZ_{0,1}=\hnw\rpt_{0,1}$ with the $\cQ=0$ constant, $\rZ_{0,1}(z;0)=\hw(z)\,E/(E+z)$.
\end{cor}

This is the first half of (F2) of Theorem~A (\cref{thm:int-leaf}). A second, purely formal proof of the regularity, by eikonal and transport equations, and the first order of $\rr_{0,1}$ in $t$ are in \Lref{Lemma~3.26, Proposition~3.27 and Remark~3.28}.

% s4_leaves_a. First half of the abridged Section 4 (long version: Sections 4.1-4.5); label prefix: lea.
\section{The leaf structure at generic weights}
\label{sec:lea}

This section sketches the proof of Theorem~A (\cref{thm:int-leaf}). The weights are complex numbers satisfying (ND) of \cref{def:set-ND}, or indeterminates, and all factors are the frame-B factors of \cref{def:set-rows,def:set-factors} at their true idempotents. From \cref{sec:row} we use, in the notation $\Fhyp$, $\Ephi$, $\Hcoef_n$, $\Gauss$, $\alg$ introduced there: the closed form $\rr_{0,1}=(E/(E+z))\,\Fhyp\,\Ephi\,\alg$ (\cref{thm:row-main}), the regularity at $z=0$ of its $\tau$-coefficients (\cref{cor:row-reg}), the value $\rr_{0,1}(z;0)=E/(E+z)$ (\cref{cor:row-F2}), and $\Lv_1=\muq+\dphi'$ with $\muq\coloneqq E/\tau-\wA$ \Lref{Lemma~3.29}. The singular idempotent comes first; the regular colour, the legs, the assembly of the proof, the closed forms and the case $\alpha_0=0$ follow, from \cref{thm:lea-row2reg} on. Along the way the following is obtained.

\begin{thm}[{\Lref{Theorem~4.3}: quantum differential equation and symplecticity}]\genusstatus{AI-proposed, unchecked}
\label{thm:lea-qde}
\emergencystretch=2em
\begin{enumerate}
\item For $i=1,2$ the second row of the quantum differential equation holds to all orders in $z$: $z\Dop_i\vv_i=(\Aop-\Lv_i)\vv_i$, with $\vv_i$, $\Dop_i$ as in \cref{def:set-rows} and $\Aop\coloneqq\begin{bsmallmatrix}0&1\\-\alpha_1\alpha_2/\tau&E/\tau-\alpha_0\end{bsmallmatrix}$ acting on covectors (\cref{prop:lea-row2,thm:lea-row2reg}).
\item $\langle\vv_i(w),\vv_j(-w)\rangle=\delta_{ij}\Delta_i$ for $i,j\in\{1,2\}$, to all orders in $w$ (\cref{prop:lea-K1}, \cref{cor:lea-orth}, \Lref{Corollary~4.29}).
\item $\vv_2(z;0)=\hreg(z)\,(1,\alpha_1\alpha_2/E)$ with $\hreg=h_{(\alpha_1,\alpha_2,\wA,\wB,-E,-E)}$ (\cref{prop:lea-regrows}; the scalar factor is the analytic input there, the direction is formal).
\end{enumerate}
\end{thm}

%-------------------------------------------------------------------------------
\subsection{The key formula, Row 2 and the symplectic constant at the singular idempotent}
\label{sec:lea-qde}

\begin{prop}[{\Lref{Proposition~4.5}: the Hessian in the Euler coordinate; $\rho$; the key formula}]\genusstatus{AI-proposed, unchecked}
\label{prop:lea-key}
\leavevmode
\begin{enumerate}
\item $\Delta_1=\alpha_1\wA\epsilon_1\epsilon_2\,\cW_2(U_1)=-\epsilon_1\epsilon_2E^3\,\alg^2\,\tau^{-2}$.
\item $\rho=\tau/(t\alg)=1+\uq t/(2E^2)+O(t^2)$ is a unit of $\Oan$, and $t^2\rho^2\Delta_1=-\epsilon_1\epsilon_2E^3$ exactly ((F1) for $\Delta_1$).
\item (Key formula.)
\begin{equation}
\label{eq:lea-key}
\Gk(z;t)\coloneqq\rho\,\rr_{0,1}(z;t)=\frac{E}{E+z}\,\frac{\tau}{t}\,\Fhyp(z,\tau)\,\Ephi(z,\tau);
\end{equation}
$\Gk$ lies in $\bbC\fps{z,t}$, $\Gk(z;0)=E/(E+z)$, and every $\Gk_k\coloneqq\coef{t^k}\Gk$ is a rational function of $z$ without pole at $z=0$ whose Taylor series is $\coef{t^k}(\rho\,\rr_{0,1})$.
\end{enumerate}
\end{prop}

\begin{proof}[Sketch of proof]
(i) is the fourth form of \eqref{eq:set-Delta} with \cref{lem:row-hess}. So $\rho^{-2}=(t\alg/\tau)^2$, and multiplying \cref{thm:row-main} by $\rho$ cancels $\alg$; the properties of $\Gk_k$ are \cref{cor:row-reg} with (T) of \cref{lem:set-rings}.
\end{proof}

\begin{prop}[{\Lref{Proposition~4.7}: the second row of the quantum differential equation at the singular idempotent}]\genusstatus{AI-proposed, unchecked}
\label{prop:lea-row2}
To all orders in $z$, as identities of meromorphic functions of $t$ near $0$: $\rr_{1,1}=(\Lv_1+z\Dop_1)\rr_{0,1}$ (Row~1, which is the definition \eqref{eq:set-Hrow}), and $(\Lv_1+z\Dop_1)\rr_{1,1}=-(\alpha_1\alpha_2/\tau)\,\rr_{0,1}+(E/\tau-\alpha_0)\,\rr_{1,1}$ (Row~2). In matrix form $z\Dop_1\vv_1=(\Aop-\Lv_1)\vv_1$, and for $\leaVn\coloneqq t\rho\,\vv_1$, an element of $\bbC\fps{z,t}^2$,
\begin{equation}
\label{eq:lea-qde}
z\,\leaVn'=(\Aop-\Lv_1)\leaVn .
\end{equation}
\end{prop}

\begin{proof}[Sketch of proof]
By \cref{prop:lea-key}(ii), $\Lv_1+z\Dop_1=(t\rho)^{-1}(z\,\dd/\dd t+\Lv_1)(t\rho)$, and since $z\Ephi'=-\dphi'\Ephi$ and $\Lv_1-\dphi'=\muq$, conjugation by $\tau\Ephi$ turns $z\,\dd/\dd t+\Lv_1$ into $\Dz+\omq$, $\omq\coloneqq(E+z)/\tau-(\wA+z)$. By \eqref{eq:lea-key}, $t\rho\,\rr_{0,1}=(E/(E+z))\,\tau\Ephi\Fhyp$, hence $t\rho\,\rr_{1,1}=(E/(E+z))\,\tau\Ephi\,(\Dz+\omq)\Fhyp$, and Row~2 is the vanishing on $\Fhyp$ of the left side of the operator identity \Lref{Lemma~4.6}
\[
(\Dz+\omq)^2-\Big(\frac{E}{\tau}-\alpha_0\Big)(\Dz+\omq)+\frac{\alpha_1\alpha_2}{\tau}=\frac{(1-\tau)z^2}{\tau}\,\Gauss ,
\]
and $\Gauss[\Fhyp]=0$: Row~2 is Gauss's equation in operator form. The regular idempotent is not used.
\end{proof}

\emph{The operator $\Aop$} \Lref{Proposition~4.8}. $\Aop$ is the transpose of the matrix of $H\qprod$ in the basis $(1,H)$, and by \cref{prop:set-frob,prop:set-idem,lem:set-eig}: the covectors $\chiv_i\coloneqq(1,\Lv_i)$ satisfy $\Aop\chiv_i=\Lv_i\chiv_i$ and $\langle\chiv_i,\chiv_j\rangle=\delta_{ij}\Delta_i$; $\Aop$ is self-adjoint for \eqref{eq:set-pairing}, $\langle\Aop a,b\rangle=\langle a,\Aop b\rangle$; and $\Aop-\Lv_1$ is $(\Lv_2-\Lv_1)$ times the orthogonal projector onto $\bbC\chiv_2$, of rank one: with $\pim(a)\coloneqq\langle a,\chiv_2\rangle$,
\begin{equation}
\label{eq:lea-proj}
\langle(\Aop-\Lv_1)a,b\rangle=\frac{\pim(a)\pim(b)}{\Thf}\quad\text{for all covectors }a,b,\qquad
\Thf\coloneqq\frac{\Delta_2}{\Lv_2-\Lv_1}=-\epsilon_1\epsilon_2\wB\tau\sast\Lv_2 .
\end{equation}

\begin{prop}[{\Lref{Proposition~4.9}: the symplectic constant $\Kone=1$}]\genusstatus{AI-proposed, unchecked}
\label{prop:lea-K1}
Let $\Kone(z)\coloneqq\langle\vv_1(z),\vv_1(-z)\rangle/\Delta_1$, $\brow(\psi)\coloneqq\rho\,\vv_1(-\psi)$ and $\gl(\psi;t)\coloneqq\Gk(-\psi;t)$.
\begin{enumerate}
\item The series $\Kone$ lies in $\bbC\fps{z,t}$ and is independent of $t$.
\item $\brow=\big(\gl,\ (\Lv_1-\psi/t)\gl-\psi\gl'\big)$, the series $t\brow$ lies in $\bbC\fps{\psi,t}^2$, $\coef{t^{-1}}\brow=(0,E)$, and $\gl|_{t=0}=E/(E-\psi)=E\lf_1(\psi)$.
\item $\Kone=1$, that is, $\langle\vv_1(w),\vv_1(-w)\rangle=\Delta_1$ to all orders in $w$.
\item Consequently the subtracted term in \eqref{eq:set-edge} is $\Pij_{11}(-w,w)=-\Delta_1$; moreover $\langle\brow(\psi),\brow(-\psi)\rangle=\rho^2\Delta_1=-\epsilon_1\epsilon_2E^3/t^2$ and $-\psi\,\frac{\dd}{\dd t}\big(t\brow(\psi)\big)=(\Aop-\Lv_1)\big(t\brow(\psi)\big)$.
\end{enumerate}
\end{prop}

\begin{proof}[Sketch of proof]
(i) $\Kone=-\langle\leaVn(z),\leaVn(-z)\rangle/(\epsilon_1\epsilon_2E^3)$, and its $t$-derivative vanishes by \eqref{eq:lea-qde} at $z$ and at $-z$ and the self-adjointness of $\Aop$. (ii) is Row~1 with the key formula: $t\rho\,\rr_{1,1}=(t\Lv_1+z)\Gk+zt\Gk'$, with value $(E+z)\cdot E/(E+z)=E$ at $t=0$, since $t\Lv_1\to E$; the $z$-dependence cancels because of the geometric form of (F2). (iii) By (ii), $\Kone|_{t=0}=-\langle(0,E),(0,E)\rangle/(\epsilon_1\epsilon_2E^3)=1$; conclude by (i). (iv) follows from (iii) and \eqref{eq:lea-qde} at $z=-\psi$.
\end{proof}

%-------------------------------------------------------------------------------
\subsection{The dilaton leaves: (L1)}
\label{sec:lea-L1}

This subsection and the next are written over $\Kal$, by (P) of \cref{lem:set-rings}.

\begin{thm}[{\Lref{Proposition~4.15 and Theorem~4.16}: (L1), all $k$}]\genusstatus{AI-proposed, unchecked}
\label{thm:lea-L1}
Put $\Rk_k(z)\coloneqq\coef{t^k}(z\Gk)$ and $q^{(k)}\coloneqq\coef{t^k}(\rho\qt)$. Then $\Rk_0=Ez/(E+z)$ and $q^{(0)}=\qo=E-E^2\lf_1$. For $k\geq1$ there are $a^{(k)}_j\in\Kal$, $1\leq j\leq k+1$, with
\[
\Rk_k(z)=\sum_{j=1}^{k+1}\frac{a_j^{(k)}}{E+jz},\qquad a_j^{(k)}=j\Res_{z=-E/j}\Rk_k,\qquad \sum_{j=1}^{k+1}a_j^{(k)}=0,\qquad
q^{(k)}=\sum_{j=1}^{k+1}a_j^{(k)}\lf_j\in\Lsp_{k+1}.
\]
Moreover $\gl_0=E\lf_1$ and, for $k\geq1$, $\gl_k\coloneqq\coef{t^k}\gl=-\frac1E\sum_{j=1}^{k+1}j\,a_j^{(k)}\lf_j\in\Lsp_{k+1}$, with $\gl_k=O(\psi^{-2})$ at $\psi=\infty$.
\end{thm}

\begin{proof}[Sketch of proof]
As $\qt=-\psi\,\rr_{0,1}(-\psi)$, $q^{(k)}(\psi)=\Rk_k(-\psi)$; by \eqref{eq:lea-key} and (T), $\Rk_k=\frac{Ez}{E+z}\sum_{n\leq k}d_{k,n}\Hcoef_n$ with $d_{k,n}=\coef{t^k}(\tau^{n+1}/t)$. Three facts give the result \Lref{Lemmas~4.12 and~4.13, Remark~4.14}. \emph{Poles away from $0$}: they are simple and lie in $\{-E/j:1\leq j\leq k+1\}$; they come from the coefficients $\Fhyp_n$ of $\Fhyp$ ($2\leq j\leq n+1$) and from $Ez/(E+z)$, the coefficients of $\Ephi$ being polynomials in $1/z$. \emph{No pole at $z=0$}: this is \cref{cor:row-reg}; the separate terms do have poles there. \emph{Value at $\infty$}: $\Fhyp_n(\infty)=1/(n+1)$, independently of the weights, and $\coef{\tau^m}\Ephi\to\delta_{m0}$, so $\Rk_k(\infty)=E\,\coef{t^k}\big(t^{-1}\sum_{n\geq0}\tau^{n+1}/(n+1)\big)=E\,\delta_{k0}$, because $-\log(1-\tau)=t$; this is how the normalisation by $\rho$ removes the constant leaf. Now (PF) of \cref{lem:set-leafid} gives the expansion of $\Rk_k$, $\Rk_k(0)=0$ the sum rule, and $\Gk_k=\Rk_k/z$ the formula for $\gl_k$.
\end{proof}

In particular every $\Dil_1(d)$ is analytic at $t=0$ ((F3) for $\Dil_1$), $\qt|_{t=0}=\qo$, and $\qt,\rho\qt\in\Dom$.

%-------------------------------------------------------------------------------
\subsection{The \texorpdfstring{$H$}{H}-row in leaf form and the \texorpdfstring{$\Edge_{11}$}{E11} flow: (F3) for \texorpdfstring{$\Edge_{11}$}{E11} and (L3)}
\label{sec:lea-flow}

Let $\Dg\coloneqq\frac{\dd}{\dd t}(t\gl)=\sum_k(k+1)\gl_kt^k$, so that $\brow=(\brow_0,\brow_1)=\big(\gl,\ \Lv_1\gl-\frac{\psi}{t}\Dg\big)$ by \cref{prop:lea-K1}(ii), and write $\Lv_1=E/t+\sum_{m\geq0}\lrem_mt^m$ (\cref{lem:set-eig}).

\begin{prop}[{\Lref{Propositions~4.17 and~4.18}: the leaf structure of $\brow$ and its projection}]\genusstatus{AI-proposed, unchecked}
\label{prop:lea-sigma}
$\coef{t^{-1}}\brow=(0,E)$, and for $k\geq0$ both components of $\brow^{(k)}\coloneqq\coef{t^k}\brow$ lie in $\Lsp_{k+1}$: $\brow^{(k)}_0=\gl_k$ and $\brow^{(k)}_1=\big(E-(k+2)\psi\big)\gl_{k+1}+\sum_{m=0}^{k}\lrem_m\,\gl_{k-m}$. In particular $\brow^{(0)}=(E\lf_1,\ \alpha_1\alpha_2\lf_1)$; with $\coef{t^{-1}}\brow$ this is the second half of (F2). Moreover $\pim(\brow(\psi))=-\Thf\,\frac{\psi}{t}\,\Dg(\psi)$.
\end{prop}

\begin{proof}[Sketch of proof]
By (LAM3), the factor $E-(k+2)\psi$ turns the top leaf $\lf_{k+2}$ of $\gl_{k+1}=\sum_jg_j\lf_j$ into a constant, and the constants add up to $(k+2)\sum_jg_j/j$, which vanishes because $\gl_{k+1}=O(\psi^{-2})$ (\cref{thm:lea-L1}, (LAM4)); this works because the coefficient $E$ of the pole of $\Lv_1$ equals the leaf scale. For the projection, $\brow=\gl\,\chiv_1-(\psi/t)\Dg\,(0,1)$, $\pim(\chiv_1)=0$ and $\pim((0,1))=\Thf$.
\end{proof}

By \cref{prop:lea-K1}(iv) and \eqref{eq:set-edge}, with $\Yv\coloneqq t\brow\in\Kal\fps{t}\fps{\psi}^2$,
\[
\Ebar(\psi,\psi')\coloneqq t^2\rho^2\Edge_{11}(\psi,\psi')=\frac{-\epsilon_1\epsilon_2E^3-\langle\Yv(\psi),\Yv(\psi')\rangle}{\psi+\psi'} ,
\]
symmetric, and in $\Kal\fps{t}\fps{\psi,\psi'}$ as the numerator vanishes on $\psi'=-\psi$; so $\Edge_{11}(a,b)=\Edge_{11}(b,a)$.

\begin{thm}[{\Lref{Theorem~4.20}: the $\Edge_{11}$ flow}]\genusstatus{AI-proposed, unchecked}
\label{thm:lea-flow}
In $\Kal\fps{t}\fps{\psi,\psi'}$,
\begin{equation}
\label{eq:lea-flow}
\frac{\dd}{\dd t}\Ebar(\psi,\psi')=\Thf(t)\,\Dg(\psi)\,\Dg(\psi'),\qquad
\Thf=\frac{\Delta_2}{\Lv_2-\Lv_1}=\epsilon_1\epsilon_2(\alpha_1\alpha_2-E\Lv_2)=\epsilon_1\epsilon_2(E\dphi'-\alpha_2\wA).
\end{equation}
\end{thm}

\begin{proof}[Sketch of proof]
By \cref{prop:lea-K1}(iv), $-\psi\,\Yv(\psi)'=(\Aop-\Lv_1)\Yv(\psi)$. So, by self-adjointness, $\psi\psi'$ times the $t$-derivative of the numerator is $(\psi+\psi')\langle(\Aop-\Lv_1)\Yv(\psi),\Yv(\psi')\rangle$: the denominator cancels, and by \eqref{eq:lea-proj} and \cref{prop:lea-sigma} the pairing is $\pim(\Yv(\psi))\pim(\Yv(\psi'))/\Thf=\Thf\,\psi\psi'\,\Dg(\psi)\Dg(\psi')$. The right side of \eqref{eq:lea-flow} is a product because $\Aop-\Lv_1$ has rank one. The forms of $\Thf$ are in \cref{lem:set-eig}; the last uses \eqref{eq:set-envelope}.
\end{proof}

\emph{Integration from $t=0$; (F3) for $\Edge_{11}$.} Since $\coef{t^0}\Yv=(0,E)$ and $\langle(0,E),(0,E)\rangle=-\epsilon_1\epsilon_2E^3$, one has $\coef{t^0}\Ebar=0$, so $\Ebar=\int_0^t\Thf\,\Dg(\psi)\Dg(\psi')\,\dd u$ termwise. Moreover $\Thf=O(t)$, as $\Delta_2$ is regular while $\Lv_2-\Lv_1=-\delta/\tau$ has a simple pole; with $\Thf_k\coloneqq\coef{t^k}\Thf$, $\Thf_1=-\epsilon_1\epsilon_2\alpha_1\alpha_2\wA\wB/E^2$. Hence $\Ebar\in t^2\Kal\fps{t}\fps{\psi,\psi'}$, and $\Edge_{11}=(t\rho)^{-2}\Ebar$ is regular at $t=0$.

\begin{thm}[{\Lref{Theorem~4.21}: (L3), all $k$}]\genusstatus{AI-proposed, unchecked}
\label{thm:lea-L3}
For every $k\geq0$,
\[
\coef{t^k}(\rho^2\Edge_{11})(\psi,\psi')=\frac{1}{k+2}\sum_{\substack{k_0\geq1,\ k_1,k_2\geq0\\ k_0+k_1+k_2=k+1}}\Thf_{k_0}(k_1+1)(k_2+1)\,\gl_{k_1}(\psi)\gl_{k_2}(\psi') ,
\]
which lies in $\mathrm{span}\{\lf_i(\psi)\lf_l(\psi'):i,l\geq1,\ (i-1)+(l-1)\leq k\}$ and is symmetric under $\psi\leftrightarrow\psi'$. At $t=0$ the loop has rank one:
$\Edge_{11}(z,w)|_{t=0}=-\tfrac12\,\epsilon_1\epsilon_2\alpha_1\alpha_2\wA\wB\,\lf_1(z)\lf_1(w)$.
\end{thm}

\begin{proof}[Sketch of proof]
$\coef{t^k}(\rho^2\Edge_{11})=\coef{t^{k+2}}\Ebar=\coef{t^{k+1}}\big(\Thf\,\Dg(\psi)\Dg(\psi')\big)/(k+2)$ by the integrated flow. Every term has $k_1+k_2\leq k$, and $\gl_{k_i}\in\Lsp_{k_i+1}$ by \cref{thm:lea-L1}.
\end{proof}

As $\rho^{e-2}\in\Kal\fps{t}$, the same support bound holds for $\coef{t^k}(\rho^e\Edge_{11})$, for every integer $e$; and since $f_{1,b}=\coef{\psi'^b}\Edge_{11}$, (LAM1) shows that $\coef{t^k}(\rho f_{1,b})$ and $\coef{t^k}f_{1,b}$ lie in $\Lsp_{k+1}$ for every $b\geq0$: this is (L2) for $j=1$ \Lref{Corollary~4.22}.

% s4_leaves_b. Second half of the abridged Section 4 (continues s4_leaves_a.tex); label prefix: lea.

\subsection{The regular colour}
\label{sec:lea-regular}

The factors of idempotent $2$ are those of idempotent $1$ transported by the involution $U\mapsto U_1+U_2-U$ of $\QA$ (\cref{def:set-QA}), which exchanges the characters $\ev_1,\ev_2$ and commutes with $\dd/\dd\cQ$ and with the $t$-independent scalars $\hnw$, $\hw^{-1}$, the same series at both idempotents (this is where ``the same scalar at both colours'' is needed): in every defining expression $U_1$ is replaced by $U_2$ \Lref{Lemma~4.23}. Since $U_2$ is analytic at $t=0$ with $U_2(0)=-\alpha_2/\wB$ and $1-U_2(0)=E/\wB\neq0$, one gets \Lref{Lemma~4.24}: \textbf{(R)} every $\rpt_{0,2}[m]$, $\rr_{0,2}[m]$, $\rr_{1,2}[m]$ is analytic at $t=0$; hence so are $\Dil_2(n)$ and $\Edge_{22}(m,n)$ ((F3) for these factors). Row~1 at $i=2$ is the definition \eqref{eq:set-Hrow}; Row~2 needs a separate argument, since the closed form concerns the singular row only and $\stil\to\infty$ as $\tau\to0$.

\begin{thm}[{\Lref{Theorem~4.26}: Row 2 of the quantum differential equation at the regular idempotent}]\genusstatus{AI-proposed, unchecked}
\label{thm:lea-row2reg}
In $\Man\fps{z}$ (in $\Oan\fps{z}$ after multiplication by $\tau$),
\[
(\Lv_2+z\Dop_2)\rr_{1,2}=-\frac{\alpha_1\alpha_2}{\tau}\rr_{0,2}+\Bigl(\frac E\tau-\alpha_0\Bigr)\rr_{1,2},\qquad\text{equivalently}\qquad z\Dop_2\rr_{1,2}=\Lv_1(\rr_{1,2}-\Lv_2\rr_{0,2}).
\]
With Row~1 this is $z\Dop_2\vv_2=(\Aop-\Lv_2)\vv_2$, for the frame-B rows, the rows $\rZ$ and the point-frame rows alike.
\end{thm}

\begin{proof}[Sketch of proof]
Work coefficientwise in $z$ near an arbitrary $\tau_0\in\Ddisc''$, away from $\tau=0$, where $\stil$ is a non-degenerate critical point of $\phi_\tau$. Put $f_0\coloneqq\phi_\tau(\stil)$, $f_2\coloneqq\phi_\tau''(\stil)$, $\omq\coloneqq(E+z)/\tau-(\wA+z)$, $\leaDt\coloneqq\dd/\dd t+f_0'/z-f_2'/(2f_2)$.
(A) As in \cref{prop:row-wickrow}, by (P1) of \cref{prop:row-calc}, $\rpt_{0,2}=U_2\,\leaLt$ with $\leaLt\coloneqq\Lnorm[\phi_\tau,(1-\tau s)^{-1};\stil]$.
(B) $(\Lv_2+z\Dop_2)\circ U_2=U_2\circ(z\leaDt+\omq)$, by the envelope formula $\Lv_2=f_0'+E/\tau-\wA$ at $\stil$ and by $\Delta_2=\alpha_1\wA\epsilon_1\epsilon_2U_2^2f_2/\tau^2$ (the fourth form of \eqref{eq:set-Delta} at $U_2$).
(C) $\leaLt$ satisfies Gauss's equation of \cref{prop:row-gaussode} with $\dd/\dd\tau$ replaced by $\leaDt/(1-\tau)$, by the argument given there at $\stil$: (P2), the amplitude identity \Lref{Lemma~3.15} and (P3).
(D) By the operator identity \Lref{Lemma~4.6}, valid for every symbol with the Leibniz rule of $z\,\dd/\dd t$, (C) says that $(z\leaDt+\omq)^2-(E/\tau-\alpha_0)(z\leaDt+\omq)+\alpha_1\alpha_2/\tau$ annihilates $\leaLt$; by (A), (B) the same polynomial in $\Lv_2+z\Dop_2$ annihilates $\rpt_{0,2}$, hence $\rZ_{0,2}$ and $\rr_{0,2}$. This is Row~2 on $\Ddisc''$, and in $\Man$ by (R).
\end{proof}

\begin{cor}[{\Lref{Corollary~4.27}: orthogonality}]\genusstatus{AI-proposed, unchecked}
\label{cor:lea-orth}
$\langle\vv_1(w),\vv_2(-w)\rangle=0$ in $\Man\fps{w}$, to all orders in $w$.
\end{cor}

\begin{proof}[Sketch of proof]
Let $Y(z)\coloneqq\langle\vv_1(z),\vv_2(-z)\rangle=\sum_mY_mz^m$. By $z\Dop_i\vv_i=(\Aop-\Lv_i)\vv_i$ for $i=1,2$ (\cref{prop:lea-row2,thm:lea-row2reg}) and the self-adjointness of $\Aop$ \Lref{Proposition~4.8}, $zY'=\big[(\Lv_2-\Lv_1)+\tfrac z2(\Delta_1'/\Delta_1+\Delta_2'/\Delta_2)\big]Y$. As $\Lv_2-\Lv_1=-\delta/\tau$ is a unit of $\Man$ and $Y$ is a power series in $z$, $Y_0=0$ and inductively $Y_m=0$.
\end{proof}

\begin{prop}[{\Lref{Proposition~4.28}: the regular rows at $t=0$}]\genusstatus{AI-proposed, unchecked}
\label{prop:lea-regrows}
\begin{enumerate}
\item[(a)] (Analytic input.) In the point frame $\rpt_{0,2}(z;0)=h_{(\alpha_2,-E,\wB)}(z)$; equivalently, in frame B, $\rr_{0,2}(z;0)=\hreg(z)$. Here $\hreg\coloneqq h_{(\alpha_1,\wA,-E)}h_{(\alpha_2,-E,\wB)}$ is the series \eqref{eq:set-h} of the multiset $(\alpha_1,\alpha_2,\wA,\wB,-E,-E)$, and $\hreg(z)\hreg(-z)=1$.
\item[(b)] $\vv_2(z;0)=\hreg(z)\,(1,\alpha_1\alpha_2/E)$; equivalently $\coef{t^1}\rr_{0,2}(z;t)=0$.
\end{enumerate}
\end{prop}

\begin{proof}[Sketch of proof]
(a) By (R) the value at $t=0$ is the normalised expansion, in the coordinate $\log U$, of the superpotential at $\cQ=1$, $E\log(1-U)-\alpha_2\log U$, at its critical point $-\alpha_2/\wB$: this is \cref{lem:row-betaphase} with $(a,b)=(-\alpha_2,E)$, whose hypothesis $ab(a+b)\neq0$ reads $\alpha_2E\wB\neq0$, the second use of input (i) of \cref{sec:int-inputs}; frame B multiplies by $h_{(\alpha_1,\wA,-E)}$. (b) In Row~2, read as $(\Dop_2\rr_{1,2})[m-1]=\Lv_1(\rr_{1,2}[m]-\Lv_2\rr_{0,2}[m])$, the left side is analytic at $t=0$ and $\Lv_1=E/t+O(1)$, so the bracket vanishes at $t=0$; and $\Lv_2(0)=\alpha_1\alpha_2/E$.
\end{proof}

Moreover $\Ktwo(z)\coloneqq\langle\vv_2(z),\vv_2(-z)\rangle/\Delta_2$ is independent of $t$, as $\Kone$ in \cref{prop:lea-K1}, and equals $\hreg(z)\hreg(-z)=1$ by (a), (b) \Lref{Corollary~4.29}; with \cref{cor:lea-orth} this completes \cref{thm:lea-qde}.

\subsection{The legs and the assembly of Theorem~A}
\label{sec:lea-legs}

Two formal facts are used. \emph{Extraction of a leg} \Lref{Lemma~4.31}: for $\Edge\in\bbC\fps{\psi,w}$ and $b\geq0$, in $\bbC\lser{\psi}$,
\begin{equation}
\label{eq:lea-legext}
\coef{w^b}\Edge=\sum_{n=0}^{b}(-1)^n\psi^{-n-1}\coef{w^{b-n}}\big((\psi+w)\Edge\big).
\end{equation}
\emph{The leaf form of the singular covector} \Lref{Proposition~4.17}: $\brow(\psi)=\rho\vv_1(-\psi)$ (\cref{prop:lea-K1}) has the expansion $\brow=t^{-1}(0,E)+\sum_{k\geq0}t^k\brow^{(k)}$ with both components of $\brow^{(k)}$ in $\Lsp_{k+1}$.

\begin{prop}[{\Lref{Proposition~4.32}: (L2) for $j=2$ and the regularity of $\Edge_{12}$}]\genusstatus{AI-proposed, unchecked}
\label{prop:lea-L2}
Assume \textup{(R)}. Then \textup{(i)} $\rho\Edge_{12}$ and $\Edge_{12}$ are $O(t^0)$, that is, every $\Edge_{12}(a,b)$ is regular at $t=0$ ((F3) for $\Edge_{12}$); and \textup{(ii)} for all $k\geq0$ and $b\geq0$, $\coef{t^k}(\rho f_{2,b})$ lies in $\Lsp_{k+1}$.
\end{prop}

\begin{proof}[Sketch of proof]
Put $\leaNtw(\psi,w)\coloneqq\rho\Pij_{12}(\psi,w)=-\langle\brow(\psi),\vv_2(-w)\rangle$; by \eqref{eq:set-edge}, $(\psi+w)\rho\Edge_{12}$ is the difference $\leaNtw(\psi,w)-\leaNtw(-w,w)$. (i) By (R) the only negative power of $t$ in $\leaNtw$ is $t^{-1}$, with coefficient $-\langle(0,E),\vv_2(-w;0)\rangle$, independent of $\psi$; it cancels in the difference. (ii) By \eqref{eq:lea-legext}, $\coef{t^k}(\rho f_{2,b})=\sum_{n\leq b}(-1)^n\psi^{-n-1}Y_{b-n}(\psi)$, where, by (R) and the leaf form of $\brow$, each $Y_m$ is a constant plus an element of $\Lsp_{k+1}$. The right side is a rational function of $\psi$ with at most simple poles at $E/j$, $j\leq k+1$, vanishing at $\infty$, whose apparent pole at $\psi=0$ cancels because the left side is a power series in $\psi$; by partial fractions (\cref{lem:set-leafid}) it lies in $\Lsp_{k+1}$.
\end{proof}

\emph{Transposition} \Lref{Proposition~4.34}. By \cref{cor:lea-orth}, $\Pij_{12}(-u,u)=0$, so $\Edge_{12}(z,w)=\Pij_{12}(z,w)/(z+w)$, and by the symmetry of the pairing $\Edge_{21}(z,w)=\Edge_{12}(w,z)$, that is, $\Edge_{21}(b,a)=\Edge_{12}(a,b)$. Hence $\Edge_{21}$ is regular at $t=0$, and an edge joining a singular vertex to a regular one gives the singular vertex the leg $f_{2,b}$ for either orientation; this is where Row~2 at the regular idempotent is needed. By (T) of \cref{lem:set-rings}, $\coef{t^k}f_{j,b}$ lies in $\Lsp_{k+1}$ as well, for $j\in\{1,2\}$ \Lref{Corollary~4.35}.

\begin{proof}[Proof of Theorem~A (\cref{thm:int-leaf}) from the results above; {\Lref{Section~4.8}}]
(F1): \cref{prop:lea-key} for $\Delta_1$ and $\rho$, \cref{lem:set-eig} for $\Delta_2$. (F2): \cref{cor:row-F2} and the orders $t^{-1}$, $t^0$ of the leaf form of $\brow$ at $\psi=-z$. (F3): $\Dil_1$ by \cref{thm:lea-L1}; $\Edge_{11}$ by the flow integrated from $t=0$ (\cref{thm:lea-flow}); $\Dil_2$, $\Edge_{22}$ by (R); $\Edge_{12}$ by \cref{prop:lea-L2}; $\Edge_{21}$ by the transposition. (L1): \cref{thm:lea-L1}. (L2): $j=1$ from \cref{thm:lea-L3} by (LAM1) of \cref{lem:set-leafid}, $j=2$ by \cref{prop:lea-L2}. (L3): \cref{thm:lea-L3}, and (T) for the powers $\rho^e$. Rationality in the weights is (P) of \cref{lem:set-rings}; the sharpness is discussed after \cref{thm:lea-closed}.
\end{proof}

\emph{Where the inputs enter} \Lref{Remark~4.37}. The analytic input (i) of \cref{sec:int-inputs} enters through the key formula of \cref{prop:lea-key}, hence in everything that uses $E/(E+z)$: (F2), (L1) at $t^0$, $\Kone=1$, (F3) for $\Edge_{11}$. Condition (ND) is used through $E\neq0$ (the leaf scale), $\alpha_1\wA\neq0$ (the non-degenerate Euler critical point), $\alpha_2\wB\neq0$ ($U_2(0)$ finite and non-zero, $\Delta_2(0)\neq0$) and $\epsilon_1\epsilon_2\neq0$ ($\eta$ invertible). The point $\cQ=0$ enters only through the definition of $\hnw$ (\cref{lem:ide-norm}).

\begin{rmk}[{\Lref{Remark~4.2}: where the Calabi--Yau condition enters}]
\label{rmk:lea-cy}
No step of the proof of Theorem~A uses $E=\epsilon_1+\epsilon_2$: the frame-B rows do not involve $\epsilon_1,\epsilon_2$ (\cref{lem:set-hmult}), and $\Delta_i$, $\Edge_{ij}$ and the pairing carry $\epsilon_1\epsilon_2$ as an overall factor. Theorem~A therefore holds at every weight point of the five-torus with (ND), with $\bw\coloneqq(E,-\epsilon_1,-\epsilon_2)$ and $E\coloneqq\alpha_0+\alpha_1+\alpha_2$. The Calabi--Yau condition makes $\bw$ a sum-zero triple whose first weight is the leaf scale $E$, which is what the two cited theorems of \cite{LP2} (\cref{thm:reg-L,thm:reg-VI}) need: it is needed for the regularity theorem built on Theorem~A, not for Theorem~A itself.
\end{rmk}

\subsection{Closed forms, sharpness, and the case \texorpdfstring{$\alpha_0=0$}{alpha0 = 0}}
\label{sec:lea-closed}

\begin{thm}[{\Lref{Theorem~4.38}, part of \Lref{Theorem~4.4}: closed form of the dilaton leaf coefficients}]\genusstatus{AI-proposed, unchecked}
\label{thm:lea-closed}
Let $a_j^{(k)}$ be the leaf coefficients of \cref{thm:lea-L1}, with $a_1^{(0)}\coloneqq-E^2$ and $a_j^{(k)}\coloneqq0$ for $j>k+1$. Put $\Psh_k\coloneqq\prod_{i=1}^{k}\big((k+1)\alpha_2-iE\big)\big((k+1)\wA-iE\big)$ ($\Psh_0=1$), $\Cj_j\coloneqq\Psh_{j-1}/(j!\,(j-1)!\,E^{2j-4})$, $\zq\coloneqq\tau\re^{\dphi/E}=t+O(t^2)$ and, for $j\geq1$,
\[
\Aleaf_j(t)\coloneqq(-1)^j\,\Cj_j\,t^{-1}\zq(t)^j\,{}_2F_1\Bigl(\frac{j\wB}{E},\frac{j\alpha_1}{E};\,j;\,\tau\Bigr)=O(t^{j-1}).
\]
Then $a_j^{(k)}=\coef{t^k}\Aleaf_j$ for all $j\geq1$, $k\geq0$; equivalently $\rho\,\qt=E+\sum_{j\geq1}\Aleaf_j(t)\lf_j$, as formal series in $t$ with finitely many $j$ at each order.
\end{thm}

\begin{proof}[Sketch of proof]
By \cref{thm:lea-L1}, $a_j^{(k)}=j\Res_{z=-E/j}\coef{t^k}(z\rho\,\rr_{0,1})$, and $z\rho\,\rr_{0,1}=(\tau/t)\frac{Ez}{E+z}\Ephi\sum_n\Fhyp_n\tau^n$ by the key formula of \cref{prop:lea-key}, $\Fhyp_n$ being the $\tau$-coefficients of $\Fhyp$. At $z=-E/j$, $j\geq2$, $\Fhyp_n$ has a simple pole exactly for $n\geq j-1$, from the factor $E+jz$ of its denominator; after the reindexing $n=j-1+m$ the residues sum to a constant times $\tau^{j-1}{}_2F_1(j\wB/E,j\alpha_1/E;j;\tau)$, and $\Ephi(-E/j)=\re^{j\dphi/E}$. For $j=1$ the pole is that of $Ez/(E+z)$.
\end{proof}

Inserting these series in the integrated flow of \cref{thm:lea-flow} gives every loop coefficient in closed form \Lref{Theorem~4.41}. \emph{Sharpness} \Lref{Corollaries~4.40 and~4.42}: the top coefficient is $a_{k+1}^{(k)}=(-1)^{k+1}\Cj_{k+1}$, so the top leaf $\lf_{k+1}$ occurs in $\coef{t^k}(\rho\qt)$ if and only if $\Psh_k\neq0$; and the coefficient of $\lf_i(\psi)\lf_l(\psi')$ in $\coef{t^k}(\rho^2\Edge_{11})$ with $(i-1)+(l-1)=k$ is non-zero if and only if $\Psh_{i-1}\Psh_{l-1}\neq0$. So the bounds of (L1) and (L3) are sharp at generic weights.

\begin{cor}[{\Lref{Corollary~4.45}: $\alpha_0=0$}]\genusstatus{AI-proposed, unchecked}
\label{cor:lea-a0}
Let $\alpha_0=0$ and $\alpha_1\alpha_2E\epsilon_1\epsilon_2\neq0$ (then $\wA=\alpha_2$, $\wB=\alpha_1$, $E=\alpha_1+\alpha_2$), and put $\xc\coloneqq\alpha_1/\alpha_2$, $\kap\coloneqq\epsilon_1\epsilon_2\alpha_2^4/\alpha_1$. Let the superscript $\TFsup$ denote the frame-B objects at the fibrewise weight point $(\alpha_0,\alpha_1,\alpha_2,\epsilon_1,\epsilon_2)=(0,\xc,1,\xc,1)$. Then, for $i,j\in\{1,2\}$ and all $t$,
\begin{gather*}
\rr_{0,i}(z;t)=\rr^{\TFsup}_{0,i}(z/\alpha_2;t),\qquad \rr_{1,i}(z;t)=\alpha_2\,\rr^{\TFsup}_{1,i}(z/\alpha_2;t),\qquad \Lv_i=\alpha_2\Lv_i^{\TFsup},\qquad \Delta_i=\kap\,\Delta_i^{\TFsup},\\
\rho=\rho^{\TFsup},\qquad \Dil_i(d)=\alpha_2^{1-d}\Dil_i^{\TFsup}(d),\qquad \Edge_{ij}(m,n)=\kap\,\alpha_2^{-1-m-n}\,\Edge_{ij}^{\TFsup}(m,n),\qquad \lf_j(\psi)=\alpha_2^{-1}\lf_j^{\TFsup}(\psi/\alpha_2).
\end{gather*}
Consequently Theorem~A at $\alpha_0=0$ is Theorem~A at the fibrewise point $(0,\xc,1,\xc,1)$ read through these rescalings.
\end{cor}

\begin{proof}[Sketch of proof]
At $\alpha_0=0$ the superpotential is $\alpha_2$ times the fibrewise one. The normalised expansions obey the rescaling rule $\Lnorm[\lambda f,1](z)=\Lnorm[f,1](z/\lambda)$ \Lref{Definition~3.7}, so $\rpt_{0,i}(z)={\rpt}^{\TFsup}_{0,i}(z/\alpha_2)$; likewise $h_{(\alpha_1,\alpha_2,-E)}(z)=h_{(\xc,1,-1-\xc)}(z/\alpha_2)$ for the frame-B twist. The pairing \eqref{eq:set-pairing} on covectors $(a_0,\alpha_2a_1)$ and the second form of \eqref{eq:set-Delta} scale by $\kap$.
\end{proof}

The direct proof of Theorem~A covers $\alpha_0=0$ as well: it never uses the point $\cQ=0$, and $\hnw$ enters it only through the multiset identity of \cref{lem:set-hmult}, whereas the description of $\hnw$ as the constant at $\cQ=0$ (\cref{lem:ide-norm}) needs $\alpha_0\neq0$. The rescaling to a fibrewise point is what \cref{cor:rig-3T} uses.

% s5_identification. Abridged Section 5 (identification of the graph sum with the potential); label prefix: ide.
\section{Identification of the graph sum with the Gromov--Witten potential}
\label{sec:ide}

This section sketches the proof of Theorem~B (\cref{thm:int-ident}; \Lref{Proposition~5.1}). The weights are the generic point of $\TCY$; no specialisation and no analytic continuation in $\cQ$ is used. On the three-torus $\Tthree$ ($\alpha_0=0$), which contains the fibrewise torus $\TF$, the point $\cQ=0$ is not semisimple and the argument of this section does not apply directly \Lref{Remark~5.18}; that case is treated in \cref{sec:rig} by a ring map (\cref{lem:rig-3Tid}). The two roots of $\Pq$ are labelled here by their behaviour at $\cQ=0$ (\cref{lem:set-limits}(ii)): $U_+\to-\alpha_2/\alpha_0$ and $\ideQU\coloneqq\cQ U_-\to\alpha_0/\wB$, both in $\KCY\fps{\cQ}$; $i,j\in\{+,-\}$. No matching with the labels $1,2$ fixed at $\cQ=1$ is needed (\cref{lem:ide-terms}(iii)).

Let $\Omega=(\Omega_{g,n})$ be the $\TCY$-equivariant Gromov--Witten classes of $Z$ (\cref{def:set-inv}), so that $\sF_g=\int_{\Mbar_g}\Omega_{g,0}$, and let $\ideV\coloneqq\KCY\ideone\oplus\KCY H$. By \Lref{Lemma~5.2}, which rests on \cite{GP99}, on \cite[App.~1~(iv)--(vi)]{CG07} and on the fundamental-class and splitting axioms of the virtual class, used in standard form, $\Omega$ is the $e_{\TCY}(R\pi_*f^*\Nb)^{-1}$-twisted theory of $\bbP^1$ in the sense of \cite[Section~1]{CG07}, a cohomological field theory with flat unit $\ideone$ over $\KCY\fps{\cQ}$ for the pairing $\eta(a,b)=\int_Zab$, whose inverse is the $\eta^{-1}$ of \cref{prop:set-frob}.

\begin{lem}[{\Lref{Lemma~5.3}: the genus-zero potential}]\genusstatus{AI-proposed, unchecked}
\label{lem:ide-g0}
In the flat coordinates $t_0,t_1$ along $\ideone,H$, with $\cQ=\re^{t_1}$, the genus-zero potential is $\cF_0=\frac16\int_Z(t_0+t_1H)^3+\Li_3(\re^{t_1})/(\epsilon_1\epsilon_2)$; that is, $\GW_{0,d}=1/(d^3\epsilon_1\epsilon_2)$ for $d\geq1$ and $\langle H,H,H\rangle_{0,d}=d^3\GW_{0,d}=1/(\epsilon_1\epsilon_2)$. Consequently $H\qprod H=(\Lv_++\Lv_-)H-\Lv_+\Lv_-\ideone$ with $\Lv_++\Lv_-=E/\tau-\alpha_0$ and $\Lv_+\Lv_-=\alpha_1\alpha_2/\tau$, where $\Lv_\pm$ are the two eigenvalues of $H\qprod$.
\end{lem}

\begin{lem}[{\Lref{Lemma~5.4}: semisimplicity at $\cQ=0$}]\genusstatus{AI-proposed, unchecked}
\label{lem:ide-ss}
The eigenvalues $\Lv_\pm$ of $H\qprod$ lie in $\KCY\fps{\cQ}$, with $\Lv_+(0)=0$, $\Lv_-(0)=-\alpha_0$, and $\Lv_+-\Lv_-$ is a unit of $\KCY\fps{\cQ}$. The idempotents $\eid_\pm|_{\cQ=0}$ are the point classes of $p_\pm$. The $\Delta_i=1/\eta(\eid_i,\eid_i)$ are given by \eqref{eq:set-Delta}, and $\Delta_\pm(0)=\eul(\wts_\pm)$, both non-zero in $\KCY$. So $\Omega$ is semisimple over $\KCY\fps{\cQ}$, and the ratio $\Delta_+(0)/\Delta_-(0)=-\alpha_1\alpha_2/(\wA\wB)$ is a non-constant element of $\KCY$.
\end{lem}

\emph{Idea of the proofs.} The value of $\GW_{0,d}$ is the Aspinwall--Morrison formula \cite[p.~3]{FP00}, after localisation on the factor $\bbC^2$, and $\langle H,H,H\rangle_{0,d}=d^3\GW_{0,d}$ is the divisor equation of the twisted theory \cite[App.~1~(iv)]{CG07}; \cref{lem:ide-ss} is Hensel's lemma, as $H\qprod$ has characteristic polynomial $\Lv(\Lv+\alpha_0)$ modulo $\cQ$ and $\alpha_0\neq0$.

% Hodge classes as an R-matrix action; Teleman's classification; flatness
Let $\TFT$ be the semisimple topological field theory over $\KCY\fps{\cQ}$ with $\TFT_{g,n}(\eid_i^{\otimes n})=\Delta_i^{g-1}$ and vanishing mixed insertions. An $R$-matrix is a series $\Rm(z)=1+\Rm_1z+\cdots$ in $\End(\ideV)\fps z$ which is symplectic, $\eta(\Rm(z)a,\Rm(-z)b)=\eta(a,b)$; $\Rm.\TFT$ is the unit-preserving action of \cite[Definitions~2.2, 2.13]{PPZ15}. A scalar series $\exp(h)$, $h=\sum_{k\geq1}b_kz^{2k-1}$, acts in the same way on the trivial rank-one theory $\ideone$, for instance $h_{\mathsf{v}}$ of \eqref{eq:set-h}.

\begin{lem}[{\Lref{Lemmas~5.5--5.7}: Hodge classes as an $R$-action; freeness}]\genusstatus{AI-proposed, unchecked}
\label{lem:ide-mumford}
\begin{enumerate}
\item For all stable $(g,n)$, $\exp(h).\ideone=\exp\big(\sum_{k\geq1}b_k\frac{(2k)!}{B_{2k}}\mathrm{ch}_{2k-1}(\Hdg)\big)$, and $\Hg_{\mathsf{v}}=h_{\mathsf{v}}.\ideone$ for every multiset $\mathsf{v}$ of non-zero weights.
\item $\mathrm{ch}_{2k-1}(\Hdg)\neq0$ on $\Mbar_k$ for $k\geq2$, and on $\Mbar_{k,1}$ for $k\geq1$.
\item Let $\TFT$ be a semisimple topological field theory over a field of characteristic $0$ with $\TFT(\eid_i^{\otimes n})=\Delta_i^{g-1}$ (of any rank). If $\Rm$ is symplectic and $\Rm.\TFT=\TFT$, then $\Rm=1$.
\end{enumerate}
\end{lem}

\begin{proof}[Sketch of proof]
(i) Mumford's formula \cite[Section~1.1]{FP00}, written with markings, expresses $\mathrm{ch}_{2l-1}(\Hdg)$ by the translation, leg and edge terms of the infinitesimal action \cite[Prop.~2.14]{PPZ15}; Mumford's relation $c(\Hdg)c(\Hdg^\vee)=1$ \cite[(5.4), (5.5)]{Mum83} then identifies $\log\Hg_{\mathsf{v}}$ with the class attached to $\log h_{\mathsf{v}}$. (ii) On $\Mbar_k$, $\mathrm{ch}_{2k-1}(\Hdg)$ is a non-zero multiple of $\lambda_k\lambda_{k-1}$, and $\int_{\Mbar_k}\lambda_k\lambda_{k-1}\lambda_{k-2}\neq0$ by \cite[Theorem~4]{FP00}. (iii) The lowest coefficient $\Rm_m\neq0$ of $\Rm-1$ is diagonal in the idempotents, $m=2k-1$ is odd, and on $\Mbar_{k,1}$ it contributes a non-zero multiple of $(\Rm_m)_{ii}\,\mathrm{ch}_{2k-1}(\Hdg)$ to $(\Rm.\TFT)(\eid_i)$.
\end{proof}

\begin{prop}[{\Lref{Proposition~5.8}: existence and field of definition of $\Rm$}]\genusstatus{AI-proposed, unchecked}
\label{prop:ide-teleman}
There is a unique symplectic $\Rm\in\End(\ideV)\otimes\KCY\lser{\cQ}\fps{z}$ with $\Omega=\Rm.\TFT$.
\end{prop}

\begin{proof}[Sketch of proof]
The field $\KCY\lser{\cQ}$ embeds in $\bbC$, and there Teleman's classification \cite[(1.15), Theorem~(2), Prop.~8.7(ii)]{Tel12} gives a symplectic $\Rm$ with $\Omega=\Rm.\TFT$; the definitions and the group property \cite[Defs.~2.2, 2.5, 2.13 and Prop.~2.14]{PPZ15} hold over any field of characteristic $0$. Uniqueness is \cref{lem:ide-mumford}(iii) with the group property; as automorphisms of $\bbC$ over $\KCY\lser{\cQ}$ fix $\Omega$ and $\TFT$, it shows that $\Rm$ descends to $\KCY\lser{\cQ}$.
\end{proof}

Let $\fid_i=\Delta_i^{1/2}\eid_i$. The square roots lie in $\ideKt\fps{\cQ}$, where $\ideKt$ is $\KCY$ with the square roots of $\eul(\wts_\pm)$ adjoined; the field $\ideFz\coloneqq\ideKt\lser{\cQ}$ carries $\thQ=\cQ\,\dd/\dd\cQ=\dd/\dd t$, with constants $\ideKt$. Let $\Psi\colon\ideFz^2\to\ideV\otimes\ideFz$, $\mathbf{e}_i\mapsto\fid_i$, an isometry with $H\qprod\Psi=\Psi\ideLam$, $\ideLam\coloneqq\mathrm{diag}(\Lv_+,\Lv_-)$. The \emph{flatness equation} is
\begin{equation}
\label{eq:ide-flat}
z\,\thQ\big(\Rm(z)\Psi\big)=\big(H\qprod\Rm(z)-\Rm(z)\,H\qprod\big)\Psi .
\end{equation}

\begin{lem}[{\Lref{Lemma~5.9}: ambiguity of flat symplectic solutions}]\genusstatus{AI-proposed, unchecked}
\label{lem:ide-ambig}
Let $\Rm_1,\Rm_2\in\End(\ideV)\otimes\ideFz\fps{z}$ be symplectic, with $\Rm_1(0)=\Rm_2(0)=1$, both solving \eqref{eq:ide-flat}. Then $\Rm_2=\Rm_1\Psi\exp\big(\mathrm{diag}\sum_{k\geq1}a_k^{(i)}z^{2k-1}\big)\Psi^{-1}$ with constants $a_k^{(i)}\in\ideKt$.
\end{lem}

\begin{prop}[{\Lref{Proposition~5.10}: flatness of $\Rm$}]\genusstatus{AI-proposed, unchecked}
\label{prop:ide-flat}
The $R$-matrix $\Rm$ of \cref{prop:ide-teleman} satisfies \eqref{eq:ide-flat}. In terms of the covectors $\vv_i\coloneqq\eid_i^*\circ\Rm^{-1}(-z)$, $\eid_i^*=\Delta_i\eta(\eid_i,\cdot)$, the equation reads $z\Dop_i\vv_i=(\Aop-\Lv_i)\vv_i$, with $\Dop_i$ as in \eqref{eq:set-Hrow} and $\Aop$ the transpose of the matrix of $H\qprod$, acting on covectors by $(\Aop a)(x)=a(H\qprod x)$.
\end{prop}

\begin{proof}[Sketch of proof of \cref{lem:ide-ambig,prop:ide-flat}]
\emph{Ambiguity.} For $\ideZ\coloneqq\Psi^{-1}\Rm_1^{-1}\Rm_2\Psi$ the two equations give $z\thQ\ideZ=[\ideLam,\ideZ]$. Order by order in $z$, since $\Lv_+-\Lv_-$ is a unit, $\ideZ$ is diagonal with $\thQ$-constant entries; symplecticity makes their logarithms odd.

\emph{Flatness.} Put $\Rinv=\Rm^{-1}$. On the moduli space $M_{g,n}$ of smooth curves only the edgeless graph survives, with the legs $\fid_i^*(\Rinv(\psi_l)x_l)$. In the push-forward of $\Omega_{g,n+1}(x,H)$ along the universal curve over $M_{g,n}$, the edgeless graph and the bubbles carrying the legs $l$ and $n+1$ replace $\Rinv(\psi_l)$ on leg $l$ by $\psi_l^{-1}[H\qprod,\Rinv(\psi_l)]$, up to terms with unchanged legs; by the divisor equation of the twisted theory \cite[App.~1~(iv)]{CG07} this push-forward is $\thQ\Omega_{g,n}(x)$. Hence the \emph{defect} $\Mdef(\psi)$, defined by $\thQ(\Psi^{-1}\Rinv)-\Psi^{-1}\psi^{-1}[H\qprod,\Rinv]=\Mdef\,\Psi^{-1}\Rinv$ and independent of $g$, satisfies at $n=2$ a relation in $H^*(M_{g,2})\otimes\ideFz$. By the stability range \cite[2.17--2.19, (2.15)]{Tel12}, as reported there from \cite{Loo96,MW07,Har85,Iva90}, the monomials in $\kappa_j$ ($j\geq1$), $\psi_1,\psi_2$ are linearly independent in $H^*(M_{g,2};\bbQ)$ in cohomological degree $<(g-1)/2$; letting $g$ grow gives $\Mdef(\psi)=\Mdef(0)$ and $\Delta_+^g\Mdef(0)_{+-}+\Delta_-^g\Mdef(0)_{-+}=0$ for every $g\geq2$. Now $\Mdef(0)$ is skew, and $\Delta_+(0)/\Delta_-(0)$ is a non-constant element of $\KCY$ (\cref{lem:ide-ss}), hence not a root of unity; so $\Mdef=0$. Taking $\eta$-adjoints gives \eqref{eq:ide-flat}, whose $i$-th column, lowered by $\eta$, is the equation for $\vv_i$.
\end{proof}

% The rows, the constants, frame B, and the proof of Theorem B
\begin{lem}[{\Lref{Lemma~5.12}: normalisation constants at $\cQ=0$}]\genusstatus{AI-proposed, unchecked}
\label{lem:ide-norm}
The $z$-coefficients of $\rpt_{0,\pm}$ lie in $\KCY\fps{\cQ}$, and $\rpt_{0,+}(z;\cQ=0)=h_{(\alpha_2,-\wA,\alpha_0)}(z)$, $\rpt_{0,-}(z;\cQ=0)=h_{(-\alpha_0,-\alpha_1,\wB)}(z)$. Hence $h_{\wts_+}/\rpt_{0,+}(z;0)=h_{\wts_-}/\rpt_{0,-}(z;0)=\hnw$, $\nw=(\alpha_1,\wA,-\epsilon_1,-\epsilon_2)$, the same scalar at both colours.
\end{lem}

Put $\rZ_{0,i}\coloneqq\hnw\rpt_{0,i}$, $\rZ_{1,i}\coloneqq(\Lv_i+z\Dop_i)\rZ_{0,i}$, $\vv^Z_i\coloneqq(\rZ_{0,i},\rZ_{1,i})$, and let $\RZ$ be the endomorphism with $\eid_i^*\circ\RZ^{-1}(-z)=\vv^Z_i$, $i=\pm$; then $\RZ(0)=1$.

\begin{prop}[{\Lref{Proposition~5.13(ii)--(iv)}: the rows}]\genusstatus{AI-proposed, unchecked}
\label{prop:ide-rows}
\begin{enumerate}
\item With $\Dz=z\,\dd/\dd t$ at fixed $U$, Euler's integration by parts is the identity in $U$ and $\cQ$
\begin{equation}
\label{eq:ide-euler}
\idePF\,\re^{\cW/z}=\thU\big(\idekm\,\re^{\cW/z}\big),\quad \idePF\coloneqq(\cQ-1)\big[\Dz^2-(\Lv_++\Lv_-)\Dz+\Lv_+\Lv_-\big],\quad \idekm\coloneqq-z\alpha_1\cQ\,\frac{1-U}{1-\cQ U},
\end{equation}
and consequently $z\Dop_i\vv^Z_i=(\Aop-\Lv_i)\vv^Z_i$, as an identity in $\KCY(\Ccov)\fps{z}$ and in $\KCY\fps{\cQ}\fps{z}$.
\item The rows lie in $\KCY\fps{\cQ}\fps{z}^2$, and $\vv^Z_i|_{\cQ=0}=h_{\wts_i}(z)\,(1,\Lv_i(0))$: that is, $\RZ|_{\cQ=0}=\mathrm{diag}(h_{\wts_+},h_{\wts_-})$ on the idempotents $\eid_\pm(0)$.
\item $\langle\vv^Z_+(z),\vv^Z_-(-z)\rangle=0$ and $\langle\vv^Z_i(z),\vv^Z_i(-z)\rangle=\Delta_i$ for the pairing \eqref{eq:set-pairing}. Hence $\RZ$ is symplectic and solves \eqref{eq:ide-flat}, and $\Rm=\RZ\,\Psi\exp\big(\mathrm{diag}\sum_{k\geq1}a_k^{(i)}z^{2k-1}\big)\Psi^{-1}$ with constants $a_k^{(i)}\in\ideKt$.
\end{enumerate}
\end{prop}

\begin{proof}[Sketch of proof of \cref{lem:ide-norm,prop:ide-rows}]
\emph{Constants.} At $\cQ=0$ the derivatives $\cW_k(U_\pm)$, $k\geq2$, become those of a phase $a\log U+b\log(1-U)$ at its critical point, with $(a,b)=(-\alpha_2,\wA)$ at $U_+$ and, in the coordinate $\ideQU$, $(a,b)=(\alpha_0,\alpha_1)$ at $U_-$; in both cases $ab(a+b)\neq0$ in $\KCY$. By the Wick form \eqref{eq:set-rpt} the value of the row at $\cQ=0$ is the normalised expansion of this Beta phase, which is $h_{(-a,-b,a+b)}$ by \cref{lem:row-betaphase} (this is where Stirling's series enters).

\emph{Rows.} (i) In the coordinate $\log U$ the row $\rpt_{0,i}$ is the normalised expansion of $\re^{\cW/z}$ with amplitude $1$ at $\log U_i$ (as in \cref{prop:row-wickrow}). By the rule (P2) of \cref{prop:row-calc}, applied with the parameter $t$ and the moving critical point $\log U_i$, the operator $z^{-1}(\Lv_i+z\Dop_i)$ acts on such expansions through $q\mapsto\de_tq+q\,\de_t\cW/z$ on the amplitudes. Hence $(\Lv_i+z\Dop_i)^2-(\Lv_++\Lv_-)(\Lv_i+z\Dop_i)+\Lv_+\Lv_-$ sends $\rpt_{0,i}$ to the expansion with amplitude $(\cQ-1)^{-1}\re^{-\cW/z}\idePF\re^{\cW/z}$, the amplitude of a total derivative by \eqref{eq:ide-euler}; it vanishes by (P3). The computation, made at complex weights with (ND), is transferred to $\KCY(\Ccov)$ by (P) of \cref{lem:set-rings}. (ii) follows from \cref{lem:ide-norm}. (iii) By (i) and the self-adjointness of $\Aop$, the first pairing satisfies a first-order linear equation which forces it to vanish order by order in $z$, and $\langle\vv^Z_i(z),\vv^Z_i(-z)\rangle/\Delta_i$ is independent of $\cQ$, with value $1$ at $\cQ=0$ by (ii). Then apply \cref{prop:ide-flat,lem:ide-ambig}.
\end{proof}

\begin{prop}[{\Lref{Proposition~5.14}: $\Rm=\RZ$}]\genusstatus{AI-proposed, unchecked}
\label{prop:ide-const}
The constants of \cref{prop:ide-rows}(iii) vanish: $\Rm=\RZ$. Consequently, for $g\geq2$, $\sF_g=\int_{\Mbar_g}\Omega_{g,0}$ is the point-frame graph sum for the rows $\vv^Z_i$, that is, the sum \eqref{eq:set-graphsum}, \eqref{eq:set-Gsum} with the vertex class $\Hg_{\bw}$ replaced by $1$ and the factors of \cref{def:set-factors} computed from the rows $\vv^Z_i$.
\end{prop}

\begin{proof}[Sketch of proof]
By \cref{prop:ide-rows}(ii), (iii), $\Rm|_{\cQ=0}=\mathrm{diag}(h_{\wts_p}\idesp_p)$ with $\idesp_p\coloneqq\exp\sum_ka_k^{(p)}z^{2k-1}$, so the $\eid_p$-component of $\Omega|_{\cQ=0}$ is $\Delta_p(0)^{g-1}((h_{\wts_p}\idesp_p).\ideone)_{g,n}$. But only constant maps contribute at $\cQ=0$, and the localisation formula \cite{GP99} on $\Mbar_{g,n}\times\{p_\pm\}$ gives $\Omega_{g,n}(\eid_p,\dots,\eid_p)|_{\cQ=0}=\eul(\wts_p)^{g-1}\Hg_{\wts_p}$, with $\Hg_{\wts_p}=h_{\wts_p}.\ideone$ (\cref{lem:ide-mumford}(i)). By the group property \cite[Prop.~2.14]{PPZ15}, $\idesp_p.\ideone=\ideone$. If $k$ were minimal with $a^{(p)}_k\neq0$, the degree-$(2k-1)$ part of $(\idesp_p.\ideone)_{k,1}$ would be a non-zero multiple of $\mathrm{ch}_{2k-1}(\Hdg)$, against \cref{lem:ide-mumford}(ii).
\end{proof}

\begin{prop}[{\Lref{Proposition~5.15}: the frame-B rows}]\genusstatus{AI-proposed, unchecked}
\label{prop:ide-frameB}
Let $\RB\coloneqq\RZ\hw^{-1}$ with $\bw=(E,-\epsilon_1,-\epsilon_2)$, whose rows are $\eid_i^*\circ\RB^{-1}(-z)=\hw(z)^{-1}\vv^Z_i=(\rr_{0,i},\rr_{1,i})$, the frame-B rows of \cref{def:set-rows}. Then, for $g\geq2$, $\sF_g$ is the Hodge-vertex graph sum \eqref{eq:set-Gsum} with vertex class $\Hg_{\bw}$ and the factors computed from the rows $(\rr_{0,i},\rr_{1,i})$.
\end{prop}

\begin{proof}[Sketch of proof]
By \cref{lem:ide-mumford}(i), $\Hg_{\bw}=\hw.\ideone$, and $\hw$ is a scalar symplectic $R$-matrix, so $\hw.\TFT=\Hg_{\bw}\otimes\TFT$. The group property \cite[Prop.~2.14]{PPZ15} gives $\RB.(\Hg_{\bw}\otimes\TFT)=(\RB\hw).\TFT=\RZ.\TFT=\Omega$, and the graph sum of \cref{prop:ide-const} applies with the vertex values $\Delta_i^{g_v-1}\Hg_{\bw}$; $\hw^{-1}\hnw=h_{(\alpha_1,\wA,-E)}$ by \cref{lem:set-hmult}.
\end{proof}

\begin{lem}[{\Lref{Lemma~5.16}: description of the coloured terms}]\genusstatus{AI-proposed, unchecked}
\label{lem:ide-terms}
Let $\Ccov$ be the double cover $\ydc^2=\Disc(\cQ)$, with function field $\KCY(\Ccov)=\QA$ (\cref{lem:set-disc,def:set-QA}); the two colours are the two roots of $\Pq$, exchanged by the deck involution. Then:
\begin{enumerate}
\item every $z$-coefficient of $\rr_{0,i}$, $\rr_{1,i}$, every $\Dil_i(d)$, every $\Edge_{ij}(a,b)$ and $\Delta_i^{\pm1}$ is an element of $\KCY(\Ccov)$, whose denominators are products of powers of $\cW_2(U_i)$, $1-U_i$, $1-\cQ U_i$, $\ydc$, $\wB\cQ$ and the constants $\alpha_1,\wA,\epsilon_1,\epsilon_2,E$;
\item hence, for $g\geq2$, each coloured term $\GCont(\Gamma,\col)$ is an element of $\KCY(\Ccov)$, regular over every $\cQ\notin\{0,1,\infty\}$ with $\Disc(\cQ)\neq0$;
\item the factors of colour $2$ are the same expressions as those of colour $1$ evaluated at the other root, so the colour sum $\sum_{\col}\GCont(\Gamma,\col)$ is invariant under the deck involution and lies in $\KCY(\cQ)$;
\item the two points of $\Ccov$ over $\cQ=0$ are unramified, and the expansion of every coloured term at either of them lies in $\KCY\fps{\cQ}$; so the total $\Gsum_g$ is regular at $\cQ=0$.
\end{enumerate}
\end{lem}

\begin{proof}[Sketch of proof of \cref{lem:ide-terms} and of Theorem~B]
By \cref{prop:ide-teleman,prop:ide-flat,prop:ide-rows,prop:ide-const,prop:ide-frameB}, $\sF_g=\int_{\Mbar_g}\Omega_{g,0}$ is the frame-B sum $\Gsum_g$, an identity in $\KCY\fps{\cQ}$ whose left side is $\sum_{d\geq0}\GW_{g,d}\cQ^d$ by definition. In \cref{lem:ide-terms} the listed denominators vanish, at finite $\cQ$, only over $\cQ\in\{0,1\}$ or $\Disc=0$, the scalars $\hnw$, $\hw^{-1}$ are the same at both colours, and (iv) is \cref{prop:ide-rows}(ii). So the colour sum lies in $\KCY(\cQ)$, is regular at $\cQ=0$ and has this expansion there.
\end{proof}

% s6_regularity (abridged version). Label prefix: reg. Abridges Section 6 of the long version.
\section{Regularity at the conifold point and the polar part}
\label{sec:reg}

This section sketches the proof of Theorem~C (\cref{thm:int-reg}) from Theorem~A and the two theorems cited from \cite{LP2} (\cref{thm:reg-L,thm:reg-VI}). Everything up to \cref{thm:reg-main} concerns the decoration sum $\GCont(\Gamma,\col)$ of \cref{def:set-graphsum}, as defined, and its Laurent expansion in $t$; its \emph{polar part} is the sum of the terms of negative degree. The scope is $\TCY$, at the generic point ($\RRing\coloneqq\KCY\fps t$) or at any weight point with (ND) ($\RRing\coloneqq\bbC\fps t$); every statement about a graph is per colouring; $g\geq2$ is the genus of the graph sum and $h\geq2$ the genus of a vertex.

\subsection{The Hodge side: dilaton homogeneity and the string identity}
\label{sec:reg-hodge}

We use the objects and the notation of \cref{def:set-fpot}. The string and dilaton equations hold for $\Hg_{\bw}$; this is its only property used in this subsection, and it holds for every triple $\bw$. The series $\rho\qt$ lies in $\Dom$, but its $\psi$-coefficient is $-\rho$, not $-1$: this is why $\Dom$, and not only the dilaton-shifted form, is needed. At the regular idempotent $\qt^{(2)}\notin\Dom$ (\cref{prop:lea-regrows}); there only the dilaton-shifted form is used, and nothing is rescaled.

\begin{lem}[{\Lref{Lemma~6.1}: dilaton homogeneity on $\Dom$}]\genusstatus{AI-proposed, unchecked}
\label{lem:reg-H0}
Let $h\geq2$, $q\in\Dom$, and let $\nu\in1+t\RRing$ be independent of $\psi$. Then $\nu q\in\Dom$, and for $f_1,\dots,f_n\in\RRing\fps\psi$,
\begin{equation}
\label{eq:reg-H0}
\dd^n\Fpot_h(q)[f_1,\dots,f_n]=\nu^{2h-2}\,\dd^n\Fpot_h(\nu q)[\nu f_1,\dots,\nu f_n].
\end{equation}
\end{lem}

\begin{proof}[Sketch of proof]
Since $\qo=\Tdil-\psi$, an insertion of $\qo$ at the base point splits into an insertion of $\Tdil$, which shifts the number $m$ of dilaton leaves, and one of $-\psi$, evaluated by the dilaton equation; the dependence on $m$ cancels, and $\dbase^{s+1}(\mathsf{b}_1,\dots,\mathsf{b}_s,\qo)=(2-2h-s)\,\dbase^s(\mathsf{b}_1,\dots,\mathsf{b}_s)$. Summing the Taylor series at $q=\qo+\upsilon$ gives the Euler identity $\dd^{n+1}\Fpot_h(q)[f,q]=(2-2h-n)\,\dd^n\Fpot_h(q)[f]$ on $\Dom$, which is integrated along the segment from $q$ to $\nu q$.
\end{proof}

Only $\nu=\rho$ is used. The string equation removes an insertion of $1$ and acts on a series $f$ at another marking by $f\mapsto(f-f(0))/\psi$; on the leaves, $(\lf_j-\lf_j(0))/\psi=(j/E)\lf_j$ and $\Tdil/\psi=-\psi\lf_1$.

\begin{lem}[{\Lref{Lemmas~6.3 and~6.4}: string identity; vanishing on $1^p$}]\genusstatus{AI-proposed, unchecked}
\label{lem:reg-string}
For $h\geq2$ and every multiset $\msU$ of integers, $\Xiz_h(1,\msU)=(J_{\msU}/E^2)\,\Xiz_h(\msU)$; iterating, $\Xiz_h(1^p,\msU)=\prod_{i=0}^{p-1}(J_{\msU}+i)\cdot E^{-2p}\,\Xiz_h(\msU)$. In particular $\Xiz_h(1^p)=0$ for $p\geq1$, that is, $\ZZ[h,n]=0$ for $n\geq1$.
\end{lem}

\begin{proof}[Sketch of proof]
By (LAM2) of \cref{lem:set-leafid}, $\lf_1=E^{-1}(1+\psi\lf_1)$. The string equation applied to the insertion $1$ returns $j_i/E$ times the integral without the extra marking from each probe $\lf_{j_i}$, and from the dilaton leaves the integral with one insertion of $-\psi\lf_1$, which cancels the second term. For $\msU=\emptyset$ the factor $J_\emptyset=0$ kills every $p\geq1$.
\end{proof}

\emph{No constant leaf} \Lref{Remarks~6.5 and~6.6}. A constant component (a multiple of $\lf_0=1/E$) in a factor attached to a singular vertex could not be stripped by \cref{lem:reg-string}, and would be a genuine obstruction to the counting below. None occurs: by (L1)--(L3) the coefficients lie in the spaces $\Lsp_m$ and in the span of the products $\lf_i(\psi)\lf_l(\psi')$, $i,l\geq1$, which contain no constants; the constant $E$ in $\qo$ is the base point of the Taylor expansion on $\Dom$, not an insertion.

\subsection{The two cited theorems}
\label{sec:reg-cited}

The companion paper \cite{LP2} works at scale $1$: for an indeterminate $\tau$ (in this subsection only; not $1-1/\cQ$) it puts $\st\coloneqq\tau(1+\tau)$, $\Hg_{(1,\tau,-1-\tau)}=\Lamv(1)\Lamv(\tau)\Lamv(-1-\tau)/(-\st)^h$, $\lf^\tau_j(\psi)=1/(1-j\psi)$, $\Tdil^\tau(\psi)=-\psi^2/(1-\psi)$, and, for a finite multiset $\msU$ of integers with $r=\#\msU$ and $2h-2+r>0$,
\[
\Xit_h(\msU)\coloneqq\sum_{m\geq0}\frac1{m!}\int_{\Mbar_{h,r+m}}\Hg_{(1,\tau,-1-\tau)}\prod_{j\in\msU}\lf^\tau_j(\psi)\prod_{k=1}^{m}\Tdil^\tau(\psi_{r+k}).
\]

\begin{lem}[{\Lref{Lemma~6.7}: scale law}]\genusstatus{AI-proposed, unchecked}
\label{lem:reg-scale}
Let $E,\mathsf{v}_2,\mathsf{v}_3$ lie in a field of characteristic $0$, with $E+\mathsf{v}_2+\mathsf{v}_3=0$ and $E\mathsf{v}_2\mathsf{v}_3\neq0$, and let $\Xiz^{\mathsf{v}}_h(\msU)$ be defined as $\Xiz_h(\msU)$ in \cref{def:set-fpot} with the class $\Hg_{\mathsf{v}}$, $\mathsf{v}=(E,\mathsf{v}_2,\mathsf{v}_3)$, the leaves $\lf_j=1/(E-j\psi)$ and $\Tdil=-\psi^2/(E-\psi)$. Then, with $\tau\coloneqq\mathsf{v}_2/E$, one has $\mathsf{v}=E(1,\tau,-1-\tau)$, $\st=-\mathsf{v}_2\mathsf{v}_3/E^2\neq0$, and for every finite multiset $\msU$ with $r=\#\msU$ and $2h-2+r>0$
\begin{equation}
\label{eq:reg-scale}
\Xiz^{\mathsf{v}}_h(\msU)=E^{-(3h-3+2r)}\,\Xit_h(\msU)\big|_{\tau=\mathsf{v}_2/E}.
\end{equation}
For $\mathsf{v}=\bw=(E,-\epsilon_1,-\epsilon_2)$ on $\TCY$ one has $\tauw=-\epsilon_1/E$, $\st=-\epsilon_1\epsilon_2/E^2$, and $\Xiz^{\bw}_h=\Xiz_h$.
\end{lem}

This is a weight count under $\psi\mapsto\psi/E$ \cite[Corollary~8.4]{LP2}.

\begin{thm}[{\Lref{Theorem~6.8}(i), (ii), (iv): vanishing of the Hodge integrals; cited}]\genusstatus{AI-proposed, unchecked}
\label{thm:reg-L}
\begin{enumerate}
\item[(i)] For every $h\geq2$ and every non-empty multiset $\msM$ of integers $\geq2$ with $\dg(\msM)\leq2h-3$, one has $\Xit_h(\msM)=0$ in $\bbQ(\tau)$ \textup{\cite[Theorem~8.1]{LP2}}.
\item[(ii)] \textup{\cite[Corollary~8.2]{LP2}} For every multiset $\msU$, the empty one included,
\[
\Xit_h(\msU;\tau)=\sum_{k=0}^{h}\Big(-\frac1{\st}\Big)^k\,\Xi^{(k)}_h(\msU),\qquad \Xi^{(k)}_h(\msU)\in\bbQ\ \text{ independent of }\tau .
\]
For $\msM$ as in \textup{(i)} every $\Xi^{(k)}_h(\msM)$ vanishes, so $\Xit_h(\msM;\tau_0)=0$ at every $\tau_0$ in a field of characteristic $0$ with $\st(\tau_0)\neq0$, the diagonal $\tau_0=-1/2$ included.
\item[(iv)] \textup{(\cref{lem:reg-scale} and \cite[Corollary~8.4]{LP2})} Consequently $\Mirr(h)$ of \cref{def:set-mirr} holds for every $h\geq2$, for the class $\Hg_{\bw}$, $\bw=(E,-\epsilon_1,-\epsilon_2)$, at the generic point of $\TCY$ and at every weight point with $\epsilon_1\epsilon_2E\neq0$, the diagonal $\epsilon_1=\epsilon_2$ included: for every non-empty $\msM$ with parts $\geq2$ and $\dg(\msM)\leq2h-3$, $\Xiz_h(\msM)=E^{-(3h-3+2r)}\,\Xit_h(\msM)|_{\tau=\tauw}=0$.
\end{enumerate}
\end{thm}

The theorem rests on the Mari\~no--Vafa formula \cite[eq.~(1)--(3)]{LLZ03} and textbook facts; its proof is not reproduced, neither here nor in [L]. Part (iv) is where the Calabi--Yau condition is used (\cref{rmk:lea-cy}), together with the hypothesis $E=\epsilon_1+\epsilon_2$ of \cref{thm:reg-VI}: the hypothesis of (iv) is that $\bw$ is a sum-zero triple whose \emph{first} weight is the leaf scale $E$ of the leaves produced by Theorem~A. No sign, rationality or integrality condition on the weights is needed, only $\st\neq0$ and $E\neq0$.

\begin{thm}[{\Lref{Theorem~6.9}(ii): the value of the vertex; cited}]\genusstatus{AI-proposed, unchecked}
\label{thm:reg-VI}
For $\bw=(E,-\epsilon_1,-\epsilon_2)$ with $E=\epsilon_1+\epsilon_2$ and every $g\geq2$,
\begin{equation}
\label{eq:reg-VI}
(-\epsilon_1\epsilon_2E^3)^{g-1}\,\ZZ[g,0]=\beta_{2g-2}(\epsilon_1,\epsilon_2),
\end{equation}
identically in $(\epsilon_1,\epsilon_2)$, and at every weight point with $\epsilon_1\epsilon_2E\neq0$ \textup{\cite[Theorem~9.1]{LP2}}.
\end{thm}

Part (i) of \Lref{Theorem~6.9} is the scale-$1$ form \cite[Theorem~9.3]{LP2}; its proof there uses the Mari\~no--Vafa formula at the profile $1^a$ only and nothing from \cref{thm:reg-L}.

\begin{lem}[{\Lref{Lemma~6.10}: the polynomials $\beta_{2g-2}$ and $\Pg$}]\genusstatus{AI-proposed, unchecked}
\label{lem:reg-beta}
Let $\Ssh(y)\coloneqq y/(2\sinh(y/2))$. Then $\Pg=\coef{t^{2g}}\,\Ssh(\epsilon_1t)\Ssh(\epsilon_2t)$ is a symmetric homogeneous polynomial of degree $2g$ in $\epsilon_1,\epsilon_2$, a polynomial in $\epsilon_1^2,\epsilon_2^2$, and $\beta_{2g-2}(\epsilon_1,\epsilon_2)=(2g-3)!\,\Pg/(\epsilon_1\epsilon_2)$ for $g\geq2$.
\end{lem}

\subsection{Theorem R: regularity and the polar part of the decoration sums}
\label{sec:reg-thmR}

The hypotheses $\Mirr(h)$ below are discharged by \cref{thm:reg-L}(iv). Recall that (F1) is the exact identity $\Delta_1\rho^2=-\epsilon_1\epsilon_2E^3/t^2$, with $\rho\in1+t\RRing$ free of $z$ (\cref{prop:lea-key}).

\begin{thm}[{\Lref{Theorem~6.12}: Theorem R on the Calabi--Yau torus}]\genusstatus{AI-proposed, unchecked}
\label{thm:reg-main}
Fix $g\geq2$. Work in frame B with the class $\bw=(E,-\epsilon_1,-\epsilon_2)$, on $\TCY$ at the generic point or at any weight point with (ND).
\begin{enumerate}
\item[(a)] If $\Mirr(h)$ holds for every $2\leq h\leq g-1$, then every coloured stable graph $(\Gamma,\col)$ of genus $g$ with at least one edge is regular at $t=0$ after summing over decorations.
\item[(b)] If $\Mirr(g)$ holds, then the single-vertex graph is regular at colour $2$, and at colour $1$ its decoration sum has the polar part
\begin{equation}
\label{eq:reg-pole}
(-\epsilon_1\epsilon_2E^3)^{g-1}\,\Fpot_g(\qo)\,t^{2-2g}=(-\epsilon_1\epsilon_2E^3)^{g-1}\,\ZZ[g,0]\,t^{2-2g},
\end{equation}
and no other polar term: the coefficients of $t^{2-2g+k}$ vanish for $1\leq k\leq2g-3$.
\end{enumerate}
\end{thm}

\begin{proof}[Sketch of proof]
By \cref{def:set-graphsum,def:set-fpot}, the sum over the dilaton decorations at a vertex of genus $h\geq2$ and colour $i$ whose $n$ half-edges carry the series $f_1,\dots,f_n$ is $\Delta_i^{h-1}\dd^n\Fpot_h(\qt^{(i)})[f_1,\dots,f_n]$, multilinear in the slots and applied term by term to $t$-series of leaves.

\emph{(i) Homogeneity.} Call a colour-$1$ vertex of genus $h\geq2$ \emph{rescaled}. By \eqref{eq:reg-H0} at $\nu=\rho$ and (F1),
\begin{equation}
\label{eq:reg-resc}
\Delta_1^{\,h-1}\,\dd^n\Fpot_h(\qt)[f_1,\dots,f_n]
=\Big(\frac{-\epsilon_1\epsilon_2E^3}{t^2}\Big)^{h-1}\dd^n\Fpot_h(\rho\qt)[\rho f_1,\dots,\rho f_n]:
\end{equation}
every half-edge at a rescaled vertex is rescaled by exactly one factor $\rho$, and the prefactor $\Delta_1^{h-1}$ becomes exactly $(-\epsilon_1\epsilon_2E^3)^{h-1}t^{2-2h}$.

\emph{(ii) Leaf expansion.} By (L1) (\cref{thm:lea-L1}), $\rho\qt=\qo+\sum_{k\geq1}t^kq^{(k)}$ with $q^{(k)}\in\Lsp_{k+1}$. A half-edge of a rescaled vertex $v$ whose edge ends at a regular vertex, or at a singular vertex of genus $0$ or $1$ (not rescaled; the edge factor is then $\rho\Edge_{11}=\rho^{-1}\cdot\rho^2\Edge_{11}$), carries a leg $\rho f_{j,b}$, and $\coef{t^k}(\rho f_{j,b})\in\Lsp_{k+1}$ by (L2) (\cref{prop:lea-L2} for $j=2$). An edge between two rescaled vertices, or a self-loop at $v$, carries $\rho^2\Edge_{11}(\psi,\psi')$, whose coefficient of $t^k$ lies in the span of the products $\lf_i(\psi)\lf_l(\psi')$ with $i,l\geq1$ and $(i-1)+(l-1)\leq k$, by (L3) (\cref{thm:lea-L3}); the bound is split between the two ends (both at $v$ for a self-loop). So the order-$t^k$ coefficient of each factor attached to $v$ is a combination of leaves $\lf_j$, $j\geq1$, of total degree $\dg\leq k$, without constant leaf. Expanding the Taylor series of $\dd^n\Fpot_h$ at $\qo$ by multilinearity, the coefficient of $t^{2-2h+k}$ of \eqref{eq:reg-resc} is a finite combination of numbers $\Xiz_h(\msU)$ with all parts $\geq1$, $\dg(\msU)\leq k$ and $\#\msU\geq n$. At the other vertices only regularity is needed, which is (F1) and (F3): $\Delta_1^{-1}=O(t^2)$ at a singular vertex of genus $0$, $\Delta_1^0=1$ in genus $1$, and $\Delta_2^{h-1}$, $\Dil_i(d)$, $\Edge_{ij}(a,b)$ are analytic at $t=0$.

\emph{(iii) Vanishing.} Write $\msU=(1^p,\msM)$, $\msM$ the parts $\geq2$. By \cref{lem:reg-string}, $\Xiz_h(\msU)$ is a multiple of $\Xiz_h(\msM)$, $\dg(\msM)=\dg(\msU)$, and it vanishes if $\msM=\emptyset$ and $p\geq1$. If $\msM\neq\emptyset$ and $k\leq2h-3$, then $\Xiz_h(\msM)=0$ by $\Mirr(h)$. Hence a non-zero term has total leaf degree $\regD_v\geq2h_v-2$ at every rescaled vertex $v$ carrying at least one leaf.

\emph{(iv) Counting; proof of (a).} In a connected stable graph of genus $g$ with an edge, every vertex has a half-edge and $g_v\leq g-1$ \Lref{Lemma~6.2}, so only $\Mirr(h)$, $2\leq h\leq g-1$, is needed. By (ii) the leaf degree of each factor at a rescaled end is bounded by its $t$-order, so the $t$-order of a non-zero term is at least $\sum_v(2-2h_v+\regD_v)+2n_0$, the sum running over the rescaled vertices and $n_0$ being the number of singular vertices of genus $0$; every summand is $\geq0$ by (iii).

\emph{(v) The single vertex; proof of (b).} At colour $2$ the term is $\Delta_2^{g-1}$ times a finite sum of products of the $\Dil_2(d)$, analytic at $t=0$. At colour $1$ it is $(-\epsilon_1\epsilon_2E^3)^{g-1}t^{2-2g}\,\Fpot_g(\rho\qt)$ by (i); $\coef{t^0}\Fpot_g(\rho\qt)=\Fpot_g(\qo)=\ZZ[g,0]$, and for $k\geq1$ the coefficient of $t^k$ is a combination of numbers $\Xiz_g(\msU)$ with $\msU$ non-empty, parts $\geq1$ and $\dg(\msU)\leq k$, which vanish for $k\leq2g-3$ by (iii) at $h=g$.
\end{proof}

\subsection{Theorem C}
\label{sec:reg-thmC}

Theorem~C (\cref{thm:int-reg}) is \Lref{Corollary~6.13}. Parts (a) and (b) are \cref{thm:reg-main} with \cref{thm:reg-L}(iv), which holds at the generic point of $\TCY$ and at every weight point with (ND), the value \eqref{eq:reg-VI} and \cref{lem:reg-beta}. For the first half of (c), $\Gsum_g$ is a finite sum of Laurent series in $t$, of which only the single-vertex graph at colour $1$ has negative powers. For the second half, each $\GCont(\Gamma,\col)$ is an element of $\KCY(\Ccov)$ (\cref{lem:ide-terms}); as $\Disc(1)=E^2\neq0$, the point $\cQ=1$ is not a branch point, the expansions at the two points of $\Ccov$ over $\cQ=1$ in the parameter $t=\log\cQ$ are the formal series considered above, and $\sF_g=\Gsum_g$ by Theorem~B. The value $\gapg$ is the coefficient of $t^{2-2g}$ in $\Pg\Li_{3-2g}(\re^t)/(\epsilon_1\epsilon_2)$ (\cref{lem:rig-li}); this is how Theorem~C enters \cref{sec:rig}.

Thus parts (a), (b) and the first half of (c) use Theorem~A and the two cited theorems of \cite{LP2} (input (iii) of \cref{sec:int-inputs}) only, and no result of \cref{sec:ide}. At a weight point with (ND) the same proof applies \Lref{Remark~6.14} and $\gapg$ is a number, which may vanish at special complex weights; no proof step uses its non-vanishing. The first values of $\beta_{2g-2}$ and the instances of $\Mirr$ used in each genus are in \Lref{Remarks~6.11 and~6.15}; the value of the pole uses \cref{thm:reg-VI} and nothing from \cref{thm:reg-L}.

% s7_rigidity (abridged). Label prefix: rig.
\section{Rigidity}
\label{sec:rig}

This section sketches the proof of Theorem~D (\cref{thm:int-rig}; \Lref{Theorem~7.1}) and its consequences. The weights are those of $\TCY$, $\KCY=\Kred(\alpha_0)$ with $\Kred=\bbQ(\alpha_1,\alpha_2,\epsilon_1)$, $\sQ\coloneqq1/\cQ$, and $\sG_g\coloneqq\sF_g-\sF_g(0)$. The poles of $\sF_g$ are located in \cref{lem:rig-poles,lem:rig-disc,lem:rig-inf}, which use Theorem~B, the genus-zero data of \cref{sec:set} and the polynomiality lemma (\cite{GP99}) only; the comparison with $\Li_{3-2g}$ at $\cQ=1$ is the only place where Theorem~C, and through it Theorem~A, enters. Every statement is algebraic in $\KCY(\Ccov)$; no analytic continuation in $\cQ$ is used \Lref{Remark~7.3}.

\begin{lem}[{\Lref{Lemma~7.2}: the possible poles}]\genusstatus{AI-proposed, unchecked}
\label{lem:rig-poles}
$\sF_g$ lies in $\KCY(\cQ)$, is regular at $\cQ=0$ with expansion $\sum_{d\geq 0}\GW_{g,d}\cQ^d$ there, and its poles lie over $\cQ=1$, over $\cQ=\infty$, or over the two roots of $\Disc(\cQ)$.
\end{lem}

\begin{proof}[Sketch of proof]
By \cref{lem:ide-terms} each coloured term of the graph sum lies in $\KCY(\Ccov)$ and is regular over every $\cQ\notin\{0,1,\infty\}$ with $\Disc(\cQ)\neq0$; the colour sum is invariant under the deck involution, hence lies in $\KCY(\cQ)$; and at the two points of $\Ccov$ over $\cQ=0$, which is unramified (\cref{lem:set-disc}) and semisimple (\cref{lem:ide-ss}), every coloured term expands in $\KCY\fps\cQ$ (\cref{prop:ide-rows}). By Theorem~B (\cref{thm:int-ident}) the colour sum is $\sF_g$.
\end{proof}

\begin{lem}[{\Lref{Lemma~7.4}: polynomiality}]\genusstatus{AI-proposed, unchecked}
\label{lem:rig-poly}
For $d\geq 1$ and all $g\geq 0$, the five-torus invariant $\epsilon_1\epsilon_2\GW^{\Tmax}_{g,d}$ is a polynomial: it lies in the ring $\bbQ[\alpha_0,\alpha_1,\alpha_2,\epsilon_1,\epsilon_2]$, is homogeneous of degree $2g$, and is symmetric under $\epsilon_1\leftrightarrow\epsilon_2$, under $\alpha_1\leftrightarrow\alpha_2$ and under $(\alpha_0,\alpha_1,\alpha_2)\mapsto(-\alpha_0,\alpha_0+\alpha_1,\alpha_0+\alpha_2)$. The $\TCY$-invariant $\GW_{g,d}$ is its restriction $\epsilon_2\mapsto\alpha_0+\alpha_1+\alpha_2-\epsilon_1$. Hence, on $\TCY$, $\GW_{g,d}\in(\epsilon_1\epsilon_2)^{-1}\,\bbQ[\alpha_0,\alpha_1,\alpha_2,\epsilon_1]\subset\Kred\fps{\alpha_0}$ for $d\geq 1$.
\end{lem}

\begin{proof}[Sketch of proof]
For $d\geq1$ the space $M=\Mbar_g(X,d)=\Mbar_g(\bbP^1,d)$ is proper and $H^0(C,f^*\cO(-1))=0$, so $R^1_a\coloneqq R^1\pi_*f^*\cO(-1)_a$ ($a=\alpha_1,\alpha_2$) are vector bundles and $e_{\Tmax}(R\pi_*f^*\Nb)^{-1}$ is $(\epsilon_1\epsilon_2)^{-1}$ times the honest equivariant class $e_{\Tmax}(R^1_{\alpha_1})e_{\Tmax}(R^1_{\alpha_2})e_{\Tmax}(\Hdg^\vee_{-\epsilon_1})e_{\Tmax}(\Hdg^\vee_{-\epsilon_2})$ on $M$. By the localisation formula \cite{GP99} the invariant of \cref{def:set-inv} is the $\Tmax$-equivariant integral of this class over the proper space $M$, an element of $H^*_{\Tmax}(\mathrm{pt})$; the degree is a dimension count and the symmetries are induced by automorphisms of $Z$ \Lref{Lemma~2.2}. The same argument gives contravariance: on a subtorus $\mgp{T}'\subset\Tmax$ on which none of the weights of $R\pi_*f^*\Nb$ on the fixed loci vanishes, $\epsilon_1\epsilon_2\GW^{\mgp{T}'}_{g,d}$ is the restriction of the polynomial $\epsilon_1\epsilon_2\GW^{\Tmax}_{g,d}$ (so on $\TCY$, $\Tthree$ and $\TF$).
\end{proof}

\begin{lem}[{\Lref{Lemmas~7.6 and~7.7}: the small root; no pole on $\Disc=0$}]\genusstatus{AI-proposed, unchecked}
\label{lem:rig-disc}
$\Disc$ has a unique root $\cQ_1$ in $\Kred\fps{\alpha_0}$ with $\cQ_1=O(\alpha_0)$; it has $\alpha_0$-adic valuation $2$, $\cQ_1=-\alpha_0^2/(4\alpha_1\alpha_2)+O(\alpha_0^3)$. Consequently $\sG_g$, and hence $\sF_g$, has no pole at the roots of $\Disc$.
\end{lem}

\begin{proof}[Sketch of proof]
Modulo $\alpha_0$, $\Disc=4\alpha_1\alpha_2\cQ+(\alpha_1-\alpha_2)^2\cQ^2$ has the simple root $\cQ=0$, which lifts uniquely to $\Kred\fps{\alpha_0}$. By \cref{lem:rig-poles} and the irreducibility of $\Disc$ over $\KCY$ (\cref{lem:set-disc}), write $\sG_g=N/((1-\cQ)^n\Disc^m)$ with $N\in\KCY[\cQ]$, $n,m\geq0$, $\Disc\nmid N$, and suppose $m\geq1$. \emph{1.} The identity $N=(1-\cQ)^n\Disc^m\sG_g$ holds in $\KCY\fps\cQ\subset\Kred\lser{\alpha_0}\fps\cQ$. On the right, $(1-\cQ)^n\Disc^m$ has coefficients in $\Kred[\alpha_0]$ and $\sG_g=\sum_{d\geq1}\GW_{g,d}\cQ^d$ has coefficients in $\Kred\fps{\alpha_0}$ by \cref{lem:rig-poly}; hence $N\in\Kred\fps{\alpha_0}[\cQ]$. \emph{2.} Because $\val(\cQ_1)=2>0$, the substitution $\cQ\mapsto\cQ_1$ is a well-defined ring homomorphism $\ev\colon\Kred\fps{\alpha_0}\fps\cQ\to\Kred\fps{\alpha_0}$, which on the polynomial $N$ is the ordinary evaluation. Applied to the identity of step~1 it gives $N(\cQ_1)=(1-\cQ_1)^n\Disc(\cQ_1)^m\ev(\sG_g)=0$. The value of the rational function $\sG_g$ at $\cQ_1$ is not used. \emph{3.} $\Disc/(\alpha_1-\alpha_2)^2$ is irreducible over $\KCY$ with root $\cQ_1$, hence its minimal polynomial, so $\Disc$ divides $N$ in $\KCY[\cQ]$: a contradiction. Therefore $m=0$.
\end{proof}

\begin{lem}[{\Lref{Lemma~7.8}: regularity at infinity}]\genusstatus{AI-proposed, unchecked}
\label{lem:rig-inf}
$\sF_g$ is regular at $\cQ=\infty$ (here $\alpha_1\neq\alpha_2$ in $\KCY$ is used).
\end{lem}

\begin{proof}[Sketch of proof]
The point $\cQ=\infty$ is unramified for $\Ccov$ (the leading coefficient $(\alpha_1-\alpha_2)^2$ of $\Disc$ is non-zero), and the two points over $\sQ=0$ are the two colours. The roots of the critical equation at $\sQ=0$ are simple, so $U_1$ and $V_2\coloneqq\cQ U_2$ are power series in $\sQ$ over $\KCY$ (\cref{lem:set-limits}(iii)). Then the $\cW_k(U_i)$, $k\geq2$ (polynomials in $(1-U_i)^{-1}$, $(1-\cQ U_i)^{-1}$ by \cref{lem:set-crit}), the $\Lv_i$, and $\Delta_i^{\pm1}$ (with $\Delta_i(\infty)\neq0$) are power series in $\sQ$; $\dd/\dd t=-\sQ\,\dd/\dd\sQ$ preserves power series in $\sQ$; the $h$-series and $\eta^{-1}$ are constants. Hence every $\Dil_i(d)$, every $\Edge_{ij}(a,b)$, every coloured term and $\sF_g$ itself are power series in $\sQ$.
\end{proof}

\begin{lem}[{\Lref{Lemma~7.9}: the polylogarithms $\Li_{-n}$}]\genusstatus{AI-proposed, unchecked}
\label{lem:rig-li}
For $n\geq 1$ the series $\Li_{-n}(\cQ)=\sum_{d\geq1}d^n\cQ^d$ equals $\Aeul_n(\cQ)/(1-\cQ)^{n+1}$ with $\Aeul_n\in\bbZ[\cQ]$, $\deg\Aeul_n=n$, $\Aeul_n(0)=0$. Hence $\Li_{-n}$ is a rational function of $\cQ$ whose only pole is at $\cQ=1$, of order $n+1$, vanishing at $\cQ=0$ and at $\cQ=\infty$, and $\Li_{-n}(1/\cQ)=(-1)^{n+1}\Li_{-n}(\cQ)$. Under $\cQ=\re^{t}$ its expansion at $t=0$ is $(-1)^{n+1}\,n!\,t^{-n-1}$ plus a power series in $t$, with no other negative power of $t$. For $n=2g-3$ (odd) the polar part is $(2g-3)!\,t^{2-2g}$.
\end{lem}

\begin{proof}[Sketch of proof]
Induction on $n$ with $\Li_{-m-1}=\thQ\Li_{-m}$; $\thQ$ anticommutes with $\cQ\mapsto1/\cQ$; and $\Li_{-n}(\re^t)=(\dd/\dd t)^n\Li_0(\re^t)$, $\Li_0(\re^t)=-1/t+\cdots$. For $n=0$, $\Li_0(1/\cQ)=-1-\Li_0(\cQ)$ does not vanish at $\infty$; this is why $g\geq2$ is needed below.
\end{proof}

\begin{proof}[Sketch of proof of Theorem~D]
This is \Lref{Section~7.5}. Let $g\geq2$ and consider the rational function $D\coloneqq\sF_g-\sF_g(0)-\Pg\Li_{3-2g}/(\epsilon_1\epsilon_2)$ in $\KCY(\cQ)$. By \cref{lem:rig-poles,lem:rig-disc,lem:rig-inf,lem:rig-li}, $D$ has no pole at a finite $\cQ\neq1$ and none at $\cQ=\infty$. At $\cQ=1$, under $\cQ=\re^t$, the polar part of $\sF_g$ is $\gapg t^{2-2g}$ by Theorem~C (\cref{thm:int-reg}(c)), and that of $\Pg\Li_{3-2g}/(\epsilon_1\epsilon_2)$ is $(2g-3)!\,\Pg\,t^{2-2g}/(\epsilon_1\epsilon_2)$ by \cref{lem:rig-li}, which is again $\gapg t^{2-2g}$ by \cref{lem:reg-beta}. Here it matters that Theorem~C gives the polar part exactly, with no terms $t^{3-2g},\dots,t^{-1}$, and that $\Li_{3-2g}(\re^t)$ has none either. So $D$ has no poles on $\bbP^1$, is a constant, and $D(0)=0$: this is the first identity of Theorem~D, and its $\cQ^d$-coefficients give the second in $\KCY$. Both sides of the second identity are polynomials on $\mathrm{Lie}(\TCY)$, by \cref{lem:rig-poly} and because $\Pg\in\bbQ[\epsilon_1^2,\epsilon_2^2]$; so it extends to every point of $\mathrm{Lie}(\TCY)$ by polynomial continuation, not by taking limits of the graph sum.
\end{proof}

\subsection{Genus zero, genus one and the conjecture}
\label{sec:rig-low}

\begin{prop}[{\Lref{Proposition~7.10}: genus zero}]\genusstatus{AI-proposed, unchecked}
\label{prop:rig-g0}
For every $d\geq1$, $\GW_{0,d}=1/(d^3\epsilon_1\epsilon_2)$ at every weight point with $\epsilon_1\epsilon_2\neq0$, i.e.\ $\epsilon_1\epsilon_2\GW_{0,d}=\Ph{0}\,d^{-3}$ with $\Ph{0}=1$ at every point of $\mathrm{Lie}(\TCY)$. This is the Aspinwall--Morrison formula \cite[p.~3]{FP00}.
\end{prop}

\begin{lem}[{\Lref{Lemma~7.11}: genus one}]\genusstatus{AI-proposed, unchecked}
\label{lem:rig-g1}
For every $d\geq 1$, on $\TCY$, $\epsilon_1\epsilon_2\,\GW_{1,d}=\Ph{1}/d=-(\epsilon_1^2+\epsilon_2^2)/(24\,d)$. For $d\geq 2$ this uses, as an input, the formula of \cite{BS} for the resolved conifold on the subtori on which $\alpha_i=\epsilon_j$; for $d=1$ it is proved by direct localisation, without that input.
\end{lem}

\begin{proof}[Sketch of proof]
In genus one $\Lamv(-\epsilon)=-\epsilon-\lambda_1$ and $\lambda_1^2=0$ (Mumford's relation \cite{Mum83}), so in the integral of \cref{lem:rig-poly} $\epsilon_1\epsilon_2\GW^{\Tmax}_{1,d}=\epsilon_1\epsilon_2N_d+(\epsilon_1+\epsilon_2)L_d$ with $N_d\in\bbQ$ and $L_d$ a linear form in $(\alpha_0,\alpha_1,\alpha_2)$; the symmetries of \cref{lem:rig-poly} force $L_d=c\,(\alpha_0+\alpha_1+\alpha_2)$, $c\in\bbQ$. For $d\geq2$: on $\TF$, where the invariants are the restrictions of the five-torus ones, the formula of \cite{BS} gives $\epsilon_1\epsilon_2\GW_{1,d}=\epsilon_1\epsilon_2/(12d)-(\epsilon_1+\epsilon_2)^2/(24d)$, and the linear independence of $\epsilon_1\epsilon_2$ and $(\epsilon_1+\epsilon_2)^2$ gives $N_d=1/(12d)$, $c=-1/(24d)$. Hence $\epsilon_1\epsilon_2\GW^{\Tmax}_{1,d}=\epsilon_1\epsilon_2/(12d)-(\epsilon_1+\epsilon_2)(\alpha_0+\alpha_1+\alpha_2)/(24d)$ on the five-torus, which on $\TCY$ is $\Ph{1}/d$. For $d=1$ the fixed loci of $\Mbar_1(\bbP^1,1)$ have a contracted genus-one component at $p_+$ or at $p_-$; the formula of \cite{GP99} and $\int_{\Mbar_{1,1}}\psi=\int_{\Mbar_{1,1}}\lambda_1=1/24$ give the same five-torus identity at $d=1$ directly.
\end{proof}

\begin{proof}[Sketch of proof of \cref{cor:int-conj} \textup{(\Lref{Corollary~7.13})}]
The conjecture of \cite{BS}, by content: for a Calabi--Yau threefold $X$, a connected reductive group $\mgp{G}$ acting on $Z=X\times\bbC^2$ such that the fixed loci of its maximal torus on the moduli spaces of stable maps are proper, a rigid class $\beta$ and a Calabi--Yau action, $\epsilon_1\epsilon_2\sum_{g\geq0}\GW_g(Z,\mgp{G};\beta)$ lies in $\bbQ\fps{\epsilon_-^2,\epsilon_+^2}$, $\epsilon_\pm=\frac12(\epsilon_1\pm\epsilon_2)$; the class $\beta=0$ is excluded. The identity $\epsilon_1\epsilon_2\GW_{g,d}=\Pg d^{2g-3}$ is Theorem~D for $g\geq2$, \cref{lem:rig-g1} for $g=1$ and \cref{prop:rig-g0} for $g=0$; summing over $g$ gives the closed formula, and each $\Pg$ is a symmetric polynomial in $\epsilon_1^2,\epsilon_2^2$. For a subtorus of $\TCY$ on which $R\pi_*f^*\Nb$ is moving on the fixed loci the invariants are the restrictions of those of $\TCY$ (contravariance, \cref{lem:rig-poly}), and for a connected $\mgp{G}$ with such a maximal torus $\mgp{T}'$ and Weyl group $W$ the $\mgp{G}$-invariant is the element of $R_{\mgp{G}}=R_{\mgp{T}'}^{W}$ with image the $\mgp{T}'$-invariant.
\end{proof}

\subsection{The flop, the constant term and the three-torus}
\label{sec:rig-flop}

\begin{prop}[{\Lref{Proposition~7.15}, parts (i) and (iii): the flop identity $\sF_g(\infty)=\sF_g(0)$}]\genusstatus{AI-proposed, unchecked}
\label{prop:rig-flop}
Let $g\geq 2$ and let $\sF_g(\infty)$ denote the value of $\sF_g$ at $\sQ=0$ (which exists by \cref{lem:rig-inf}).
\begin{enumerate}
\item $\sF_g(\infty)=\sF_g(0)$ in $\KCY$.
\item The constant-map terms of $Z$ and of its flop $\Zfl=\Xfl\times\bbC^2$, $\Xfl$ the other small resolution of the conifold, agree: $\sum_{p=\pm}\int_{\Mbar_g}\prod_{a=1}^{5}\Lamv(w^{(p)}_a)/w^{(p)}_a$ is the same for $(\wts_+,\wts_-)$ and for the tangent weights $(\wts'_+,\wts'_-)$ of $\Zfl$ at its two fixed points,
\[
\wts'_-\coloneqq(\alpha_2-\alpha_1,\alpha_1,\wB,-\epsilon_1,-\epsilon_2),\qquad
\wts'_+\coloneqq(\alpha_1-\alpha_2,\alpha_2,\wA,-\epsilon_1,-\epsilon_2),
\]
an identity among constant-map Hodge integrals for which no direct proof is given here.
\end{enumerate}
\end{prop}

\begin{proof}[Sketch of proof]
(i) is Theorem~D with $\Li_{3-2g}(\infty)=0$ (\cref{lem:rig-li}). For (ii) one computes the data at $\sQ=0$ (this is part (ii) of the long statement, omitted above). By \cref{lem:rig-inf} every factor is a power series in $\sQ$ and $\dd/\dd t=-\sQ\,\dd/\dd\sQ$ kills its value at $\sQ=0$, so $\rr_{1,i}=\Lv_i\rr_{0,i}$ there. The values $\cW_k(U_i)(\sQ=0)$, $k\geq2$, are the $\thU^k$-derivatives of a limiting phase $a\log U+b\log(1-U)$ at its critical point $a/(a+b)$, with $(a,b)=(\alpha_1-\alpha_2,\wA)$ at colour~$1$ and, in the variable $V=\cQ U$, $(a,b)=(-\alpha_2,\alpha_1)$ at colour~$2$. By the Beta-phase lemma (\cref{lem:row-betaphase}) the normalised expansion of such a phase is $h_{(-a,-b,a+b)}(z)$; multiplying by $\hnw$ and cancelling the pairs $(v,-v)$ gives $\rZ_{0,1}(z;\infty)=h_{\wts'_-}(z)$ and $\rZ_{0,2}(z;\infty)=h_{\wts'_+}(z)$, and $\Delta_1(\infty)=\eul(\wts'_-)$, $\Delta_2(\infty)=\eul(\wts'_+)$ by \cref{lem:set-limits}(iii). So at $\sQ=0$ the data of the graph sum of $Z$ are those of the graph sum of $\Zfl$ at its own point $\cQ'=0$, where it takes the constant-map value (\cref{lem:ide-norm,prop:ide-const} read for $\Zfl$); combine with (i). Part (ii) is subject to a reservation, which does not concern (i): in [L] the weights of the two small resolutions are stated without reference, and the identification of the factors of $Z$ at $\sQ=0$ with those of $\Zfl$ at $\cQ'=0$ (the eigenvalues and the pairing are those of $Z$ on one side and of $\Zfl$ on the other, and enter only through the rows, the $\Delta_i$ and the orthogonality of the eigen-covectors) is not spelled out in full and is marked there as a step to be checked.
\end{proof}

\emph{The three-torus.} On $\Tthree=\{\alpha_0=0\}\cap\TCY$, with function field $\Kthree=\bbQ(\alpha_1,\alpha_2,\epsilon_1)$, the point $\cQ=0$ is not semisimple; Theorem~B is obtained there by specialisation.

\begin{lem}[{\Lref{Lemma~7.17}: Theorem B on $\Tthree$}]\genusstatus{AI-proposed, unchecked}
\label{lem:rig-3Tid}
Let $g\geq 2$. On $\Tthree$, with $\alpha_1\alpha_2\epsilon_1\epsilon_2E\neq 0$, the frame-B graph sum \eqref{eq:set-graphsum} formed with the rows $\rr_{0,i}=h_{(\alpha_1,\alpha_2,-E)}\rpt_{0,i}$, $\rr_{1,i}=(\Lv_i+z\Dop_i)\rr_{0,i}$ and $\bw=(E,-\epsilon_1,-\epsilon_2)$ lies in $\Kthree(\cQ)$ and equals $\sF^{\Tthree}_g=\sum_{d\geq 0}\GW^{\Tthree}_{g,d}\cQ^d$, the series of the $\Tthree$-equivariant invariants in the localisation definition.
\end{lem}

\begin{cor}[{\Lref{Corollary~7.18}: Theorem C on $\Tthree$}]\genusstatus{AI-proposed, unchecked}
\label{cor:rig-3T}
Let $g\geq 2$. On $\Tthree$, at the generic point and at every specialisation with $\alpha_1\alpha_2\epsilon_1\epsilon_2E\neq 0$, Theorem~C holds for the frame-B graph sum of \cref{lem:rig-3Tid}; in particular the polar part of $\sF^{\Tthree}_g$ at $\cQ=1$ is exactly $\gapg t^{2-2g}=\beta_{2g-2}(\epsilon_1,\epsilon_2)t^{2-2g}$.
\end{cor}

\begin{proof}[Sketch of proof of \cref{lem:rig-3Tid,cor:rig-3T}]
Apply the ring map $\alpha_0\mapsto0$ to the identity of Theorem~B in $\KCY(\cQ)$. For $d\geq1$ the invariants are regular along $\alpha_0=0$ and restrict to the $\Tthree$-invariants by \cref{lem:rig-poly}; for $d=0$ the two simple poles of $\GW_{g,0}$ at $\alpha_0=0$, from the factors $\Lamv(\pm\alpha_0)/(\pm\alpha_0)$ at $p_\pm$, have opposite residues, and the value is the $\Tthree$-invariant. On the graph side, at generic $\cQ$ every factor is regular along $\alpha_0=0$, $\hnw\hw^{-1}$ specialises to $h_{(\alpha_1,\alpha_2,-E)}$, and the frame change of \cref{prop:ide-frameB} applies verbatim. For the corollary: by \cref{cor:lea-a0}, at $\alpha_0=0$ the frame-B rows and factors are the fibrewise ones at $\xc=\alpha_1/\alpha_2$ rescaled by $z\mapsto z/\alpha_2$, and Theorem~A holds there (its proof never uses the point $\cQ=0$). The proof of \cref{thm:reg-main} then applies verbatim, with \cref{thm:reg-L} at $\tauw=-\epsilon_1/E$ and the value of \cref{thm:reg-VI}: the rows carry $(\alpha_1,\alpha_2)$ only, the vertex class carries $(\epsilon_1,\epsilon_2)$ only, and no relation between them is needed, because $\Mirr(h)$ is a vanishing statement for each multiset separately. Rigidity on $\Tthree$ is the restriction of Theorem~D by polynomiality; the $\alpha_0$-adic argument of \cref{lem:rig-disc} is not available there \Lref{Remark~7.19}.
\end{proof}

% s8_evidence (abridged). Label prefix: evi.
\section{Evidence and what remains conditional}
\noindent\genusstatus{numerically supported} The finite-order checks recorded here are computations (exact arithmetic unless stated); they support the statements at the orders and genera given, and are not proofs. The code and outputs are in \texttt{code/} of this project.

\label{sec:evi}

\subsection*{Finite-order checks}

The long version records, in \Lref{Appendix~A}, the finite-order computations that accompanied the work. None of them is a proof step, and no statement of \cref{sec:set,sec:row,sec:lea,sec:ide,sec:reg,sec:rig} refers to them; they are evidence at the orders and weight points tested. The arithmetic is exact, in $\bbQ$ or in a prime field $\bbF_p$ ($p=2^{127}-1$ or $2^{89}-1$, where a zero test is probabilistic); several implementations written separately were used, together with \emph{controls}, deliberately wrong inputs (a wrong class, a missing constant, a point off $\TCY$) that must make a check fail and did. In outline:
\begin{enumerate}
\item \emph{Identities and rows.} The identities of \cref{sec:set,sec:row,sec:lea} at random rational points of $\TCY$; the closed form of the singular row (\cref{thm:row-main}) through $\tau^8z^{11}$; in an independent computation to $t^{14}z^{23}$ at nine points of $\TCY$ (including $\alpha_1=\alpha_2$, $\uq=0$, $\epsilon_1=\epsilon_2$): (F1)--(F3), (L1) for $k\leq13$, (L2) and (L3) for $k\leq9$, \cref{thm:lea-qde,thm:lea-flow}, and the rescaling at $\alpha_0=0$ (\cref{cor:lea-a0}).
\item \emph{Graph regularity.} The graph sum \eqref{eq:set-graphsum} evaluated per coloured stable graph, summed over decorations: all $20$, $350$, $8596$ and $284806$ coloured stable graphs with an edge of genus $2,3,4,5$ are regular at $t=0$ (genus $2$--$4$ at several generic and special points of $\TCY$, genus $5$ at a single point); the single-vertex graph has the single polar term $\gapg t^{2-2g}$ at colour $1$ and no pole at colour $2$. Other classes at colour $1$, or the point frame, used as controls, give singular edge graphs; the class at colour $2$ is not constrained by these tests.
\item \emph{Direct localisation}, independent of the $R$-matrix: $\GW_{g,d}=\Pg d^{2g-3}/(\epsilon_1\epsilon_2)$ for $g\leq7$ at $d=1$, $g\leq5$ at $d\leq3$ and $g\leq4$ at $d\leq4$, at points of $\TCY$; the control off $\TCY$ fails for every $g\geq1$ tested.
\item \emph{Rational reconstruction in $\cQ$} at genus $2,3$: the individual graphs have poles on $\Disc=0$, of orders $4$ and $8$, which cancel in the total (compare \cref{lem:rig-disc}).
\item \emph{The flop identity} of \cref{prop:rig-flop}(ii) for $g=2,\dots,6$ at two points of $\TCY$; it fails at a point off $\TCY$. The constant term $\sF_g(0)$ was observed to depend on $\alpha$.
\item \emph{Off the Calabi--Yau torus}, with the class $(E_\alpha,-\epsilon_1,-\epsilon_2)$, $E_\alpha=\alpha_0+\alpha_1+\alpha_2\neq\epsilon_1+\epsilon_2$: the leaf structure survives in the tests (\cref{rmk:lea-cy}), but at genus $2$ the vertex graph acquires a $t^{-1}$ term, and at genus $3$ and $4$ coloured graphs with an edge are singular ($4$ of $350$, $54$ of $8596$). This is evidence only; nothing is claimed off $\TCY$.
\end{enumerate}
All checks are at finite genus, at finite orders in $t$, $z$, $\psi$, and at finitely many weight points. The $\alpha_0$-adic step of \cref{lem:rig-disc} and the regularity at $\cQ=\infty$ (\cref{lem:rig-inf}) were tested only through the totals and the rational reconstruction. Not tested at all: the stability argument and the descent in the Teleman step of \cref{sec:ide}, the analytic input \cref{lem:row-betaphase} beyond the values it predicts, and the identification behind \cref{prop:rig-flop}(ii) beyond the Hodge-integral identity itself.

\subsection*{What remains conditional}

The results are conditional only on the cited theorems (i)--(iv) of \cref{sec:int-inputs}, whose proofs are not reproduced; which result uses which input is stated there and is unchanged by the abridgement. None of the arguments uses topological recursion, a mirror theorem, a degeneration argument or a finite-order computation; the one input proved by degeneration is (iv), used in genus $1$ for $d\geq2$ only. The two theorems of \cite{LP2} (\cref{thm:reg-L,thm:reg-VI}) are quoted, not proved, and enter Theorem~C only; the numbering of \cite{LP2} quoted here is provisional, the locators in \cite{Tel12} are those of the arXiv version, and the results of \cite{BS} are referred to by content. The fundamental-class and splitting axioms of the virtual class are used in standard form, without a reference (\cref{sec:ide}). What is not claimed is listed in \cref{sec:int-notclaimed}.

\subsection*{Steps of the long version still to be checked}

The long version is a draft, and it marks the following steps as written out in more detail than in its sources, or as still to be checked by the authors. They are not further inputs from the literature, but the statements that rest on them should be read with this reservation.
\begin{enumerate}
\item The Stirling step in the proof of the Beta-phase lemma (\cref{lem:row-betaphase}), on which input (i) rests in the form used.
\item In \cref{sec:ide}: the list of weights showing that $R\pi_*f^*\Nb$ is moving on the fixed loci, behind \cref{def:set-inv}; and the projection-formula argument for $\Hg_{\bw}\otimes\TFT=\hw.\TFT$ in the proof of \cref{prop:ide-frameB}.
\item In \cref{lem:rig-poly}: the restriction of the invariants to a general subtorus (contravariance), which the sources state for $\TCY$ and $\Tthree$ only; it is used in \cref{lem:rig-g1} on $\TF$, in \cref{cor:int-conj} for subtori, and in \cref{lem:rig-3Tid}.
\item In \cref{lem:rig-g1}: the sign conventions of the node and edge factors in the localisation formula at $d=1$; its normalisation agrees with the genus-zero value, and the resulting five-torus value was confirmed by a separate direct-localisation computation recorded in \Lref{Appendix~A}.
\item In \cref{prop:rig-flop}(ii): the weights of the two small resolutions and the identification of the factors at $\sQ=0$ with those of the flop, as said after its sketch. The identity itself was checked for $2\leq g\leq6$ at two points of $\TCY$ (see the finite-order checks above).
\item In \cref{lem:rig-3Tid}: that the ring map $\alpha_0\mapsto0$ may be applied to Theorem~B as stated over $\KCY$, and that for $d\geq1$ the restricted invariants are the $\Tthree$-invariants; \cref{cor:rig-3T} depends on this lemma only for its statement about $\sF^{\Tthree}_g$.
\end{enumerate}
Theorems~A--D on $\TCY$ at genus $g\geq2$ depend, among these, on (i) and (ii) only; (iii) and (iv) concern genus one and the extension to subtori, (v) the flop identity, and (vi) the three-torus.


\bibliographystyle{amsalpha}
\bibliography{refs}

\end{document}
